% =====================================================================
%  Chapitre « Méthodes de déploiement d'OpenStack »
%  Fichier importé par main.tex avec % =====================================================================
%  Chapitre « Méthodes de déploiement d'OpenStack »
%  Fichier importé par main.tex avec % =====================================================================
%  Chapitre « Méthodes de déploiement d'OpenStack »
%  Fichier importé par main.tex avec % =====================================================================
%  Chapitre « Méthodes de déploiement d'OpenStack »
%  Fichier importé par main.tex avec \include{chap_deploiement}.
%  Il ne contient ni préambule ni \begin{document} : les environnements
%  (definition, resultat, remarque, pointcle, aretenir, savoirfaire, exercice,
%  correction), les styles TikZ (pcli, pfoc, psvc, pdat, pgris, fluxo, fluxr,
%  fluxd, etiqf, callout, num, acteur, vie, msg, rep) et les couleurs (insarouge,
%  insagris, insaorange, insapeche) sont définis dans main.tex. Un chapitre = une \section.
%  Sources : support « Méthodes de déploiement » (chapitre 4) et documentation
%  officielle de DevStack, Kolla, Kolla-Ansible, OpenStack-Ansible, OpenStack-Helm,
%  du guide d'installation, des guides de haute disponibilité et de mise à niveau
%  (voir sources.bib). Versions décrites : celles de la documentation en vigueur
%  début octobre 2026 (version 2026.2 « Hibiscus » sortie le 30 septembre 2026).
%  Écarts avec le support : TripleO est retiré du projet (dépôt archivé en 2024) ;
%  Kolla-Ansible ne déploie plus Swift et ne déploie pas Ceph ; les images publiées par
%  Kolla ne sont pas destinées à la production ; la documentation de Kolla-Ansible ne
%  décrit pas de retour arrière après une mise à niveau.
%  Chapitre détaillé : 15 pages au plus.
% =====================================================================

\part{Mettre en œuvre}\label{part:miseenoeuvre}
\partdesc{Passer de la théorie à un cloud en fonctionnement : choix de l'outil de déploiement, topologie, installation, exploitation et mises à niveau.}

% =====================================================================
\section{Méthodes de déploiement d'OpenStack}\label{chap:deploiement}
% =====================================================================

Les chapitres précédents ont présenté chaque service du point de vue de l'utilisateur, puis de son architecture. Celui-ci se place
du point de vue de l'administrateur, qui doit installer la plateforme et la faire vivre. Il explique pourquoi ce travail est délicat,
décrit ce que tout déploiement assemble, présente les outils les plus répandus (DevStack, Kolla-Ansible, OpenStack-Ansible,
installation manuelle, approches fondées sur Kubernetes), aide à choisir, puis termine par l'exploitation d'un cloud installé :
haute disponibilité, sécurité, mises à niveau.

% ---------------------------------------------------------------------
\subsection{Pourquoi déployer OpenStack est délicat}\label{sec:dep-pourquoi}
% ---------------------------------------------------------------------

OpenStack n'est pas un logiciel unique mais une constellation de services (chapitres~\ref{chap:keystone} à~\ref{chap:trove}) :
chacun a son code, ses processus, son fichier de configuration et, le plus souvent, sa propre base de données. Tous reposent sur la
même infrastructure : une base SQL, une file de messages et un cache (\gls{memcached}) \parencite{install-env}. Un cloud réel
s'étend de plus sur plusieurs machines aux rôles différents, dont les versions doivent rester cohérentes entre elles. Tout installer,
configurer et mettre à jour à la main devient vite ingérable au-delà de quelques nœuds : c'est ce que les outils de déploiement
automatisent (figure~\ref{fig:dep-pile}).

\begin{definition}{Outil de déploiement}
Logiciel qui installe, configure et relie les services d'un cloud sur un ensemble de machines à partir d'une description
(inventaire, fichier de variables, modèles), de façon reproductible ; la plupart savent aussi le mettre à jour.
\end{definition}

\begin{figure}[!ht]
\centering
\adjustbox{max width=\linewidth}{%
\begin{tikzpicture}[>=Stealth, font=\footnotesize]
  \node[pcli, minimum width=9.6cm] (u) at (5.0,7.0) {Utilisateurs \ $\cdot$ \ ligne de commande \ $\cdot$ \ Horizon \ $\cdot$ \ API REST};
  \node[pfoc, minimum width=9.6cm, minimum height=1.0cm] (s) at (5.0,5.4)
    {Services OpenStack\\ {\scriptsize\mdseries Keystone, Nova, Neutron, Glance, Cinder, Octavia, Heat\dots}};
  \node[psvc, minimum width=9.6cm, minimum height=1.0cm] (i) at (5.0,3.8)
    {Services d'infrastructure partagés\\ {\scriptsize\mdseries MariaDB \ $\cdot$ \ RabbitMQ \ $\cdot$ \ memcached \ $\cdot$ \ HAProxy}};
  \node[pgris, minimum width=9.6cm, minimum height=1.0cm] (e) at (5.0,2.2)
    {Mode d'exécution\\ {\scriptsize\mdseries paquets \ $\cdot$ \ environnements Python \ $\cdot$ \ conteneurs \ $\cdot$ \ Kubernetes}};
  \node[pgris, minimum width=9.6cm, minimum height=1.0cm] (m) at (5.0,0.6)
    {Machines et réseaux\\ {\scriptsize\mdseries contrôleurs, calcul, stockage ; réseaux de gestion, de tunnel et externe}};
  \node[draw=insarouge, thick, rounded corners=4pt, fill=insarouge!8, minimum width=1.7cm, minimum height=6.0cm] (o) at (11.9,3.0) {};
  \node[rotate=90, font=\footnotesize\bfseries, text=insarouge] at (11.9,3.0) {Outil de déploiement};
  \foreach \b in {s,i,e,m} \draw[fluxr] (o.west |- \b.east) -- (\b.east);
  \node[callout, minimum width=13.2cm] at (6.3,-1.0)
    {L'outil de déploiement automatise toute la pile, pas seulement l'installation des services.};
\end{tikzpicture}}
\caption{Ce qu'un déploiement doit assembler. Les services (chapitres~\ref{chap:keystone} à~\ref{chap:trove}) reposent sur une
infrastructure partagée, installée sur des machines selon un mode d'exécution au choix ; l'outil de déploiement prend en charge
les quatre étages.}
\label{fig:dep-pile}
\end{figure}

Avant d'installer quoi que ce soit, il faut trancher les questions du tableau~\ref{tab:dep-decisions}. La version compte
autant que l'outil : OpenStack publie une version tous les six mois, la plus récente (2026.2 \emph{Hibiscus}) est sortie le
30~septembre 2026 et 2027.1 \emph{Indri} est en développement \parencite{openstack-releases}. Les outils sont eux-mêmes
versionnés par série, et l'on installe la branche de l'outil qui correspond à la version visée \parencite{kolla-ansible-operating}.

\begin{table}[!ht]
\centering
\renewcommand{\arraystretch}{1.3}
\begin{tabular}{@{}>{\raggedright\arraybackslash}p{3.0cm}>{\raggedright\arraybackslash}p{6.3cm}>{\raggedright\arraybackslash}p{6.1cm}@{}}
\toprule
\textbf{Décision} & \textbf{Question à trancher} & \textbf{Exemples de réponses} \\
\midrule
Topologie & Quels rôles, sur combien de nœuds ? & Un nœud pour l'essai ; trois contrôleurs pour la production. \\
Réseaux & Quelles interfaces pour la gestion, les tunnels, l'accès externe ? & Une interface par réseau, ou des VLAN sur un lien. \\
Stockage & Où vivent les volumes, les images, les objets ? & Disques locaux (LVM), cluster Ceph, NFS. \\
Mode d'installation & Paquets, sources, conteneurs ou Kubernetes ? & Voir le tableau~\ref{tab:dep-outils}. \\
Version & Quelle version, à quel rythme la mettre à niveau ? & Toutes les versions (six mois) ou une sur deux (\gls{slurp}). \\
Exploitation & Haute disponibilité, TLS, secrets, sauvegardes, supervision ? & Section~\ref{sec:dep-exploitation}. \\
\bottomrule
\end{tabular}
\caption{Les décisions à prendre avant de déployer.}
\label{tab:dep-decisions}
\end{table}

% ---------------------------------------------------------------------
\subsection{Anatomie d'un déploiement}\label{sec:dep-anatomie}
% ---------------------------------------------------------------------

Quelle que soit la méthode, on retrouve les mêmes rôles de nœuds, les mêmes services d'infrastructure et les mêmes réseaux
(figure~\ref{fig:dep-topo}).

\begin{itemize}
  \item \textbf{Hôte de déploiement} : la machine d'où l'on lance l'outil (\gls{ansible}, le plus souvent) ; en production, une machine
        distincte du cloud \parencite{osa-deploy-host,kolla-ansible-production}.
  \item \textbf{Contrôleurs} : API des services, ordonnanceurs, base de données et file de messages ; un nombre impair, pour
        établir un quorum \parencite{kolla-ansible-production}.
  \item \textbf{Réseau} : agents Neutron, routeurs ; avec Kolla-Ansible, HAProxy et keepalived s'y exécutent aussi
        \parencite{kolla-ansible-production}.
  \item \textbf{Calcul} : hyperviseur et \texttt{nova-compute}, là où vivent les machines virtuelles.
  \item \textbf{Stockage} : \texttt{cinder-volume} et ses disques ; les volumes peuvent aussi résider sur un cluster \gls{ceph} externe.
\end{itemize}

Le guide d'installation décrit une architecture minimale : un contrôleur, qui exécute l'identité, les images, Placement, les parties
de gestion de Nova et de Neutron, des agents réseau et le tableau de bord, et un nœud de calcul, avec en option des nœuds de stockage
bloc et objet. Il précise qu'elle n'est pas destinée à la production : les agents réseau y résident sur le contrôleur et le trafic
des tunnels emprunte le réseau de gestion \parencite{install-overview}.

\begin{figure}[!ht]
\centering
\adjustbox{max width=\linewidth}{%
\begin{tikzpicture}[>=Stealth, font=\footnotesize,
    nd/.style={draw=insaorange, fill=insapeche, thick, rounded corners=3pt, text width=3.3cm, align=center,
               inner sep=4pt, minimum height=2.3cm}]
  \node[pcli, minimum width=5.2cm] (dh) at (8.4,3.5) {Hôte de déploiement\\ {\scriptsize\mdseries Ansible, inventaire, variables}};
  \node[nd] (c) at (1.8,0)  {\textbf{Contrôleurs} {\scriptsize$\times 3$}\\[2pt] {\scriptsize API des services, ordonnanceurs, MariaDB, RabbitMQ, memcached}};
  \node[nd] (r) at (6.0,0)  {\textbf{Réseau}\\[2pt] {\scriptsize agents Neutron, routeurs, HAProxy et keepalived}};
  \node[nd] (k) at (10.2,0) {\textbf{Calcul} {\scriptsize$\times N$}\\[2pt] {\scriptsize hyperviseur, \texttt{nova-compute}, agent du réseau virtuel}};
  \node[nd] (s) at (14.4,0) {\textbf{Stockage}\\[2pt] {\scriptsize \texttt{cinder-volume}, disques locaux ou accès à Ceph}};
  \foreach \n in {c,r,k,s} \draw[fluxd] (dh.south) -- ++(0,-0.4) -| (\n.north);
  \node[etiqf] at (4.6,2.775) {Ansible, par SSH};
  % réseaux
  \foreach \y/\col/\tc/\nom in {-2.5/insarouge/insarouge/Gestion, -3.3/insaorange/insaorange/Tunnel,
      -4.1/insarouge!45/insarouge/Stockage, -4.9/insagris/insagris/Externe}
    { \draw[line width=2pt, \col] (0.2,\y) -- (16.1,\y);
      \node[anchor=east, font=\scriptsize\bfseries, text=\tc] at (-0.1,\y) {\nom}; }
  % branchements : décalage selon le réseau
  \foreach \cx/\y/\off/\col in {
      1.8/-2.5/-1.2/insarouge, 1.8/-4.9/1.2/insagris,
      6.0/-2.5/-1.2/insarouge, 6.0/-3.3/-0.4/insaorange, 6.0/-4.9/1.2/insagris,
      10.2/-2.5/-1.2/insarouge, 10.2/-3.3/-0.4/insaorange, 10.2/-4.1/0.4/insarouge!45,
      14.4/-2.5/-1.2/insarouge, 14.4/-4.1/0.4/insarouge!45}
    { \draw[thick, \col] ({\cx+\off},-1.15) -- ({\cx+\off},\y);
      \fill[\col] ({\cx+\off},\y) circle (2.6pt); }
  \node[callout, minimum width=15.2cm] at (8.2,-5.9)
    {Un même nœud peut tenir plusieurs rôles : l'inventaire de l'outil les décrit, le logiciel ne les fige pas.};
\end{tikzpicture}}
\caption{Topologie d'un déploiement multi-nœuds. L'hôte de déploiement configure les nœuds par SSH (pointillés) ; chaque rôle se
branche (points) sur les réseaux dont il a besoin.}
\label{fig:dep-topo}
\end{figure}

Les services d'infrastructure demandent une attention particulière. La base de données est le plus souvent un cluster MariaDB avec
\gls{galera} (multimaître), la file de messages un cluster \gls{rabbitmq}, et \gls{haproxy} répartit les requêtes des API vers les contrôleurs
\parencite{osa-service-arch}. Ces composants à état imposent un \gls{quorum} : deux instances n'apportent aucune haute
disponibilité (la panne de l'une paralyse l'autre), trois en tolèrent une en panne, cinq en tolèrent deux, et au-delà le gain
disparaît \parencite{kolla-ansible-production}. L'horloge de tous les nœuds doit en outre être synchronisée (\gls{ntp})
\parencite{install-env}.

Les réseaux se séparent par leur rôle (tableau~\ref{tab:dep-reseaux}) ; le trafic de gestion, qui porte les mots de passe et les
messages entre services, doit rester interne \parencite{kolla-ansible-production}.

\begin{table}[!ht]
\centering
\renewcommand{\arraystretch}{1.3}
\begin{tabular}{@{}>{\raggedright\arraybackslash}p{2.4cm}>{\raggedright\arraybackslash}p{7.6cm}>{\raggedright\arraybackslash}p{5.5cm}@{}}
\toprule
\textbf{Réseau} & \textbf{Rôle} & \textbf{Variable Kolla-Ansible} \\
\midrule
Gestion & Dialogue des services entre eux, avec la base et la file ; à garder interne. & \texttt{api\_\allowbreak interface} (défaut : \texttt{network\_\allowbreak interface}) \\
Tunnel & Trafic des réseaux de projets encapsulé (VXLAN, Geneve). & \texttt{tunnel\_\allowbreak interface} \\
Externe & Réseaux externes et adresses flottantes ; l'interface porte \texttt{br-ex} et n'a pas d'adresse IP. & \texttt{neutron\_\allowbreak external\_\allowbreak interface} \\
API publique & Points d'accès exposés aux utilisateurs derrière HAProxy. & \texttt{kolla\_\allowbreak external\_\allowbreak vip\_\allowbreak interface} \\
\bottomrule
\end{tabular}
\caption{Les réseaux d'un déploiement et leurs interfaces dans Kolla-Ansible \parencite{kolla-ansible-production}.}
\label{tab:dep-reseaux}
\end{table}

% ---------------------------------------------------------------------
\subsection{Panorama des méthodes}\label{sec:dep-panorama}
% ---------------------------------------------------------------------

Les outils se distinguent d'abord par la manière dont ils installent les services (figure~\ref{fig:dep-carte}), puis par leur public
(tableau~\ref{tab:dep-outils}). Aucun n'est « le meilleur » : chacun répond à un contexte.

\begin{figure}[!ht]
\centering
\adjustbox{max width=\linewidth}{%
\begin{tikzpicture}[>=Stealth, font=\footnotesize]
  \foreach \x/\t in {1.9/{Paquets de la\\ distribution}, 6.3/{Sources Git}, 10.7/{Images de\\ conteneurs}, 15.1/{Kubernetes}}
    \node[pgris, minimum width=3.9cm, minimum height=1.0cm] at (\x,4.6) {\t};
  \node[pdat, minimum width=3.7cm]  at (1.9,3.1)  {Installation manuelle};
  \node[pdat, minimum width=3.7cm]  at (1.9,2.2)  {Packstack (RDO)};
  \node[pdat, minimum width=3.7cm]  at (6.3,3.1)  {DevStack};
  \node[psvc, minimum width=3.7cm]  at (6.3,2.2)  {OpenStack-Ansible};
  \node[psvc, minimum width=3.7cm]  at (10.7,3.1) {Kolla-Ansible};
  \node[psvc, minimum width=3.7cm]  at (10.7,2.2) {Kayobe};
  \node[psvc, minimum width=3.7cm]  at (15.1,3.1) {OpenStack-Helm};
  \node[psvc, minimum width=3.7cm]  at (15.1,2.2) {Sunbeam (Canonical)};
  \node[psvc, minimum width=3.7cm]  at (15.1,1.3) {RHOSO (Red Hat)};
  \node[pdat, minimum width=0.5cm, inner xsep=4pt] at (3.2,0.0) {};
  \node[anchor=west] at (3.55,0.0) {essai, apprentissage, développement};
  \node[psvc, minimum width=0.5cm, inner xsep=4pt] at (9.4,0.0) {};
  \node[anchor=west] at (9.75,0.0) {production};
  \node[callout, minimum width=16.4cm] at (8.5,-1.2)
    {Ce qui change d'un outil à l'autre n'est pas la liste des services, mais la façon dont ils sont installés et mis à jour.};
\end{tikzpicture}}
\caption{Les outils de déploiement selon le mode d'installation des services. OpenStack-Ansible installe depuis les sources par
défaut (paquets possibles pour certains services) et peut isoler les services dans des conteneurs LXC \parencite{osa-config}.}
\label{fig:dep-carte}
\end{figure}

\begin{table}[!ht]
\centering
\renewcommand{\arraystretch}{1.25}
\small
\begin{tabular}{@{}>{\raggedright\arraybackslash}p{2.5cm}>{\raggedright\arraybackslash}p{5.5cm}>{\raggedright\arraybackslash}p{3.4cm}>{\raggedright\arraybackslash}p{3.8cm}@{}}
\toprule
\textbf{Outil} & \textbf{Principe} & \textbf{Usage visé} & \textbf{Ressources minimales documentées} \\
\midrule
Installation manuelle & Paquets de la distribution, guide d'installation officiel, service par service & Apprentissage, environnement isolé & Contrôleur : 4~Go ; calcul : 2~Go (preuve de concept) \parencite{install-env} \\
\gls{packstack} & Installeur du projet RDO (CentOS Stream) ; haute disponibilité et mises à niveau hors périmètre & Preuve de concept & 16~Go de mémoire \parencite{packstack} \\
DevStack & Scripts shell, installation depuis les sources Git & Développement, tests, démonstration & Machine dédiée (non précisé) \parencite{devstack-index} \\
OpenStack-Ansible & Ansible ; sources dans des environnements virtuels Python, conteneurs LXC & Production & Test (AIO) : 8~vCPU, 8~Go, 50~Go \parencite{osa-aio} \\
Kolla-Ansible & Ansible ; images de conteneurs Docker ou Podman & Production, évaluation & 2~interfaces, 8~Go, 40~Go \parencite{kolla-ansible-quickstart} \\
Kayobe & Kolla-Ansible, plus configuration des hôtes et provisionnement du matériel (Bifrost) & Production, depuis le serveur nu & --- \parencite{kayobe-index} \\
OpenStack-Helm & Charts Helm sur Kubernetes & Cloud sur Kubernetes & Kubernetes 1.33 à 1.35 \parencite{helm-readme} \\
Sunbeam & Snap, Juju, Kubernetes (Canonical OpenStack) \parencite{sunbeam-multinode} & Petits clouds, de un à plusieurs nœuds & Un nœud : 4~cœurs, 16~Gio, 100~Gio \parencite{sunbeam-tuto} \\
RHOSO & Opérateurs sur OpenShift (Red Hat) & Clients de Red Hat & --- \parencite{rhoso-deploy} \\
TripleO & Retiré du projet (dépôt archivé en 2024) & --- & --- \parencite{tripleo-docs} \\
\bottomrule
\end{tabular}
\caption{Les principaux outils de déploiement d'OpenStack.}
\label{tab:dep-outils}
\end{table}

TripleO (« OpenStack sur OpenStack » : un cloud d'administration, l'\emph{undercloud}, déployait le cloud de production,
l'\emph{overcloud}) figure encore dans de nombreux supports ; il n'est plus maintenu. Red Hat, qui en tirait son outil
\emph{director}, déploie aujourd'hui son offre par des opérateurs sur OpenShift (RHOSO) \parencite{tripleo-docs,rhoso-deploy}. La suite
détaille DevStack (section~\ref{sec:dep-devstack}), Kolla-Ansible (section~\ref{sec:dep-kolla}) et OpenStack-Ansible
(section~\ref{sec:dep-osa}), revient sur l'installation manuelle (section~\ref{sec:dep-manuel}) pour comprendre ce que les outils
automatisent, situe les approches fondées sur Kubernetes (section~\ref{sec:dep-autres}), puis aide à choisir
(section~\ref{sec:dep-choix}) et à exploiter (section~\ref{sec:dep-exploitation}).

% ---------------------------------------------------------------------
\subsection{DevStack : l'environnement tout-en-un}\label{sec:dep-devstack}
% ---------------------------------------------------------------------

\begin{definition}{DevStack}
Série de scripts qui déploient rapidement un cloud OpenStack complet sur une seule machine, à partir des dépôts Git de chaque projet
\parencite{devstack-index}.
\end{definition}

\gls{devstack} sert à construire vite des environnements de développement, à décrire des configurations qui fonctionnent, à faciliter les
contributions et à fournir l'environnement des tests d'intégration de chaque modification proposée \parencite{devstack-source}. Il
installe depuis les sources, sur la branche \texttt{master} par défaut ou sur une branche \texttt{stable/<série>} pour une version plus
ancienne : c'est ce qui en fait l'outil idéal pour essayer une modification de code (dans Nova, par exemple) avant de la proposer. Il
demande une machine dédiée et jetable. Ubuntu 24.04 est la plus testée ; le projet suit aussi les deux dernières versions LTS
d'Ubuntu, Rocky Linux 9 et openEuler. L'installation dure de 15 à 30 minutes \parencite{devstack-index}.

\begin{figure}[!ht]
\centering
\adjustbox{max width=\linewidth}{%
\begin{tikzpicture}[>=Stealth, font=\footnotesize]
  \node[pcli, minimum width=3.6cm] (lc) at (1.9,2.4) {\texttt{local.conf}\\ {\scriptsize\mdseries mots de passe, services}};
  \node[pgris, minimum width=3.6cm] (git) at (1.9,0.3) {Dépôts Git des projets\\ {\scriptsize\mdseries opendev.org}};
  \node[pfoc, minimum width=2.6cm, minimum height=1.3cm] (st) at (7.4,1.35) {\texttt{stack.sh}};
  \node[pgris, minimum width=5.6cm] (pk) at (13.0,2.8) {Paquets de la distribution\\ {\scriptsize\mdseries MySQL, RabbitMQ, Apache}};
  \node[psvc, minimum width=5.6cm] (sv) at (13.0,1.35) {Services installés dans \texttt{/opt/stack}\\ {\scriptsize\mdseries fichiers de configuration, bases}};
  \node[psvc, minimum width=5.6cm] (sd) at (13.0,-0.1) {Unités systemd \texttt{devstack@service}};
  \node[pdat, minimum width=9.4cm] (cl) at (9.4,-1.9) {Cloud à un nœud : API, Horizon, fichier \texttt{openrc}};
  \draw[fluxr] (lc.east) -| (st.north) node[num, pos=0.25] {1};
  \draw[fluxr] (git.east) -| (st.south) node[num, pos=0.25] {2};
  \foreach \t in {pk,sv,sd} \draw[fluxo] (st.east) -- (\t.west);
  \node[num] at (9.7,1.35) {3};
  \draw[fluxo] (sd.south) -- (sd.south |- cl.north);
  \node[callout, minimum width=15.6cm] at (8.5,-3.3)
    {DevStack installe depuis les sources, sur une machine dédiée : pour apprendre et développer, pas pour héberger un service réel.};
\end{tikzpicture}}
\caption{Ce que fait \texttt{stack.sh}. Il lit \texttt{local.conf} (1) et les dépôts des projets (2), puis installe les paquets de la
distribution, installe et configure les services dans \texttt{/opt/stack} et les démarre comme unités systemd (3).}
\label{fig:dep-devstack}
\end{figure}

Le script \texttt{stack.sh} ne doit pas être lancé en tant que \texttt{root} : il s'exécute sous un compte dédié, \texttt{stack}, qui
dispose de \texttt{sudo} sans mot de passe \parencite{devstack-index}. Ses services de base (Keystone, Nova, Placement, Glance,
Cinder, Neutron avec OVN et Horizon) sont actifs par défaut ; les autres s'ajoutent avec \texttt{enable\_service} (c'est le cas de Swift)
ou, pour un projet qui n'est pas dans le dépôt de DevStack, avec \texttt{enable\_plugin} \parencite{devstack-source,devstack-config,devstack-plugins}.
Chaque service tourne en unité systemd, dont on lit le journal avec \texttt{journalctl} \parencite{devstack-systemd}.

\par\noindent\begin{minipage}{\linewidth}
\begin{lstlisting}[style=terminal]
$ sudo useradd -s /bin/bash -d /opt/stack -m stack
$ sudo chmod +x /opt/stack
$ echo "stack ALL=(ALL) NOPASSWD: ALL" | sudo tee /etc/sudoers.d/stack
$ sudo -u stack -i
$ git clone https://opendev.org/openstack/devstack
$ cd devstack && nano local.conf
$ ./stack.sh
$ sudo journalctl -f --unit devstack@n-cpu.service
\end{lstlisting}
\end{minipage}\par

\par\noindent\begin{minipage}{\linewidth}
\begin{lstlisting}
[[local|localrc]]
ADMIN_PASSWORD=secret
DATABASE_PASSWORD=$ADMIN_PASSWORD
RABBIT_PASSWORD=$ADMIN_PASSWORD
SERVICE_PASSWORD=$ADMIN_PASSWORD
HOST_IP=192.168.1.10
# ajouter un projet absent des services par défaut (ici Heat)
enable_plugin heat https://opendev.org/openstack/heat
# retirer un service dont on n'a pas besoin
disable_service horizon
\end{lstlisting}
\end{minipage}\par

Le fichier \texttt{local.conf} est le seul à personnaliser : il suit un format de type INI dont la section \texttt{[[local|localrc]]}
fixe les mots de passe (des caractères alphanumériques seulement), l'adresse de l'hôte, les services actifs, et même les dépôts et
les branches à cloner (\texttt{NOVA\_REPO}, \texttt{NOVA\_BRANCH}) \parencite{devstack-config}. Une fois le script terminé, on charge les
identifiants avec \texttt{. openrc} puis on utilise la ligne de commande \parencite{devstack-source}.

Pour plusieurs machines, le contrôleur reçoit un \texttt{local.conf} complet et chaque nœud de calcul un fichier réduit, qui
désigne le contrôleur (\texttt{SERVICE\_HOST}, \texttt{MYSQL\_HOST}, \texttt{RABBIT\_HOST}) et ne garde que des services
comme \texttt{n-cpu}. Il faut ensuite associer chaque nouvel hôte à la cellule, depuis le contrôleur, avec
\texttt{tools/discover\_hosts.sh} (section~\ref{sec:nova-archi}) \parencite{devstack-multinode}.

DevStack n'est pas un installateur généraliste : la plupart des combinaisons d'options sont rarement testées, et il modifie en
profondeur le système hôte \parencite{devstack-overview,devstack-source}. Conçu pour un nœud, il n'offre pas de haute disponibilité.
\texttt{./unstack.sh} arrête les services (en laissant MySQL et RabbitMQ), \texttt{./clean.sh} efface les données.

% ---------------------------------------------------------------------
\subsection{Kolla-Ansible : OpenStack en conteneurs}\label{sec:dep-kolla}
% ---------------------------------------------------------------------

Le projet \gls{kolla} a pour mission de fournir des outils qui créent des conteneurs prêts pour la production et des outils qui déploient
des clouds OpenStack \parencite{kolla-source}. Il faut distinguer ses livrables :

\begin{itemize}
  \item \textbf{Kolla} construit les images de conteneurs de chaque service (\texttt{kolla-build}) ;
  \item \textbf{Kolla-Ansible} fournit les playbooks Ansible qui déploient ces images sur les nœuds ;
  \item \textbf{\gls{kayobe}} ajoute la configuration des hôtes et le provisionnement du matériel (avec Bifrost, qui s'appuie sur
        Ironic) à Kolla-Ansible \parencite{kayobe-index}.
\end{itemize}

\gls{kollaansible} fonctionne avec Docker ou Podman. Il prend en charge comme hôtes Debian 13, Rocky Linux 10 et Ubuntu 24.04, et les
images existent dans les mêmes trois distributions (\texttt{kolla\_base\_distro}) \parencite{kolla-ansible-support}. Par défaut il
tire les images publiées sur Quay.io (espace \texttt{openstack.kolla}). Ses auteurs avertissent que ces images, destinées aux
essais et aux démonstrations, ne subissent aucun examen de sécurité : en production, on construit et on maintient son propre jeu
d'images avec \texttt{kolla-build}, que l'on place dans un \gls{registre} local (recommandé dès que le déploiement compte plusieurs nœuds)
\parencite{kolla-source,kolla-image-building,kolla-ansible-multinode}.

\begin{figure}[!ht]
\centering
\adjustbox{max width=\linewidth}{%
\begin{tikzpicture}[>=Stealth, font=\footnotesize,
    ct/.style={draw=insaorange, fill=insapeche, rounded corners=2pt, font=\scriptsize\ttfamily, minimum height=0.55cm, inner xsep=3pt}]
  \node[pcli, minimum width=7.2cm] (dh) at (3.9,7.1)
    {Hôte de déploiement\\ {\scriptsize\mdseries \texttt{kolla-ansible}, inventaire, \texttt{globals.yml}, \texttt{passwords.yml}}};
  \node[pgris, minimum width=5.6cm] (reg) at (12.2,7.1) {Registre d'images\\ {\scriptsize\mdseries Quay.io ou registre local}};
  \foreach \y/\nom in {5.0/{Contrôleurs $\times 3$}, 3.6/{Réseau}, 2.2/{Calcul}, 0.8/{Stockage}}
    { \draw[draw=insagris, dashed, rounded corners=3pt, fill=insagris!5] (0,{\y-0.58}) rectangle (15.7,{\y+0.58});
      \node[anchor=west, font=\footnotesize\bfseries] at (0.15,\y) {\nom}; }
  \foreach \x/\t in {3.9/keystone, 5.9/nova\_api, 8.5/neutron\_server, 10.9/horizon, 12.7/mariadb, 14.5/rabbitmq}
    \node[ct] at (\x,5.0) {\t};
  \foreach \x/\t in {3.9/haproxy, 6.0/keepalived, 8.9/neutron\_l3\_agent, 12.1/neutron\_dhcp\_agent}
    \node[ct] at (\x,3.6) {\t};
  \foreach \x/\t in {4.2/nova\_compute, 6.8/nova\_libvirt, 10.0/openvswitch\_vswitchd}
    \node[ct] at (\x,2.2) {\t};
  \node[ct] at (4.2,0.8) {cinder\_volume};
  \node[pcli, minimum width=1.0cm, inner xsep=4pt] (vip) at (14.8,3.6) {VIP};
  \node[pgris, dashed, minimum width=3.8cm] (ceph) at (11.8,0.8) {Ceph (externe)};
  \draw[fluxd] (6.0,0.8) -- (ceph.west);
  \draw[fluxd] (dh.south) -- (3.9,5.65) node[etiqf, midway, right=3pt] {Ansible, par SSH};
  \draw[fluxo] (reg.south) -- (12.2,5.65) node[etiqf, midway, right=3pt] {tirer les images};
  \node[callout, minimum width=15.4cm] at (7.85,-0.75)
    {Chaque service tourne dans son conteneur, avec ses dépendances figées dans l'image : pas de conflit de versions entre services.};
\end{tikzpicture}}
\caption{Kolla-Ansible sur plusieurs nœuds. L'hôte de déploiement pilote les nœuds par Ansible (pointillés) ; ceux-ci tirent leurs
images d'un registre. Les groupes de l'inventaire fixent les conteneurs de chaque nœud. HAProxy et keepalived, qui portent
l'adresse virtuelle (VIP), s'exécutent par défaut sur les nœuds réseau. Kolla-Ansible ne déploie pas Ceph : il s'y connecte.}
\label{fig:dep-kolla}
\end{figure}

La configuration tient en trois éléments, placés sur l'hôte de déploiement, qu'installe \texttt{pip} dans un environnement virtuel
\parencite{kolla-ansible-quickstart} :
\begin{itemize}
  \item l'\textbf{\gls{inventaireansible}} Ansible, qui affecte les machines aux groupes \texttt{control}, \texttt{network}, \texttt{compute},
        \texttt{monitoring} et \texttt{storage} (le groupe \texttt{control} compte un nombre impair de nœuds) ; chaque service a lui-même
        son groupe, mais en modifier la répartition peut casser le déploiement \parencite{kolla-ansible-multinode,kolla-ansible-production} ;
  \item \texttt{/etc/kolla/globals.yml}, le fichier principal ;
  \item \texttt{/etc/kolla/passwords.yml}, dont les mots de passe vides se remplissent avec \texttt{kolla-genpwd}.
\end{itemize}

\par\noindent\begin{minipage}{\linewidth}
\begin{lstlisting}[style=terminal]
$ python3 -m venv /opt/kolla-venv && source /opt/kolla-venv/bin/activate
$ pip install git+https://opendev.org/openstack/kolla-ansible@stable/2026.1
$ sudo mkdir -p /etc/kolla && sudo chown $USER:$USER /etc/kolla
$ cp -r /opt/kolla-venv/share/kolla-ansible/etc_examples/kolla/* /etc/kolla
$ cp /opt/kolla-venv/share/kolla-ansible/ansible/inventory/multinode .
$ kolla-ansible install-deps
$ kolla-genpwd
\end{lstlisting}
\end{minipage}\par

\par\noindent\begin{minipage}{\linewidth}
\begin{lstlisting}
# /etc/kolla/globals.yml
kolla_base_distro: "ubuntu"
network_interface: "eth0"
neutron_external_interface: "eth1"
kolla_internal_vip_address: "10.1.0.250"
enable_cinder: true
enable_octavia: true
\end{lstlisting}
\end{minipage}\par

La valeur de \texttt{kolla\_internal\_vip\_address} doit être une adresse inutilisée du réseau de gestion : keepalived la porte et la
déplace en cas de panne ; l'interface externe doit être active, mais sans adresse IP \parencite{kolla-ansible-quickstart}. Kolla-Ansible
est volontairement très directif : quelques variables suffisent à un déploiement fonctionnel, et il ne multiplie pas les options.
Pour tout le reste, les fichiers déposés dans \texttt{/etc/kolla/config/} sont fusionnés avec la configuration par défaut du service
(par exemple, \texttt{nova.conf} avec \texttt{virt\_type = qemu} pour faire tourner le cloud dans des machines virtuelles sans
virtualisation imbriquée) \parencite{kolla-ansible-philosophy}. Les services de base (Keystone, Glance, Nova avec Placement, Neutron,
Heat, Horizon) sont actifs par défaut ; on ajoute les autres avec les variables \texttt{enable\_*} (tableau~\ref{tab:dep-kolla-services})
\parencite{kolla-ansible-source}.

\begin{table}[!ht]
\centering
\renewcommand{\arraystretch}{1.3}
\begin{tabular}{@{}>{\raggedright\arraybackslash}p{2.6cm}>{\raggedright\arraybackslash}p{3.4cm}>{\raggedright\arraybackslash}p{9.5cm}@{}}
\toprule
\textbf{Service} & \textbf{Variable} & \textbf{À savoir} \\
\midrule
Cinder (chap.~\ref{chap:cinder}) & \texttt{enable\_cinder} & Inactif par défaut ; disques LVM (\texttt{enable\_cinder\_\allowbreak backend\_\allowbreak lvm}) ou cluster Ceph externe. \\
Octavia (chap.~\ref{chap:octavia}) & \texttt{enable\_octavia} & Demande des certificats propres (\texttt{kolla-ansible octavia-certificates}) \parencite{kolla-ansible-source}. \\
Ironic (chap.~\ref{chap:ironic}) & \texttt{enable\_ironic} & S'appuie sur la même base de déploiement que le reste. \\
Magnum (chap.~\ref{chap:magnum}) & \texttt{enable\_magnum} & Le pilote de cluster reste à choisir (section~\ref{sec:ma-pilotes}). \\
Trove (chap.~\ref{chap:trove}) & \texttt{enable\_trove} & Les images invitées restent à construire (section~\ref{sec:tr-datastores}). \\
Swift (chap.~\ref{chap:swift}) & --- & Plus déployé par Kolla-Ansible (retiré faute de mainteneurs) ; Glance peut utiliser un Swift externe \parencite{kolla-ansible-source}. \\
\bottomrule
\end{tabular}
\caption{Activer les autres services avec Kolla-Ansible.}
\label{tab:dep-kolla-services}
\end{table}

Le déploiement suit un enchaînement fixe (figure~\ref{fig:dep-kolla-flux} et tableau~\ref{tab:dep-kolla-cmd}). La commande
\texttt{deploy} peut être relancée autant de fois que nécessaire, mais si une tâche d'amorçage échoue, un nouvel appel ne la répare pas ;
le plus rapide est alors de détruire ce déploiement raté (\texttt{kolla-ansible destroy}) \parencite{kolla-ansible-troubleshooting}.

\begin{figure}[!ht]
\centering
\adjustbox{max width=\linewidth}{%
\begin{tikzpicture}[>=Stealth, font=\footnotesize]
  \foreach \x/\c/\d [count=\i] in {1.6/bootstrap-servers/{prépare les hôtes}, 4.8/prechecks/{vérifie les prérequis}, 8.0/pull/{tire les images},
                                   11.2/deploy/{démarre les conteneurs}, 14.4/post-deploy/{écrit \texttt{clouds.yaml}}}
    { \node[pfoc, minimum width=2.9cm, font=\scriptsize\bfseries\ttfamily] (s\i) at (\x,2.0) {\c};
      \node[num] at (\x-1.25,2.55) {\i};
      \node[etiqf] at (\x,1.3) {\d}; }
  \foreach \a/\b in {1/2,2/3,3/4,4/5} \draw[fluxr] (s\a.east) -- (s\b.west);
  \node[anchor=west, font=\footnotesize\bfseries, text=insagris] at (0,-0.1) {Ensuite :};
  \foreach \x/\c in {3.6/reconfigure, 7.0/upgrade, 10.4/mariadb-backup, 13.8/destroy}
    \node[psvc, minimum width=2.9cm, font=\scriptsize\bfseries\ttfamily] at (\x,-0.1) {\c};
  \node[callout, minimum width=15.6cm] at (8.0,-1.4)
    {Les mêmes commandes servent à installer, à reconfigurer et à mettre à niveau : l'inventaire et les variables restent la seule description du cloud.};
\end{tikzpicture}}
\caption{Les cinq étapes d'un déploiement avec Kolla-Ansible, puis les commandes d'exploitation.}
\label{fig:dep-kolla-flux}
\end{figure}

\begin{table}[!ht]
\centering
\renewcommand{\arraystretch}{1.25}
\begin{tabular}{@{}>{\raggedright\arraybackslash\ttfamily\small}p{3.3cm}>{\raggedright\arraybackslash}p{12.2cm}@{}}
\toprule
\textrm{\textbf{Commande}} & \textbf{Effet} \\
\midrule
bootstrap-servers & Prépare chaque hôte : paquets système, moteur de conteneurs, utilisateurs \parencite{kolla-ansible-quickstart}. \\
prechecks & Vérifie que les hôtes sont dans un état où l'on peut déployer \parencite{kolla-ansible-troubleshooting}. \\
pull & Tire les images sur les hôtes, sans toucher aux conteneurs. \\
deploy & Génère les configurations, amorce les services et démarre tous les conteneurs. \\
post-deploy & Écrit les accès administrateur dans \texttt{/etc/kolla/clouds.yaml} \parencite{kolla-ansible-quickstart}. \\
reconfigure & Applique une modification de \texttt{globals.yml} ou de \texttt{config/}, éventuellement pour un seul service (\texttt{-t mariadb}). \\
validate-config & Valide les fichiers de configuration générés (après un premier déploiement). \\
upgrade & Met à niveau d'une série à la suivante : nouvelles images, migrations de schéma \parencite{kolla-ansible-operating}. \\
mariadb-backup & Sauvegarde complète ou incrémentale (\texttt{--incremental}) de MariaDB \parencite{kolla-ansible-mariadb-backup}. \\
destroy & Détruit les conteneurs, les volumes et la configuration des hôtes. \\
\bottomrule
\end{tabular}
\caption{Les commandes de \texttt{kolla-ansible} \parencite{kolla-ansible-operating,kolla-ansible-source}.}
\label{tab:dep-kolla-cmd}
\end{table}

% ---------------------------------------------------------------------
\subsection{OpenStack-Ansible : sources et conteneurs LXC}\label{sec:dep-osa}
% ---------------------------------------------------------------------

\gls{openstackansible} (OSA) déploie OpenStack avec Ansible, lui aussi pour la production. Par défaut, tous les services sont installés
depuis les sources dans des \glspl{venv} Python, puis isolés dans des conteneurs \gls{lxc}. Ce sont des conteneurs « machine »,
proches de petites machines virtuelles, et non des conteneurs d'application comme ceux de Docker ; ils sont facultatifs, et un
déploiement directement sur les hôtes est possible \parencite{osa-about}. Les hôtes pris en charge sont Ubuntu 24.04 et 26.04, Debian 13,
CentOS Stream 10 et Rocky Linux 10 \parencite{osa-about}.

\begin{figure}[!ht]
\centering
\adjustbox{max width=\linewidth}{%
\begin{tikzpicture}[>=Stealth, font=\footnotesize,
    ct/.style={draw=insaorange, fill=insapeche, rounded corners=2pt, font=\scriptsize\ttfamily, minimum height=0.6cm, minimum width=2.1cm},
    lx/.style={draw=insarouge!80!black, fill=insarouge!12, rounded corners=3pt, text width=1.8cm, align=center, font=\scriptsize,
               minimum height=1.35cm, inner sep=2pt}]
  \node[pcli, minimum width=3.4cm] at (3.7,5.4) {Kolla-Ansible};
  \node[pcli, minimum width=3.4cm] at (12.0,5.4) {OpenStack-Ansible};
  \draw[draw=insagris, dashed, rounded corners=3pt, fill=insagris!5] (0,0.6) rectangle (7.4,4.7);
  \node[anchor=north west, font=\footnotesize\bfseries] at (0.15,4.6) {Hôte : Docker ou Podman};
  \node[ct] at (1.4,2.9) {keystone};
  \node[ct] at (3.7,2.9) {nova\_api};
  \node[ct] at (6.0,2.9) {mariadb};
  \node[font=\scriptsize, text width=6.6cm, align=center] at (3.7,1.4)
    {un conteneur d'application par service, créé à partir d'une image (Kolla)};
  \draw[draw=insagris, dashed, rounded corners=3pt, fill=insagris!5] (8.4,0.6) rectangle (15.8,4.7);
  \node[anchor=north west, font=\footnotesize\bfseries] at (8.55,4.6) {Hôte : LXC (ou directement sur l'hôte)};
  \node[lx] at (9.9,2.8) {\textbf{galera}\\ MariaDB};
  \node[lx] at (12.1,2.8) {\textbf{keystone}\\ environnement Python};
  \node[lx] at (14.3,2.8) {\textbf{nova\_api}\\ environnement Python};
  \node[font=\scriptsize, text width=6.8cm, align=center] at (12.1,1.2)
    {conteneurs « machine » LXC ; services installés depuis les sources};
  \node[callout, minimum width=15.6cm] at (7.9,-0.6)
    {Même outil (Ansible), deux façons d'isoler les services : des conteneurs d'application ou des conteneurs « machine ».};
\end{tikzpicture}}
\caption{Kolla-Ansible (à gauche) et OpenStack-Ansible (à droite). Le premier exécute chaque service dans un conteneur
d'application ; le second installe les services dans des environnements virtuels Python, au sein de conteneurs LXC.}
\label{fig:dep-osa}
\end{figure}

Le déploiement distingue l'hôte de déploiement, qui exécute Ansible, et les hôtes cibles. On y décrit le cloud dans
\texttt{/etc/openstack\_deploy/} : \texttt{openstack\_user\_config.yml} affecte les machines aux rôles (au moins trois hôtes
d'infrastructure dans \texttt{shared\_infra\_hosts}), \texttt{user\_variables.yml} règle les options globales et
\texttt{user\_secrets.yml} contient les mots de passe, que \texttt{pw-token-gen.py} génère \parencite{osa-config}. Le choix
\texttt{install\_method} (sources ou paquets) se fait au premier déploiement et ne peut plus changer ensuite \parencite{osa-config}.
Trois \glspl{playbook} font le travail, dans cet ordre \parencite{osa-run} : préparer les hôtes et construire les conteneurs, installer l'infrastructure (memcached, dépôt des sources, Galera, RabbitMQ), puis les services d'OpenStack.

\par\noindent\begin{minipage}{\linewidth}
\begin{lstlisting}[style=terminal]
# openstack-ansible openstack.osa.setup_hosts
# openstack-ansible openstack.osa.setup_infrastructure
# openstack-ansible openstack.osa.setup_openstack
\end{lstlisting}
\end{minipage}\par

Un déploiement « \gls{aio} » (\emph{all-in-one}, AIO), qui réunit tout sur une machine, sert à l'essai ou au laboratoire ; il demande
au minimum 8~vCPU, 8~Go de mémoire et 50~Go de disque (le projet en recommande 16~Go et 80~Go) \parencite{osa-aio}. OSA automatise
aussi les mises à niveau entre versions majeures. Le projet déconseille son choix si l'on veut des conteneurs d'application à 100~\%, si l'on
préfère Puppet, ou si l'on veut installer tous les services depuis les paquets de la distribution, dont la couverture reste partielle
\parencite{osa-about}.

% ---------------------------------------------------------------------
\subsection{Installation manuelle : ce que les outils automatisent}\label{sec:dep-manuel}
% ---------------------------------------------------------------------

Le guide d'installation officiel décrit, commande par commande, un cloud minimal installé avec les paquets de la distribution
\parencite{install-guide,install-overview}. Aucun outil ne s'interpose : on tape les commandes sur chaque nœud. On commence par
l'environnement commun (horloge synchronisée, dépôt des paquets OpenStack, MariaDB, RabbitMQ, memcached)
\parencite{install-env,install-packages-ubuntu,install-sql-ubuntu,install-messaging-ubuntu}, puis on installe les services dans
l'ordre que leurs dépendances imposent : Keystone, Glance, Placement, Nova, Neutron, puis Horizon et, en option, Cinder
\parencite{install-overview}. Chaque service suit le même scénario en six étapes (figure~\ref{fig:dep-manuel}).

\begin{figure}[!ht]
\centering
\adjustbox{max width=\linewidth}{%
\begin{tikzpicture}[>=Stealth, font=\footnotesize,
    cmd/.style={font=\scriptsize\ttfamily, text width=2.55cm, align=center}]
  \node[pgris, minimum width=15.6cm] (env) at (8.0,3.9)
    {Environnement commun, une fois sur le contrôleur : \ NTP \ $\cdot$ \ paquets OpenStack \ $\cdot$ \ MariaDB \ $\cdot$ \ RabbitMQ \ $\cdot$ \ memcached};
  \node[anchor=east, font=\footnotesize\bfseries, text=insarouge] at (15.9,2.85) {Pour chaque service (Glance en exemple) :};
  \foreach \x/\t/\c [count=\i] in {1.35/{Base de\\ données}/{CREATE DATABASE glance;},
        4.05/{Enregistrement\\ dans Keystone}/{openstack endpoint create},
        6.75/{Paquets}/{apt install glance},
        9.45/{Configuration}/{glance-api.conf},
        12.15/{Schéma de\\ la base}/{glance-manage db\_sync},
        14.85/{Redémarrage}/{systemctl restart glance-api}}
    { \node[pfoc, minimum width=2.45cm, minimum height=1.0cm] (s\i) at (\x,1.7) {\t};
      \node[num] at (\x-1.0,2.3) {\i};
      \node[cmd] at (\x,0.55) {\c}; }
  \foreach \a/\b in {1/2,2/3,3/4,4/5,5/6} \draw[fluxr] (s\a.east) -- (s\b.west);
  \draw[fluxd] (s1.north |- env.south) -- (s1.north);
  \node[callout, minimum width=15.6cm] at (8.0,-0.55)
    {Six étapes à répéter pour chaque service, sur chaque nœud : c'est précisément ce que les outils de déploiement automatisent.};
\end{tikzpicture}}
\caption{Le scénario d'installation manuelle d'un service, illustré par Glance \parencite{glance-install}. L'étape~2 inscrit le service
dans le catalogue de Keystone : un utilisateur de service, le rôle \texttt{admin}, une entrée de service et trois points d'accès.}
\label{fig:dep-manuel}
\end{figure}

\par\noindent\begin{minipage}{\linewidth}
\begin{lstlisting}[style=terminal]
$ sudo rabbitmqctl add_user openstack RABBIT_PASS
$ sudo rabbitmqctl set_permissions openstack ".*" ".*" ".*"
$ openstack user create --domain default --password-prompt glance
$ openstack role add --project service --user glance admin
$ openstack service create --name glance --description "OpenStack Image" image
$ openstack endpoint create --region RegionOne \
    image public http://controller:9292
\end{lstlisting}
\end{minipage}\par

Les deux premières commandes préparent RabbitMQ une fois pour toutes \parencite{install-messaging-ubuntu} ; les quatre suivantes
enregistrent Glance dans le catalogue (chapitre~\ref{chap:keystone}), puis on répète la dernière pour les points d'accès interne et
d'administration \parencite{glance-install}. Un nœud de calcul reçoit ensuite \texttt{nova-compute} et l'agent de Neutron, et doit être
associé à sa cellule, depuis le contrôleur, par \texttt{nova-manage cell\_v2 discover\_hosts} \parencite{nova-install-compute}.

Cette méthode a une vraie valeur pédagogique : on voit où chaque mot de passe, chaque base et chaque point d'accès apparaît. Elle ne
tient pas dans la durée : des dizaines de fichiers à garder cohérents, aucune idempotence (relancer une commande qui a déjà réussi
échoue ou duplique), aucune description unique du cloud, et des mises à niveau à la main. Au-delà d'un laboratoire, on délègue ce
travail à un outil.

\begin{remarque}[Deux sens de « bare metal »]
Le support parle d'installation « bare metal ». Le terme recouvre deux idées. Installer les services \emph{directement sur le
système des machines}, sans conteneur, c'est l'installation manuelle (ou OpenStack-Ansible sans LXC). Fournir \emph{des machines
physiques aux utilisateurs} est le rôle d'Ironic (chapitre~\ref{chap:ironic}) ; Kayobe s'en sert en outre pour installer le cloud lui-même
sur des serveurs nus (section~\ref{sec:dep-kolla}).
\end{remarque}

% ---------------------------------------------------------------------
\subsection{OpenStack sur Kubernetes}\label{sec:dep-autres}
% ---------------------------------------------------------------------

Deux approches installent OpenStack \emph{sur} Kubernetes, pour des publics différents. À ne pas confondre avec Magnum
(chapitre~\ref{chap:magnum}), qui fait l'inverse : il fournit des clusters Kubernetes \emph{sur} OpenStack.

\textbf{\gls{openstackhelm}} est une collection de charts Helm qui déploient OpenStack et les services associés sur Kubernetes
\parencite{helm-index,helm-readme}. Sa documentation vise des versions récentes de Kubernetes (1.33 à 1.35) et des hôtes Ubuntu
\parencite{helm-readme}. Elle s'adresse à une équipe qui exploite déjà Kubernetes et sait le faire vivre.

\textbf{\gls{sunbeam}} (Canonical OpenStack) s'installe avec un \emph{snap} et s'appuie sur \gls{juju} et Kubernetes. Le tutoriel officiel
déploie un cloud à un nœud (4~cœurs, 16~Gio de mémoire, 100~Gio de disque, Ubuntu 24.04) en cinq commandes \parencite{sunbeam-tuto} :

\par\noindent\begin{minipage}{\linewidth}
\begin{lstlisting}[style=terminal]
$ sudo snap install openstack
$ sunbeam prepare-node-script --bootstrap | bash -x && newgrp snap_daemon
$ sunbeam cluster bootstrap --accept-defaults --role control,compute,storage
$ sunbeam configure --accept-defaults --openrc demo-openrc
$ sunbeam launch ubuntu --name test
\end{lstlisting}
\end{minipage}\par

Pour agrandir le cloud, \texttt{sunbeam cluster add <FQDN>} produit un jeton que la nouvelle machine donne à
\texttt{sunbeam cluster join --role <rôles> -}. Les rôles sont \texttt{control}, \texttt{compute}, \texttt{network} et
\texttt{storage}, et \texttt{sunbeam cluster resize} répartit les fonctions de contrôle après l'ajout de nœuds de contrôle ;
les nœuds ne se désignent que par leur nom complet (FQDN) \parencite{sunbeam-multinode}. Red Hat suit une logique voisine avec
RHOSO, dont les opérateurs tournent sur OpenShift \parencite{rhoso-deploy}.

% ---------------------------------------------------------------------
\subsection{Choisir une méthode}\label{sec:dep-choix}
% ---------------------------------------------------------------------

Tous ces outils déploient le même OpenStack : les API, les services et les concepts des chapitres précédents ne changent pas. Seules
varient l'installation et l'exploitation, et c'est là que se joue le choix (figure~\ref{fig:dep-choix}).

\begin{figure}[!ht]
\centering
\adjustbox{max width=\linewidth}{%
\begin{tikzpicture}[>=Stealth, font=\footnotesize]
  \node[pcli, text width=6.0cm] (q1) at (3.4,5.3) {Apprendre ou développer, sur une machine jetable ?};
  \node[pcli, text width=6.0cm] (q2) at (3.4,3.1) {La plateforme existante repose-t-elle déjà sur Kubernetes ?};
  \node[pcli, text width=6.0cm] (q3) at (3.4,0.8) {Préférez-vous des conteneurs d'application, un par service ?};
  \draw[fluxd] (q1.south) -- (q2.north) node[etiqf, midway, right=2pt] {non};
  \draw[fluxd] (q2.south) -- (q3.north) node[etiqf, midway, right=2pt] {non};
  \node[pdat, text width=6.2cm] (r1a) at (12.2,5.9) {\textbf{DevStack} : développer, tester};
  \node[pdat, text width=6.2cm] (r1b) at (12.2,4.7) {\textbf{Installation manuelle}, \textbf{Sunbeam} à un nœud : apprendre};
  \node[psvc, text width=6.2cm] (r2)  at (12.2,3.1) {\textbf{OpenStack-Helm}, \textbf{Sunbeam}, \textbf{RHOSO}};
  \node[psvc, text width=6.2cm] (r3a) at (12.2,1.5) {\textbf{Kolla-Ansible} (Kayobe si le matériel est à provisionner)};
  \node[psvc, text width=6.2cm] (r3b) at (12.2,0.1) {\textbf{OpenStack-Ansible} (sources, conteneurs LXC)};
  \draw[fluxr] (q1.east) -- ++(0.75,0) |- (r1a.west);
  \draw[fluxr] (q1.east) -- ++(0.75,0) |- (r1b.west);
  \node[etiqf] at (7.0,5.55) {oui};
  \draw[fluxr] (q2.east) -- (r2.west) node[etiqf, midway, above] {oui};
  \draw[fluxr] (q3.east) -- ++(0.75,0) |- (r3a.west) node[etiqf, pos=0.8, above] {oui};
  \draw[fluxr] (q3.east) -- ++(0.75,0) |- (r3b.west) node[etiqf, pos=0.8, above] {non};
  \node[pdat, minimum width=0.5cm, inner xsep=4pt] at (2.6,-1.1) {};
  \node[anchor=west] at (2.95,-1.1) {essai, apprentissage};
  \node[psvc, minimum width=0.5cm, inner xsep=4pt] at (8.0,-1.1) {};
  \node[anchor=west] at (8.35,-1.1) {production};
  \node[callout, minimum width=15.6cm] at (7.9,-2.2)
    {Le bon choix dépend des compétences de l'équipe (Ansible, Kubernetes) et du support attendu, pas d'un classement général.};
\end{tikzpicture}}
\caption{Arbre de décision pour choisir une méthode de déploiement. Ce n'est pas la seule lecture possible : un fournisseur
(Canonical, Red Hat) peut imposer son offre, un client déjà équipé d'Ansible ira vers Kolla-Ansible ou OpenStack-Ansible.}
\label{fig:dep-choix}
\end{figure}

Quelques critères aident à départager Kolla-Ansible et OpenStack-Ansible, les deux outils de production les plus répandus :

\begin{itemize}
  \item \textbf{Isolation} : Kolla-Ansible isole chaque service dans une image figée ; OpenStack-Ansible installe depuis les sources
        dans des environnements virtuels, au sein de conteneurs LXC \parencite{osa-about}.
  \item \textbf{Configuration} : Kolla-Ansible est directif (quelques variables, puis des fichiers fusionnés) ; OpenStack-Ansible
        expose toute la configurabilité de ses rôles \parencite{kolla-ansible-philosophy,osa-about}.
  \item \textbf{Images} : avec Kolla-Ansible, il faut construire et héberger ses propres images pour la production ; avec
        OpenStack-Ansible, on gère un dépôt de sources et des environnements virtuels.
  \item \textbf{Mises à niveau} : toutes deux les automatisent (\texttt{kolla-ansible upgrade}, playbooks d'OpenStack-Ansible) ;
        aucune ne dispense de lire les notes de version.
\end{itemize}

% ---------------------------------------------------------------------
\subsection{Exploiter un cloud déployé}\label{sec:dep-exploitation}
% ---------------------------------------------------------------------

Déployer n'est que le début. Ce qui distingue un cloud fiable, c'est ce qu'on met en place autour : vérification, haute
disponibilité, sécurité, sauvegardes, supervision, mises à niveau.

\textbf{Vérifier.} Avant d'ouvrir le cloud aux utilisateurs, on contrôle que chaque service est enregistré et vivant. Avec
Kolla-Ansible, \texttt{post-deploy} écrit les accès administrateur dans \texttt{/etc/kolla/clouds.yaml}, et le script \texttt{init-runonce}
crée des réseaux, une image et d'autres ressources d'exemple (démonstration seulement) \parencite{kolla-ansible-quickstart} ; avec OpenStack-Ansible,
on utilise le conteneur utilitaire \parencite{osa-verify}.

\par\noindent\begin{minipage}{\linewidth}
\begin{lstlisting}[style=terminal]
$ openstack endpoint list                # chaque service a ses points d'accès
$ openstack compute service list         # conductor, scheduler, compute : up
$ openstack network agent list           # agents Neutron : Alive
$ openstack volume service list          # cinder-scheduler, cinder-volume : up
$ nova-status upgrade check              # prérequis d'une mise à niveau de Nova
\end{lstlisting}
\end{minipage}\par

\textbf{Haute disponibilité.} Elle repose sur deux principes \parencite{ha-concepts}. Les services sans état (API, ordonnanceurs)
tournent en actif/actif derrière HAProxy et une adresse virtuelle que \gls{keepalived} déplace en cas de panne. Les services à état
(MariaDB avec \gls{galera}, \gls{rabbitmq}) se répliquent en cluster d'un nombre impair de nœuds, dont le quorum est la majorité
(\(\lfloor n/2\rfloor+1\)) ; le \gls{plancontrole} compte donc au moins trois contrôleurs. La haute disponibilité ne protège pas
les machines des utilisateurs : si un hyperviseur tombe, Nova ne les relance pas seul (on évacue les instances avec
\texttt{openstack server evacuate}).

\begin{table}[!ht]
\centering
\renewcommand{\arraystretch}{1.25}
\small
\begin{tabular}{@{}>{\raggedright\arraybackslash}p{2.3cm}>{\raggedright\arraybackslash}p{13.1cm}@{}}
\toprule
\textbf{Sujet} & \textbf{Mise en œuvre avec Kolla-Ansible} \\
\midrule
Répartiteur & HAProxy et keepalived (groupe \texttt{loadbalancer}, qui reprend par défaut les nœuds \texttt{network}) ; variables \texttt{enable\_haproxy} et \texttt{enable\_keepalived} \parencite{kolla-ansible-haproxy}. \\
TLS & \texttt{kolla\_enable\_tls\_external} et \texttt{kolla\_enable\_tls\_internal} ; certificats dans \texttt{/etc/kolla/certificates} (\texttt{haproxy.pem}, \texttt{haproxy-internal.pem}) ; \texttt{kolla-ansible certificates} produit une autorité de test ; \texttt{enable\_letsencrypt} automatise le renouvellement \parencite{kolla-ansible-tls}. \\
Secrets & Mots de passe de \texttt{passwords.yml} (à sauvegarder avant tout changement). Rotation : vider les valeurs, \texttt{kolla-genpwd -p}, puis \texttt{kolla-ansible reconfigure} ; pour RabbitMQ, \texttt{stop} puis \texttt{deploy} \parencite{kolla-ansible-passwords}. \\
Sauvegarde & \texttt{kolla-ansible mariadb-backup} (complète ou incrémentale) \parencite{kolla-ansible-mariadb-backup} ; copie de \texttt{/etc/kolla} et de l'inventaire ; volumes et images selon les chapitres~\ref{chap:cinder} et~\ref{chap:glance}. \\
Journaux & Fichiers dans \texttt{/var/log/kolla/<service>} et \texttt{docker logs} ; \texttt{enable\_central\_logging} ajoute Fluentd et OpenSearch \parencite{kolla-ansible-troubleshooting,kolla-ansible-logging}. \\
Supervision & \texttt{enable\_prometheus} déploie Prometheus, que l'on peut relier à Grafana \parencite{kolla-ansible-prometheus}. \\
\bottomrule
\end{tabular}
\caption{L'exploitation d'un cloud Kolla-Ansible : ce que l'outil apporte.}
\label{tab:dep-exploit}
\end{table}

\textbf{Sécurité et sauvegardes.} Trois règles reviennent partout : garder le réseau de gestion interne (il porte les mots de passe
et les messages), chiffrer les API exposées (TLS) et protéger les secrets (tableau~\ref{tab:dep-exploit}). Une sauvegarde ne vaut que
si on sait restaurer : la base de données de chaque service, la configuration de l'outil et les mots de passe sont à sauvegarder
régulièrement et à restaurer lors d'un exercice.

\textbf{Mises à niveau.} Une version sort tous les six mois, mais on n'est pas obligé de les suivre toutes. Une version sur deux,
la première de l'année (2025.1, 2026.1\dots), est une version \gls{slurp} : on peut passer directement de l'une à la suivante, alors
qu'une version intermédiaire ne se met à niveau que vers la version qui la suit \parencite{openstack-slurp} (figure~\ref{fig:dep-slurp}).

\begin{figure}[!ht]
\centering
\adjustbox{max width=\linewidth}{%
\begin{tikzpicture}[>=Stealth, font=\footnotesize, adj/.style={-{Stealth[length=2mm]}, thick, insagris}]
  \foreach \x/\v/\n/\k [count=\i] in {1.5/2025.1/{}/psvc, 4.5/2025.2/{}/pgris, 7.5/2026.1/{Gazpacho}/psvc, 10.5/2026.2/{Hibiscus}/pgris, 13.5/2027.1/{Indri}/psvc}
    { \node[\k, minimum width=2.5cm, minimum height=1.1cm] (v\i) at (\x,1.0) {\textbf{\v}\\ {\scriptsize\mdseries \n}}; }
  \node[font=\scriptsize] at (1.5,0.2) {SLURP};
  \node[font=\scriptsize] at (7.5,0.2) {SLURP};
  \node[font=\scriptsize] at (13.5,0.2) {SLURP, en développement};
  \foreach \a/\b in {1/2,2/3,3/4,4/5} \draw[adj] (v\a.east) -- (v\b.west);
  \draw[fluxo] (v1.north) to[out=65,in=115] node[etiqf, midway, above] {saut SLURP $\to$ SLURP} (v3.north);
  \draw[fluxo] (v3.north) to[out=65,in=115] node[etiqf, midway, above] {saut SLURP $\to$ SLURP} (v5.north);
  \node[callout, minimum width=15.6cm] at (7.5,-0.8)
    {Depuis 2025.2, il faut passer par 2026.1 : une version non SLURP ne se met à niveau que vers la suivante.};
\end{tikzpicture}}
\caption{Chemins de mise à niveau (flèches grises : version suivante, toujours possible). Les noms des versions
2026.1, 2026.2 et 2027.1 sont ceux de la documentation \parencite{openstack-releases}.}
\label{fig:dep-slurp}
\end{figure}

On met toujours le plan de contrôle à niveau avant les nœuds de calcul, que l'on traite par lots : services anciens et récents
doivent coexister un temps. Nova impose un ordre (Placement d'abord, puis \texttt{nova-conductor} en premier et \texttt{nova-api} en
dernier), puis des migrations de données « en ligne » (\texttt{nova-manage db online\_data\_migrations}) \parencite{nova-upgrades}.
Kolla-Ansible enchaîne tout cela dans une seule commande \parencite{kolla-ansible-operating} :

\par\noindent\begin{minipage}{\linewidth}
\begin{lstlisting}[style=terminal]
$ cp -r /etc/kolla /etc/kolla.sauvegarde        # configuration d'abord
$ pip install --upgrade \
    git+https://opendev.org/openstack/kolla-ansible@stable/2026.2
$ kolla-ansible install-deps
$ kolla-mergepwd --old passwords.yml.old --new passwords.yml.new \
    --final /etc/kolla/passwords.yml
$ kolla-ansible pull && kolla-ansible prechecks
$ kolla-ansible upgrade
\end{lstlisting}
\end{minipage}\par

Avant cela, on lit les notes de version de la série visée, on met à jour l'inventaire et \texttt{globals.yml} d'après les exemples
fournis, et l'on fusionne les mots de passe avec \texttt{kolla-mergepwd}. Pour un saut SLURP, RabbitMQ ne supporte pas un saut de deux
versions : il faut d'abord le mettre à niveau vers une version intermédiaire. La documentation ne décrit pas de retour arrière :
on se protège par une sauvegarde de la base et de la configuration, et par un essai préalable sur un environnement de test
\parencite{kolla-ansible-operating}.

% ---------------------------------------------------------------------
\subsection{Exercices corrigés}\label{sec:dep-exercices}
% ---------------------------------------------------------------------

\begin{exercice}{Choisir un outil et une topologie}\label{ex:dep-choix}
Une PME de vingt utilisateurs veut une plateforme interne de développement sur quatre serveurs physiques. Une personne de l'équipe
maîtrise Ansible, personne ne connaît Kubernetes. Proposer un outil, une répartition des rôles et les points à prévoir avant
l'ouverture.
\end{exercice}
\begin{correction}
Pas de réponse unique ; l'argumentation compte. Production interne, compétence Ansible, aucun Kubernetes : on écarte DevStack,
Packstack, l'installation manuelle et les approches fondées sur Kubernetes, et l'on retient \textbf{Kolla-Ansible} ou
\textbf{OpenStack-Ansible}. Trois serveurs forment les contrôleurs (quorum de MariaDB et de RabbitMQ) et portent aussi le réseau et
le calcul ; le quatrième est un nœud de calcul. Les machines des utilisateurs partagent alors les serveurs avec le plan de
contrôle : acceptable à cette échelle, à surveiller. À prévoir : un hôte de déploiement distinct, un registre d'images local, une
adresse virtuelle inutilisée sur le réseau de gestion, TLS, la sauvegarde de MariaDB et de la configuration, une mise à niveau
annuelle sur les versions SLURP.
\end{correction}

\begin{exercice}{Relire une configuration}\label{ex:dep-relire}
Cet extrait d'un déploiement Kolla-Ansible contient trois erreurs de conception ; lesquelles ? On sait que \texttt{node1} a l'adresse
\texttt{10.1.0.11} sur \texttt{eth0}.

\begin{lstlisting}
[control]
node1
node2
[network]
node1
# globals.yml
kolla_internal_vip_address: "10.1.0.11"
network_interface: "eth0"
neutron_external_interface: "eth0"
\end{lstlisting}
\end{exercice}
\begin{correction}
(1) \textbf{Deux contrôleurs} n'apportent aucune haute disponibilité : la panne de l'un fait perdre le quorum ; il en faut trois.
(2) L'\textbf{adresse virtuelle est celle de \texttt{node1}} : elle doit être inutilisée sur le réseau de gestion, car keepalived la
porte et la déplace. (3) L'\textbf{interface externe est celle de gestion} : \texttt{neutron\_external\_interface} doit être une
interface dédiée, active et sans adresse IP, qui porte \texttt{br-ex}. Au passage, le groupe \texttt{network} se réduit à
\texttt{node1} : HAProxy et keepalived, qui suivent ce groupe, n'ont alors aucun nœud de secours.
\end{correction}

\begin{exercice}{Ajouter un service à un cloud en production}\label{ex:dep-service}
Le cloud, déployé avec Kolla-Ansible, doit fournir des répartiteurs de charge (chapitre~\ref{chap:octavia}). Quelles modifications et
quelles commandes faut-il ? Que reste-t-il à faire hors de l'outil ?
\end{exercice}
\begin{correction}
\begin{lstlisting}[style=terminal]
$ vi /etc/kolla/globals.yml                # enable_octavia: true
$ kolla-ansible octavia-certificates
$ kolla-ansible prechecks
$ kolla-ansible deploy
\end{lstlisting}
\texttt{deploy} est rejouable : il applique l'état décrit par les variables, sans repartir de zéro. Les certificats propres à
Octavia sont indispensables. Hors de l'outil, il reste à fournir ce que le chapitre~\ref{chap:octavia} décrit : l'image
d'amphore, le réseau et le groupe de sécurité de gestion, la \emph{flavor} des amphores.
\end{correction}

\begin{exercice}{Planifier une mise à niveau}\label{ex:dep-niveau}
Un cloud tourne en 2025.2 et l'équipe vise 2026.2. Peut-elle y passer directement ? Décrire la première mise à niveau avec
Kolla-Ansible.
\end{exercice}
\begin{correction}
Non : 2025.2 n'est pas une version SLURP, et ne se met à niveau que vers la suivante, 2026.1 ; de là, 2026.2 est la version
suivante. Il faut donc \textbf{deux mises à niveau successives}, chacune suivie d'une validation. Pour la première : lire les notes
de version de 2026.1 ; sauvegarder \texttt{/etc/kolla} et la base ; installer la branche \texttt{stable/2026.1} de
\texttt{kolla-ansible} ; \texttt{kolla-ansible install-deps} ; mettre à jour l'inventaire et \texttt{globals.yml} d'après les
exemples et fusionner les mots de passe avec \texttt{kolla-mergepwd} ; \texttt{kolla-ansible pull} puis \texttt{prechecks} ; enfin
\texttt{kolla-ansible upgrade}, qui recrée les conteneurs et migre les schémas. Aucun retour arrière n'est documenté : la
sauvegarde et l'essai sur un environnement de test sont la seule protection.
\end{correction}

% ---------------------------------------------------------------------
\subsection{Récapitulatif du chapitre}\label{sec:dep-recap}
% ---------------------------------------------------------------------

\begin{aretenir}
\begin{itemize}
  \item Un déploiement assemble des services, une infrastructure partagée (base, file de messages, cache, répartiteur) et des
        machines en rôles ; l'outil de déploiement automatise toute la pile.
  \item Un nombre impair de contrôleurs (trois au moins) pour le quorum ; séparer au moins les réseaux de gestion, de tunnel et
        externe, et garder le réseau de gestion interne.
  \item DevStack (sources, machine jetable) : développement et apprentissage. Kolla-Ansible (images de conteneurs, à construire
        soi-même en production) et OpenStack-Ansible (sources, conteneurs LXC) : production.
  \item TripleO est retiré ; Kolla-Ansible ne déploie plus Swift et ne déploie pas Ceph.
  \item L'exploitation compte autant que l'installation : haute disponibilité, TLS, secrets, sauvegardes, supervision, mises à
        niveau (SLURP : une version sur deux, avec saut direct).
\end{itemize}
\end{aretenir}

\begin{center}
\renewcommand{\arraystretch}{1.3}
\begin{tabular}{@{}>{\raggedright\arraybackslash}p{4.0cm}>{\raggedright\arraybackslash}p{9.8cm}r@{}}
\toprule
\textbf{Tâche} & \textbf{Commande} & \textbf{§} \\
\midrule
Préparer hôtes et secrets & \texttt{kolla-ansible bootstrap-servers}, \texttt{kolla-genpwd} & \ref{sec:dep-kolla} \\
Vérifier puis déployer & \texttt{kolla-ansible prechecks}, \texttt{pull}, \texttt{deploy} & \ref{sec:dep-kolla} \\
Appliquer une modification & \texttt{kolla-ansible reconfigure [-t SERVICE]} & \ref{sec:dep-kolla} \\
Sauvegarder MariaDB & \texttt{kolla-ansible mariadb-backup [--incremental]} & \ref{sec:dep-exploitation} \\
Mettre à niveau & \texttt{kolla-ansible upgrade} & \ref{sec:dep-exploitation} \\
Contrôler les services & \texttt{openstack compute service list}, \texttt{network agent list} & \ref{sec:dep-exploitation} \\
\bottomrule
\end{tabular}
\end{center}

\begin{savoirfaire}
\begin{itemize}
  \item Choisir un outil et une topologie à partir d'un contexte, en justifiant (exercice~\ref{ex:dep-choix}).
  \item Relire un inventaire et un \texttt{globals.yml} pour y repérer les erreurs de quorum, d'adresse virtuelle et d'interfaces
        (exercice~\ref{ex:dep-relire}).
  \item Ajouter un service à un cloud Kolla-Ansible et nommer ce qui reste à faire hors de l'outil (exercice~\ref{ex:dep-service}).
  \item Planifier un chemin de mise à niveau avec les versions SLURP (exercice~\ref{ex:dep-niveau}).
\end{itemize}
\end{savoirfaire}
.
%  Il ne contient ni préambule ni \begin{document} : les environnements
%  (definition, resultat, remarque, pointcle, aretenir, savoirfaire, exercice,
%  correction), les styles TikZ (pcli, pfoc, psvc, pdat, pgris, fluxo, fluxr,
%  fluxd, etiqf, callout, num, acteur, vie, msg, rep) et les couleurs (insarouge,
%  insagris, insaorange, insapeche) sont définis dans main.tex. Un chapitre = une \section.
%  Sources : support « Méthodes de déploiement » (chapitre 4) et documentation
%  officielle de DevStack, Kolla, Kolla-Ansible, OpenStack-Ansible, OpenStack-Helm,
%  du guide d'installation, des guides de haute disponibilité et de mise à niveau
%  (voir sources.bib). Versions décrites : celles de la documentation en vigueur
%  début octobre 2026 (version 2026.2 « Hibiscus » sortie le 30 septembre 2026).
%  Écarts avec le support : TripleO est retiré du projet (dépôt archivé en 2024) ;
%  Kolla-Ansible ne déploie plus Swift et ne déploie pas Ceph ; les images publiées par
%  Kolla ne sont pas destinées à la production ; la documentation de Kolla-Ansible ne
%  décrit pas de retour arrière après une mise à niveau.
%  Chapitre détaillé : 15 pages au plus.
% =====================================================================

\part{Mettre en œuvre}\label{part:miseenoeuvre}
\partdesc{Passer de la théorie à un cloud en fonctionnement : choix de l'outil de déploiement, topologie, installation, exploitation et mises à niveau.}

% =====================================================================
\section{Méthodes de déploiement d'OpenStack}\label{chap:deploiement}
% =====================================================================

Les chapitres précédents ont présenté chaque service du point de vue de l'utilisateur, puis de son architecture. Celui-ci se place
du point de vue de l'administrateur, qui doit installer la plateforme et la faire vivre. Il explique pourquoi ce travail est délicat,
décrit ce que tout déploiement assemble, présente les outils les plus répandus (DevStack, Kolla-Ansible, OpenStack-Ansible,
installation manuelle, approches fondées sur Kubernetes), aide à choisir, puis termine par l'exploitation d'un cloud installé :
haute disponibilité, sécurité, mises à niveau.

% ---------------------------------------------------------------------
\subsection{Pourquoi déployer OpenStack est délicat}\label{sec:dep-pourquoi}
% ---------------------------------------------------------------------

OpenStack n'est pas un logiciel unique mais une constellation de services (chapitres~\ref{chap:keystone} à~\ref{chap:trove}) :
chacun a son code, ses processus, son fichier de configuration et, le plus souvent, sa propre base de données. Tous reposent sur la
même infrastructure : une base SQL, une file de messages et un cache (\gls{memcached}) \parencite{install-env}. Un cloud réel
s'étend de plus sur plusieurs machines aux rôles différents, dont les versions doivent rester cohérentes entre elles. Tout installer,
configurer et mettre à jour à la main devient vite ingérable au-delà de quelques nœuds : c'est ce que les outils de déploiement
automatisent (figure~\ref{fig:dep-pile}).

\begin{definition}{Outil de déploiement}
Logiciel qui installe, configure et relie les services d'un cloud sur un ensemble de machines à partir d'une description
(inventaire, fichier de variables, modèles), de façon reproductible ; la plupart savent aussi le mettre à jour.
\end{definition}

\begin{figure}[!ht]
\centering
\adjustbox{max width=\linewidth}{%
\begin{tikzpicture}[>=Stealth, font=\footnotesize]
  \node[pcli, minimum width=9.6cm] (u) at (5.0,7.0) {Utilisateurs \ $\cdot$ \ ligne de commande \ $\cdot$ \ Horizon \ $\cdot$ \ API REST};
  \node[pfoc, minimum width=9.6cm, minimum height=1.0cm] (s) at (5.0,5.4)
    {Services OpenStack\\ {\scriptsize\mdseries Keystone, Nova, Neutron, Glance, Cinder, Octavia, Heat\dots}};
  \node[psvc, minimum width=9.6cm, minimum height=1.0cm] (i) at (5.0,3.8)
    {Services d'infrastructure partagés\\ {\scriptsize\mdseries MariaDB \ $\cdot$ \ RabbitMQ \ $\cdot$ \ memcached \ $\cdot$ \ HAProxy}};
  \node[pgris, minimum width=9.6cm, minimum height=1.0cm] (e) at (5.0,2.2)
    {Mode d'exécution\\ {\scriptsize\mdseries paquets \ $\cdot$ \ environnements Python \ $\cdot$ \ conteneurs \ $\cdot$ \ Kubernetes}};
  \node[pgris, minimum width=9.6cm, minimum height=1.0cm] (m) at (5.0,0.6)
    {Machines et réseaux\\ {\scriptsize\mdseries contrôleurs, calcul, stockage ; réseaux de gestion, de tunnel et externe}};
  \node[draw=insarouge, thick, rounded corners=4pt, fill=insarouge!8, minimum width=1.7cm, minimum height=6.0cm] (o) at (11.9,3.0) {};
  \node[rotate=90, font=\footnotesize\bfseries, text=insarouge] at (11.9,3.0) {Outil de déploiement};
  \foreach \b in {s,i,e,m} \draw[fluxr] (o.west |- \b.east) -- (\b.east);
  \node[callout, minimum width=13.2cm] at (6.3,-1.0)
    {L'outil de déploiement automatise toute la pile, pas seulement l'installation des services.};
\end{tikzpicture}}
\caption{Ce qu'un déploiement doit assembler. Les services (chapitres~\ref{chap:keystone} à~\ref{chap:trove}) reposent sur une
infrastructure partagée, installée sur des machines selon un mode d'exécution au choix ; l'outil de déploiement prend en charge
les quatre étages.}
\label{fig:dep-pile}
\end{figure}

Avant d'installer quoi que ce soit, il faut trancher les questions du tableau~\ref{tab:dep-decisions}. La version compte
autant que l'outil : OpenStack publie une version tous les six mois, la plus récente (2026.2 \emph{Hibiscus}) est sortie le
30~septembre 2026 et 2027.1 \emph{Indri} est en développement \parencite{openstack-releases}. Les outils sont eux-mêmes
versionnés par série, et l'on installe la branche de l'outil qui correspond à la version visée \parencite{kolla-ansible-operating}.

\begin{table}[!ht]
\centering
\renewcommand{\arraystretch}{1.3}
\begin{tabular}{@{}>{\raggedright\arraybackslash}p{3.0cm}>{\raggedright\arraybackslash}p{6.3cm}>{\raggedright\arraybackslash}p{6.1cm}@{}}
\toprule
\textbf{Décision} & \textbf{Question à trancher} & \textbf{Exemples de réponses} \\
\midrule
Topologie & Quels rôles, sur combien de nœuds ? & Un nœud pour l'essai ; trois contrôleurs pour la production. \\
Réseaux & Quelles interfaces pour la gestion, les tunnels, l'accès externe ? & Une interface par réseau, ou des VLAN sur un lien. \\
Stockage & Où vivent les volumes, les images, les objets ? & Disques locaux (LVM), cluster Ceph, NFS. \\
Mode d'installation & Paquets, sources, conteneurs ou Kubernetes ? & Voir le tableau~\ref{tab:dep-outils}. \\
Version & Quelle version, à quel rythme la mettre à niveau ? & Toutes les versions (six mois) ou une sur deux (\gls{slurp}). \\
Exploitation & Haute disponibilité, TLS, secrets, sauvegardes, supervision ? & Section~\ref{sec:dep-exploitation}. \\
\bottomrule
\end{tabular}
\caption{Les décisions à prendre avant de déployer.}
\label{tab:dep-decisions}
\end{table}

% ---------------------------------------------------------------------
\subsection{Anatomie d'un déploiement}\label{sec:dep-anatomie}
% ---------------------------------------------------------------------

Quelle que soit la méthode, on retrouve les mêmes rôles de nœuds, les mêmes services d'infrastructure et les mêmes réseaux
(figure~\ref{fig:dep-topo}).

\begin{itemize}
  \item \textbf{Hôte de déploiement} : la machine d'où l'on lance l'outil (\gls{ansible}, le plus souvent) ; en production, une machine
        distincte du cloud \parencite{osa-deploy-host,kolla-ansible-production}.
  \item \textbf{Contrôleurs} : API des services, ordonnanceurs, base de données et file de messages ; un nombre impair, pour
        établir un quorum \parencite{kolla-ansible-production}.
  \item \textbf{Réseau} : agents Neutron, routeurs ; avec Kolla-Ansible, HAProxy et keepalived s'y exécutent aussi
        \parencite{kolla-ansible-production}.
  \item \textbf{Calcul} : hyperviseur et \texttt{nova-compute}, là où vivent les machines virtuelles.
  \item \textbf{Stockage} : \texttt{cinder-volume} et ses disques ; les volumes peuvent aussi résider sur un cluster \gls{ceph} externe.
\end{itemize}

Le guide d'installation décrit une architecture minimale : un contrôleur, qui exécute l'identité, les images, Placement, les parties
de gestion de Nova et de Neutron, des agents réseau et le tableau de bord, et un nœud de calcul, avec en option des nœuds de stockage
bloc et objet. Il précise qu'elle n'est pas destinée à la production : les agents réseau y résident sur le contrôleur et le trafic
des tunnels emprunte le réseau de gestion \parencite{install-overview}.

\begin{figure}[!ht]
\centering
\adjustbox{max width=\linewidth}{%
\begin{tikzpicture}[>=Stealth, font=\footnotesize,
    nd/.style={draw=insaorange, fill=insapeche, thick, rounded corners=3pt, text width=3.3cm, align=center,
               inner sep=4pt, minimum height=2.3cm}]
  \node[pcli, minimum width=5.2cm] (dh) at (8.4,3.5) {Hôte de déploiement\\ {\scriptsize\mdseries Ansible, inventaire, variables}};
  \node[nd] (c) at (1.8,0)  {\textbf{Contrôleurs} {\scriptsize$\times 3$}\\[2pt] {\scriptsize API des services, ordonnanceurs, MariaDB, RabbitMQ, memcached}};
  \node[nd] (r) at (6.0,0)  {\textbf{Réseau}\\[2pt] {\scriptsize agents Neutron, routeurs, HAProxy et keepalived}};
  \node[nd] (k) at (10.2,0) {\textbf{Calcul} {\scriptsize$\times N$}\\[2pt] {\scriptsize hyperviseur, \texttt{nova-compute}, agent du réseau virtuel}};
  \node[nd] (s) at (14.4,0) {\textbf{Stockage}\\[2pt] {\scriptsize \texttt{cinder-volume}, disques locaux ou accès à Ceph}};
  \foreach \n in {c,r,k,s} \draw[fluxd] (dh.south) -- ++(0,-0.4) -| (\n.north);
  \node[etiqf] at (4.6,2.775) {Ansible, par SSH};
  % réseaux
  \foreach \y/\col/\tc/\nom in {-2.5/insarouge/insarouge/Gestion, -3.3/insaorange/insaorange/Tunnel,
      -4.1/insarouge!45/insarouge/Stockage, -4.9/insagris/insagris/Externe}
    { \draw[line width=2pt, \col] (0.2,\y) -- (16.1,\y);
      \node[anchor=east, font=\scriptsize\bfseries, text=\tc] at (-0.1,\y) {\nom}; }
  % branchements : décalage selon le réseau
  \foreach \cx/\y/\off/\col in {
      1.8/-2.5/-1.2/insarouge, 1.8/-4.9/1.2/insagris,
      6.0/-2.5/-1.2/insarouge, 6.0/-3.3/-0.4/insaorange, 6.0/-4.9/1.2/insagris,
      10.2/-2.5/-1.2/insarouge, 10.2/-3.3/-0.4/insaorange, 10.2/-4.1/0.4/insarouge!45,
      14.4/-2.5/-1.2/insarouge, 14.4/-4.1/0.4/insarouge!45}
    { \draw[thick, \col] ({\cx+\off},-1.15) -- ({\cx+\off},\y);
      \fill[\col] ({\cx+\off},\y) circle (2.6pt); }
  \node[callout, minimum width=15.2cm] at (8.2,-5.9)
    {Un même nœud peut tenir plusieurs rôles : l'inventaire de l'outil les décrit, le logiciel ne les fige pas.};
\end{tikzpicture}}
\caption{Topologie d'un déploiement multi-nœuds. L'hôte de déploiement configure les nœuds par SSH (pointillés) ; chaque rôle se
branche (points) sur les réseaux dont il a besoin.}
\label{fig:dep-topo}
\end{figure}

Les services d'infrastructure demandent une attention particulière. La base de données est le plus souvent un cluster MariaDB avec
\gls{galera} (multimaître), la file de messages un cluster \gls{rabbitmq}, et \gls{haproxy} répartit les requêtes des API vers les contrôleurs
\parencite{osa-service-arch}. Ces composants à état imposent un \gls{quorum} : deux instances n'apportent aucune haute
disponibilité (la panne de l'une paralyse l'autre), trois en tolèrent une en panne, cinq en tolèrent deux, et au-delà le gain
disparaît \parencite{kolla-ansible-production}. L'horloge de tous les nœuds doit en outre être synchronisée (\gls{ntp})
\parencite{install-env}.

Les réseaux se séparent par leur rôle (tableau~\ref{tab:dep-reseaux}) ; le trafic de gestion, qui porte les mots de passe et les
messages entre services, doit rester interne \parencite{kolla-ansible-production}.

\begin{table}[!ht]
\centering
\renewcommand{\arraystretch}{1.3}
\begin{tabular}{@{}>{\raggedright\arraybackslash}p{2.4cm}>{\raggedright\arraybackslash}p{7.6cm}>{\raggedright\arraybackslash}p{5.5cm}@{}}
\toprule
\textbf{Réseau} & \textbf{Rôle} & \textbf{Variable Kolla-Ansible} \\
\midrule
Gestion & Dialogue des services entre eux, avec la base et la file ; à garder interne. & \texttt{api\_\allowbreak interface} (défaut : \texttt{network\_\allowbreak interface}) \\
Tunnel & Trafic des réseaux de projets encapsulé (VXLAN, Geneve). & \texttt{tunnel\_\allowbreak interface} \\
Externe & Réseaux externes et adresses flottantes ; l'interface porte \texttt{br-ex} et n'a pas d'adresse IP. & \texttt{neutron\_\allowbreak external\_\allowbreak interface} \\
API publique & Points d'accès exposés aux utilisateurs derrière HAProxy. & \texttt{kolla\_\allowbreak external\_\allowbreak vip\_\allowbreak interface} \\
\bottomrule
\end{tabular}
\caption{Les réseaux d'un déploiement et leurs interfaces dans Kolla-Ansible \parencite{kolla-ansible-production}.}
\label{tab:dep-reseaux}
\end{table}

% ---------------------------------------------------------------------
\subsection{Panorama des méthodes}\label{sec:dep-panorama}
% ---------------------------------------------------------------------

Les outils se distinguent d'abord par la manière dont ils installent les services (figure~\ref{fig:dep-carte}), puis par leur public
(tableau~\ref{tab:dep-outils}). Aucun n'est « le meilleur » : chacun répond à un contexte.

\begin{figure}[!ht]
\centering
\adjustbox{max width=\linewidth}{%
\begin{tikzpicture}[>=Stealth, font=\footnotesize]
  \foreach \x/\t in {1.9/{Paquets de la\\ distribution}, 6.3/{Sources Git}, 10.7/{Images de\\ conteneurs}, 15.1/{Kubernetes}}
    \node[pgris, minimum width=3.9cm, minimum height=1.0cm] at (\x,4.6) {\t};
  \node[pdat, minimum width=3.7cm]  at (1.9,3.1)  {Installation manuelle};
  \node[pdat, minimum width=3.7cm]  at (1.9,2.2)  {Packstack (RDO)};
  \node[pdat, minimum width=3.7cm]  at (6.3,3.1)  {DevStack};
  \node[psvc, minimum width=3.7cm]  at (6.3,2.2)  {OpenStack-Ansible};
  \node[psvc, minimum width=3.7cm]  at (10.7,3.1) {Kolla-Ansible};
  \node[psvc, minimum width=3.7cm]  at (10.7,2.2) {Kayobe};
  \node[psvc, minimum width=3.7cm]  at (15.1,3.1) {OpenStack-Helm};
  \node[psvc, minimum width=3.7cm]  at (15.1,2.2) {Sunbeam (Canonical)};
  \node[psvc, minimum width=3.7cm]  at (15.1,1.3) {RHOSO (Red Hat)};
  \node[pdat, minimum width=0.5cm, inner xsep=4pt] at (3.2,0.0) {};
  \node[anchor=west] at (3.55,0.0) {essai, apprentissage, développement};
  \node[psvc, minimum width=0.5cm, inner xsep=4pt] at (9.4,0.0) {};
  \node[anchor=west] at (9.75,0.0) {production};
  \node[callout, minimum width=16.4cm] at (8.5,-1.2)
    {Ce qui change d'un outil à l'autre n'est pas la liste des services, mais la façon dont ils sont installés et mis à jour.};
\end{tikzpicture}}
\caption{Les outils de déploiement selon le mode d'installation des services. OpenStack-Ansible installe depuis les sources par
défaut (paquets possibles pour certains services) et peut isoler les services dans des conteneurs LXC \parencite{osa-config}.}
\label{fig:dep-carte}
\end{figure}

\begin{table}[!ht]
\centering
\renewcommand{\arraystretch}{1.25}
\small
\begin{tabular}{@{}>{\raggedright\arraybackslash}p{2.5cm}>{\raggedright\arraybackslash}p{5.5cm}>{\raggedright\arraybackslash}p{3.4cm}>{\raggedright\arraybackslash}p{3.8cm}@{}}
\toprule
\textbf{Outil} & \textbf{Principe} & \textbf{Usage visé} & \textbf{Ressources minimales documentées} \\
\midrule
Installation manuelle & Paquets de la distribution, guide d'installation officiel, service par service & Apprentissage, environnement isolé & Contrôleur : 4~Go ; calcul : 2~Go (preuve de concept) \parencite{install-env} \\
\gls{packstack} & Installeur du projet RDO (CentOS Stream) ; haute disponibilité et mises à niveau hors périmètre & Preuve de concept & 16~Go de mémoire \parencite{packstack} \\
DevStack & Scripts shell, installation depuis les sources Git & Développement, tests, démonstration & Machine dédiée (non précisé) \parencite{devstack-index} \\
OpenStack-Ansible & Ansible ; sources dans des environnements virtuels Python, conteneurs LXC & Production & Test (AIO) : 8~vCPU, 8~Go, 50~Go \parencite{osa-aio} \\
Kolla-Ansible & Ansible ; images de conteneurs Docker ou Podman & Production, évaluation & 2~interfaces, 8~Go, 40~Go \parencite{kolla-ansible-quickstart} \\
Kayobe & Kolla-Ansible, plus configuration des hôtes et provisionnement du matériel (Bifrost) & Production, depuis le serveur nu & --- \parencite{kayobe-index} \\
OpenStack-Helm & Charts Helm sur Kubernetes & Cloud sur Kubernetes & Kubernetes 1.33 à 1.35 \parencite{helm-readme} \\
Sunbeam & Snap, Juju, Kubernetes (Canonical OpenStack) \parencite{sunbeam-multinode} & Petits clouds, de un à plusieurs nœuds & Un nœud : 4~cœurs, 16~Gio, 100~Gio \parencite{sunbeam-tuto} \\
RHOSO & Opérateurs sur OpenShift (Red Hat) & Clients de Red Hat & --- \parencite{rhoso-deploy} \\
TripleO & Retiré du projet (dépôt archivé en 2024) & --- & --- \parencite{tripleo-docs} \\
\bottomrule
\end{tabular}
\caption{Les principaux outils de déploiement d'OpenStack.}
\label{tab:dep-outils}
\end{table}

TripleO (« OpenStack sur OpenStack » : un cloud d'administration, l'\emph{undercloud}, déployait le cloud de production,
l'\emph{overcloud}) figure encore dans de nombreux supports ; il n'est plus maintenu. Red Hat, qui en tirait son outil
\emph{director}, déploie aujourd'hui son offre par des opérateurs sur OpenShift (RHOSO) \parencite{tripleo-docs,rhoso-deploy}. La suite
détaille DevStack (section~\ref{sec:dep-devstack}), Kolla-Ansible (section~\ref{sec:dep-kolla}) et OpenStack-Ansible
(section~\ref{sec:dep-osa}), revient sur l'installation manuelle (section~\ref{sec:dep-manuel}) pour comprendre ce que les outils
automatisent, situe les approches fondées sur Kubernetes (section~\ref{sec:dep-autres}), puis aide à choisir
(section~\ref{sec:dep-choix}) et à exploiter (section~\ref{sec:dep-exploitation}).

% ---------------------------------------------------------------------
\subsection{DevStack : l'environnement tout-en-un}\label{sec:dep-devstack}
% ---------------------------------------------------------------------

\begin{definition}{DevStack}
Série de scripts qui déploient rapidement un cloud OpenStack complet sur une seule machine, à partir des dépôts Git de chaque projet
\parencite{devstack-index}.
\end{definition}

\gls{devstack} sert à construire vite des environnements de développement, à décrire des configurations qui fonctionnent, à faciliter les
contributions et à fournir l'environnement des tests d'intégration de chaque modification proposée \parencite{devstack-source}. Il
installe depuis les sources, sur la branche \texttt{master} par défaut ou sur une branche \texttt{stable/<série>} pour une version plus
ancienne : c'est ce qui en fait l'outil idéal pour essayer une modification de code (dans Nova, par exemple) avant de la proposer. Il
demande une machine dédiée et jetable. Ubuntu 24.04 est la plus testée ; le projet suit aussi les deux dernières versions LTS
d'Ubuntu, Rocky Linux 9 et openEuler. L'installation dure de 15 à 30 minutes \parencite{devstack-index}.

\begin{figure}[!ht]
\centering
\adjustbox{max width=\linewidth}{%
\begin{tikzpicture}[>=Stealth, font=\footnotesize]
  \node[pcli, minimum width=3.6cm] (lc) at (1.9,2.4) {\texttt{local.conf}\\ {\scriptsize\mdseries mots de passe, services}};
  \node[pgris, minimum width=3.6cm] (git) at (1.9,0.3) {Dépôts Git des projets\\ {\scriptsize\mdseries opendev.org}};
  \node[pfoc, minimum width=2.6cm, minimum height=1.3cm] (st) at (7.4,1.35) {\texttt{stack.sh}};
  \node[pgris, minimum width=5.6cm] (pk) at (13.0,2.8) {Paquets de la distribution\\ {\scriptsize\mdseries MySQL, RabbitMQ, Apache}};
  \node[psvc, minimum width=5.6cm] (sv) at (13.0,1.35) {Services installés dans \texttt{/opt/stack}\\ {\scriptsize\mdseries fichiers de configuration, bases}};
  \node[psvc, minimum width=5.6cm] (sd) at (13.0,-0.1) {Unités systemd \texttt{devstack@service}};
  \node[pdat, minimum width=9.4cm] (cl) at (9.4,-1.9) {Cloud à un nœud : API, Horizon, fichier \texttt{openrc}};
  \draw[fluxr] (lc.east) -| (st.north) node[num, pos=0.25] {1};
  \draw[fluxr] (git.east) -| (st.south) node[num, pos=0.25] {2};
  \foreach \t in {pk,sv,sd} \draw[fluxo] (st.east) -- (\t.west);
  \node[num] at (9.7,1.35) {3};
  \draw[fluxo] (sd.south) -- (sd.south |- cl.north);
  \node[callout, minimum width=15.6cm] at (8.5,-3.3)
    {DevStack installe depuis les sources, sur une machine dédiée : pour apprendre et développer, pas pour héberger un service réel.};
\end{tikzpicture}}
\caption{Ce que fait \texttt{stack.sh}. Il lit \texttt{local.conf} (1) et les dépôts des projets (2), puis installe les paquets de la
distribution, installe et configure les services dans \texttt{/opt/stack} et les démarre comme unités systemd (3).}
\label{fig:dep-devstack}
\end{figure}

Le script \texttt{stack.sh} ne doit pas être lancé en tant que \texttt{root} : il s'exécute sous un compte dédié, \texttt{stack}, qui
dispose de \texttt{sudo} sans mot de passe \parencite{devstack-index}. Ses services de base (Keystone, Nova, Placement, Glance,
Cinder, Neutron avec OVN et Horizon) sont actifs par défaut ; les autres s'ajoutent avec \texttt{enable\_service} (c'est le cas de Swift)
ou, pour un projet qui n'est pas dans le dépôt de DevStack, avec \texttt{enable\_plugin} \parencite{devstack-source,devstack-config,devstack-plugins}.
Chaque service tourne en unité systemd, dont on lit le journal avec \texttt{journalctl} \parencite{devstack-systemd}.

\par\noindent\begin{minipage}{\linewidth}
\begin{lstlisting}[style=terminal]
$ sudo useradd -s /bin/bash -d /opt/stack -m stack
$ sudo chmod +x /opt/stack
$ echo "stack ALL=(ALL) NOPASSWD: ALL" | sudo tee /etc/sudoers.d/stack
$ sudo -u stack -i
$ git clone https://opendev.org/openstack/devstack
$ cd devstack && nano local.conf
$ ./stack.sh
$ sudo journalctl -f --unit devstack@n-cpu.service
\end{lstlisting}
\end{minipage}\par

\par\noindent\begin{minipage}{\linewidth}
\begin{lstlisting}
[[local|localrc]]
ADMIN_PASSWORD=secret
DATABASE_PASSWORD=$ADMIN_PASSWORD
RABBIT_PASSWORD=$ADMIN_PASSWORD
SERVICE_PASSWORD=$ADMIN_PASSWORD
HOST_IP=192.168.1.10
# ajouter un projet absent des services par défaut (ici Heat)
enable_plugin heat https://opendev.org/openstack/heat
# retirer un service dont on n'a pas besoin
disable_service horizon
\end{lstlisting}
\end{minipage}\par

Le fichier \texttt{local.conf} est le seul à personnaliser : il suit un format de type INI dont la section \texttt{[[local|localrc]]}
fixe les mots de passe (des caractères alphanumériques seulement), l'adresse de l'hôte, les services actifs, et même les dépôts et
les branches à cloner (\texttt{NOVA\_REPO}, \texttt{NOVA\_BRANCH}) \parencite{devstack-config}. Une fois le script terminé, on charge les
identifiants avec \texttt{. openrc} puis on utilise la ligne de commande \parencite{devstack-source}.

Pour plusieurs machines, le contrôleur reçoit un \texttt{local.conf} complet et chaque nœud de calcul un fichier réduit, qui
désigne le contrôleur (\texttt{SERVICE\_HOST}, \texttt{MYSQL\_HOST}, \texttt{RABBIT\_HOST}) et ne garde que des services
comme \texttt{n-cpu}. Il faut ensuite associer chaque nouvel hôte à la cellule, depuis le contrôleur, avec
\texttt{tools/discover\_hosts.sh} (section~\ref{sec:nova-archi}) \parencite{devstack-multinode}.

DevStack n'est pas un installateur généraliste : la plupart des combinaisons d'options sont rarement testées, et il modifie en
profondeur le système hôte \parencite{devstack-overview,devstack-source}. Conçu pour un nœud, il n'offre pas de haute disponibilité.
\texttt{./unstack.sh} arrête les services (en laissant MySQL et RabbitMQ), \texttt{./clean.sh} efface les données.

% ---------------------------------------------------------------------
\subsection{Kolla-Ansible : OpenStack en conteneurs}\label{sec:dep-kolla}
% ---------------------------------------------------------------------

Le projet \gls{kolla} a pour mission de fournir des outils qui créent des conteneurs prêts pour la production et des outils qui déploient
des clouds OpenStack \parencite{kolla-source}. Il faut distinguer ses livrables :

\begin{itemize}
  \item \textbf{Kolla} construit les images de conteneurs de chaque service (\texttt{kolla-build}) ;
  \item \textbf{Kolla-Ansible} fournit les playbooks Ansible qui déploient ces images sur les nœuds ;
  \item \textbf{\gls{kayobe}} ajoute la configuration des hôtes et le provisionnement du matériel (avec Bifrost, qui s'appuie sur
        Ironic) à Kolla-Ansible \parencite{kayobe-index}.
\end{itemize}

\gls{kollaansible} fonctionne avec Docker ou Podman. Il prend en charge comme hôtes Debian 13, Rocky Linux 10 et Ubuntu 24.04, et les
images existent dans les mêmes trois distributions (\texttt{kolla\_base\_distro}) \parencite{kolla-ansible-support}. Par défaut il
tire les images publiées sur Quay.io (espace \texttt{openstack.kolla}). Ses auteurs avertissent que ces images, destinées aux
essais et aux démonstrations, ne subissent aucun examen de sécurité : en production, on construit et on maintient son propre jeu
d'images avec \texttt{kolla-build}, que l'on place dans un \gls{registre} local (recommandé dès que le déploiement compte plusieurs nœuds)
\parencite{kolla-source,kolla-image-building,kolla-ansible-multinode}.

\begin{figure}[!ht]
\centering
\adjustbox{max width=\linewidth}{%
\begin{tikzpicture}[>=Stealth, font=\footnotesize,
    ct/.style={draw=insaorange, fill=insapeche, rounded corners=2pt, font=\scriptsize\ttfamily, minimum height=0.55cm, inner xsep=3pt}]
  \node[pcli, minimum width=7.2cm] (dh) at (3.9,7.1)
    {Hôte de déploiement\\ {\scriptsize\mdseries \texttt{kolla-ansible}, inventaire, \texttt{globals.yml}, \texttt{passwords.yml}}};
  \node[pgris, minimum width=5.6cm] (reg) at (12.2,7.1) {Registre d'images\\ {\scriptsize\mdseries Quay.io ou registre local}};
  \foreach \y/\nom in {5.0/{Contrôleurs $\times 3$}, 3.6/{Réseau}, 2.2/{Calcul}, 0.8/{Stockage}}
    { \draw[draw=insagris, dashed, rounded corners=3pt, fill=insagris!5] (0,{\y-0.58}) rectangle (15.7,{\y+0.58});
      \node[anchor=west, font=\footnotesize\bfseries] at (0.15,\y) {\nom}; }
  \foreach \x/\t in {3.9/keystone, 5.9/nova\_api, 8.5/neutron\_server, 10.9/horizon, 12.7/mariadb, 14.5/rabbitmq}
    \node[ct] at (\x,5.0) {\t};
  \foreach \x/\t in {3.9/haproxy, 6.0/keepalived, 8.9/neutron\_l3\_agent, 12.1/neutron\_dhcp\_agent}
    \node[ct] at (\x,3.6) {\t};
  \foreach \x/\t in {4.2/nova\_compute, 6.8/nova\_libvirt, 10.0/openvswitch\_vswitchd}
    \node[ct] at (\x,2.2) {\t};
  \node[ct] at (4.2,0.8) {cinder\_volume};
  \node[pcli, minimum width=1.0cm, inner xsep=4pt] (vip) at (14.8,3.6) {VIP};
  \node[pgris, dashed, minimum width=3.8cm] (ceph) at (11.8,0.8) {Ceph (externe)};
  \draw[fluxd] (6.0,0.8) -- (ceph.west);
  \draw[fluxd] (dh.south) -- (3.9,5.65) node[etiqf, midway, right=3pt] {Ansible, par SSH};
  \draw[fluxo] (reg.south) -- (12.2,5.65) node[etiqf, midway, right=3pt] {tirer les images};
  \node[callout, minimum width=15.4cm] at (7.85,-0.75)
    {Chaque service tourne dans son conteneur, avec ses dépendances figées dans l'image : pas de conflit de versions entre services.};
\end{tikzpicture}}
\caption{Kolla-Ansible sur plusieurs nœuds. L'hôte de déploiement pilote les nœuds par Ansible (pointillés) ; ceux-ci tirent leurs
images d'un registre. Les groupes de l'inventaire fixent les conteneurs de chaque nœud. HAProxy et keepalived, qui portent
l'adresse virtuelle (VIP), s'exécutent par défaut sur les nœuds réseau. Kolla-Ansible ne déploie pas Ceph : il s'y connecte.}
\label{fig:dep-kolla}
\end{figure}

La configuration tient en trois éléments, placés sur l'hôte de déploiement, qu'installe \texttt{pip} dans un environnement virtuel
\parencite{kolla-ansible-quickstart} :
\begin{itemize}
  \item l'\textbf{\gls{inventaireansible}} Ansible, qui affecte les machines aux groupes \texttt{control}, \texttt{network}, \texttt{compute},
        \texttt{monitoring} et \texttt{storage} (le groupe \texttt{control} compte un nombre impair de nœuds) ; chaque service a lui-même
        son groupe, mais en modifier la répartition peut casser le déploiement \parencite{kolla-ansible-multinode,kolla-ansible-production} ;
  \item \texttt{/etc/kolla/globals.yml}, le fichier principal ;
  \item \texttt{/etc/kolla/passwords.yml}, dont les mots de passe vides se remplissent avec \texttt{kolla-genpwd}.
\end{itemize}

\par\noindent\begin{minipage}{\linewidth}
\begin{lstlisting}[style=terminal]
$ python3 -m venv /opt/kolla-venv && source /opt/kolla-venv/bin/activate
$ pip install git+https://opendev.org/openstack/kolla-ansible@stable/2026.1
$ sudo mkdir -p /etc/kolla && sudo chown $USER:$USER /etc/kolla
$ cp -r /opt/kolla-venv/share/kolla-ansible/etc_examples/kolla/* /etc/kolla
$ cp /opt/kolla-venv/share/kolla-ansible/ansible/inventory/multinode .
$ kolla-ansible install-deps
$ kolla-genpwd
\end{lstlisting}
\end{minipage}\par

\par\noindent\begin{minipage}{\linewidth}
\begin{lstlisting}
# /etc/kolla/globals.yml
kolla_base_distro: "ubuntu"
network_interface: "eth0"
neutron_external_interface: "eth1"
kolla_internal_vip_address: "10.1.0.250"
enable_cinder: true
enable_octavia: true
\end{lstlisting}
\end{minipage}\par

La valeur de \texttt{kolla\_internal\_vip\_address} doit être une adresse inutilisée du réseau de gestion : keepalived la porte et la
déplace en cas de panne ; l'interface externe doit être active, mais sans adresse IP \parencite{kolla-ansible-quickstart}. Kolla-Ansible
est volontairement très directif : quelques variables suffisent à un déploiement fonctionnel, et il ne multiplie pas les options.
Pour tout le reste, les fichiers déposés dans \texttt{/etc/kolla/config/} sont fusionnés avec la configuration par défaut du service
(par exemple, \texttt{nova.conf} avec \texttt{virt\_type = qemu} pour faire tourner le cloud dans des machines virtuelles sans
virtualisation imbriquée) \parencite{kolla-ansible-philosophy}. Les services de base (Keystone, Glance, Nova avec Placement, Neutron,
Heat, Horizon) sont actifs par défaut ; on ajoute les autres avec les variables \texttt{enable\_*} (tableau~\ref{tab:dep-kolla-services})
\parencite{kolla-ansible-source}.

\begin{table}[!ht]
\centering
\renewcommand{\arraystretch}{1.3}
\begin{tabular}{@{}>{\raggedright\arraybackslash}p{2.6cm}>{\raggedright\arraybackslash}p{3.4cm}>{\raggedright\arraybackslash}p{9.5cm}@{}}
\toprule
\textbf{Service} & \textbf{Variable} & \textbf{À savoir} \\
\midrule
Cinder (chap.~\ref{chap:cinder}) & \texttt{enable\_cinder} & Inactif par défaut ; disques LVM (\texttt{enable\_cinder\_\allowbreak backend\_\allowbreak lvm}) ou cluster Ceph externe. \\
Octavia (chap.~\ref{chap:octavia}) & \texttt{enable\_octavia} & Demande des certificats propres (\texttt{kolla-ansible octavia-certificates}) \parencite{kolla-ansible-source}. \\
Ironic (chap.~\ref{chap:ironic}) & \texttt{enable\_ironic} & S'appuie sur la même base de déploiement que le reste. \\
Magnum (chap.~\ref{chap:magnum}) & \texttt{enable\_magnum} & Le pilote de cluster reste à choisir (section~\ref{sec:ma-pilotes}). \\
Trove (chap.~\ref{chap:trove}) & \texttt{enable\_trove} & Les images invitées restent à construire (section~\ref{sec:tr-datastores}). \\
Swift (chap.~\ref{chap:swift}) & --- & Plus déployé par Kolla-Ansible (retiré faute de mainteneurs) ; Glance peut utiliser un Swift externe \parencite{kolla-ansible-source}. \\
\bottomrule
\end{tabular}
\caption{Activer les autres services avec Kolla-Ansible.}
\label{tab:dep-kolla-services}
\end{table}

Le déploiement suit un enchaînement fixe (figure~\ref{fig:dep-kolla-flux} et tableau~\ref{tab:dep-kolla-cmd}). La commande
\texttt{deploy} peut être relancée autant de fois que nécessaire, mais si une tâche d'amorçage échoue, un nouvel appel ne la répare pas ;
le plus rapide est alors de détruire ce déploiement raté (\texttt{kolla-ansible destroy}) \parencite{kolla-ansible-troubleshooting}.

\begin{figure}[!ht]
\centering
\adjustbox{max width=\linewidth}{%
\begin{tikzpicture}[>=Stealth, font=\footnotesize]
  \foreach \x/\c/\d [count=\i] in {1.6/bootstrap-servers/{prépare les hôtes}, 4.8/prechecks/{vérifie les prérequis}, 8.0/pull/{tire les images},
                                   11.2/deploy/{démarre les conteneurs}, 14.4/post-deploy/{écrit \texttt{clouds.yaml}}}
    { \node[pfoc, minimum width=2.9cm, font=\scriptsize\bfseries\ttfamily] (s\i) at (\x,2.0) {\c};
      \node[num] at (\x-1.25,2.55) {\i};
      \node[etiqf] at (\x,1.3) {\d}; }
  \foreach \a/\b in {1/2,2/3,3/4,4/5} \draw[fluxr] (s\a.east) -- (s\b.west);
  \node[anchor=west, font=\footnotesize\bfseries, text=insagris] at (0,-0.1) {Ensuite :};
  \foreach \x/\c in {3.6/reconfigure, 7.0/upgrade, 10.4/mariadb-backup, 13.8/destroy}
    \node[psvc, minimum width=2.9cm, font=\scriptsize\bfseries\ttfamily] at (\x,-0.1) {\c};
  \node[callout, minimum width=15.6cm] at (8.0,-1.4)
    {Les mêmes commandes servent à installer, à reconfigurer et à mettre à niveau : l'inventaire et les variables restent la seule description du cloud.};
\end{tikzpicture}}
\caption{Les cinq étapes d'un déploiement avec Kolla-Ansible, puis les commandes d'exploitation.}
\label{fig:dep-kolla-flux}
\end{figure}

\begin{table}[!ht]
\centering
\renewcommand{\arraystretch}{1.25}
\begin{tabular}{@{}>{\raggedright\arraybackslash\ttfamily\small}p{3.3cm}>{\raggedright\arraybackslash}p{12.2cm}@{}}
\toprule
\textrm{\textbf{Commande}} & \textbf{Effet} \\
\midrule
bootstrap-servers & Prépare chaque hôte : paquets système, moteur de conteneurs, utilisateurs \parencite{kolla-ansible-quickstart}. \\
prechecks & Vérifie que les hôtes sont dans un état où l'on peut déployer \parencite{kolla-ansible-troubleshooting}. \\
pull & Tire les images sur les hôtes, sans toucher aux conteneurs. \\
deploy & Génère les configurations, amorce les services et démarre tous les conteneurs. \\
post-deploy & Écrit les accès administrateur dans \texttt{/etc/kolla/clouds.yaml} \parencite{kolla-ansible-quickstart}. \\
reconfigure & Applique une modification de \texttt{globals.yml} ou de \texttt{config/}, éventuellement pour un seul service (\texttt{-t mariadb}). \\
validate-config & Valide les fichiers de configuration générés (après un premier déploiement). \\
upgrade & Met à niveau d'une série à la suivante : nouvelles images, migrations de schéma \parencite{kolla-ansible-operating}. \\
mariadb-backup & Sauvegarde complète ou incrémentale (\texttt{--incremental}) de MariaDB \parencite{kolla-ansible-mariadb-backup}. \\
destroy & Détruit les conteneurs, les volumes et la configuration des hôtes. \\
\bottomrule
\end{tabular}
\caption{Les commandes de \texttt{kolla-ansible} \parencite{kolla-ansible-operating,kolla-ansible-source}.}
\label{tab:dep-kolla-cmd}
\end{table}

% ---------------------------------------------------------------------
\subsection{OpenStack-Ansible : sources et conteneurs LXC}\label{sec:dep-osa}
% ---------------------------------------------------------------------

\gls{openstackansible} (OSA) déploie OpenStack avec Ansible, lui aussi pour la production. Par défaut, tous les services sont installés
depuis les sources dans des \glspl{venv} Python, puis isolés dans des conteneurs \gls{lxc}. Ce sont des conteneurs « machine »,
proches de petites machines virtuelles, et non des conteneurs d'application comme ceux de Docker ; ils sont facultatifs, et un
déploiement directement sur les hôtes est possible \parencite{osa-about}. Les hôtes pris en charge sont Ubuntu 24.04 et 26.04, Debian 13,
CentOS Stream 10 et Rocky Linux 10 \parencite{osa-about}.

\begin{figure}[!ht]
\centering
\adjustbox{max width=\linewidth}{%
\begin{tikzpicture}[>=Stealth, font=\footnotesize,
    ct/.style={draw=insaorange, fill=insapeche, rounded corners=2pt, font=\scriptsize\ttfamily, minimum height=0.6cm, minimum width=2.1cm},
    lx/.style={draw=insarouge!80!black, fill=insarouge!12, rounded corners=3pt, text width=1.8cm, align=center, font=\scriptsize,
               minimum height=1.35cm, inner sep=2pt}]
  \node[pcli, minimum width=3.4cm] at (3.7,5.4) {Kolla-Ansible};
  \node[pcli, minimum width=3.4cm] at (12.0,5.4) {OpenStack-Ansible};
  \draw[draw=insagris, dashed, rounded corners=3pt, fill=insagris!5] (0,0.6) rectangle (7.4,4.7);
  \node[anchor=north west, font=\footnotesize\bfseries] at (0.15,4.6) {Hôte : Docker ou Podman};
  \node[ct] at (1.4,2.9) {keystone};
  \node[ct] at (3.7,2.9) {nova\_api};
  \node[ct] at (6.0,2.9) {mariadb};
  \node[font=\scriptsize, text width=6.6cm, align=center] at (3.7,1.4)
    {un conteneur d'application par service, créé à partir d'une image (Kolla)};
  \draw[draw=insagris, dashed, rounded corners=3pt, fill=insagris!5] (8.4,0.6) rectangle (15.8,4.7);
  \node[anchor=north west, font=\footnotesize\bfseries] at (8.55,4.6) {Hôte : LXC (ou directement sur l'hôte)};
  \node[lx] at (9.9,2.8) {\textbf{galera}\\ MariaDB};
  \node[lx] at (12.1,2.8) {\textbf{keystone}\\ environnement Python};
  \node[lx] at (14.3,2.8) {\textbf{nova\_api}\\ environnement Python};
  \node[font=\scriptsize, text width=6.8cm, align=center] at (12.1,1.2)
    {conteneurs « machine » LXC ; services installés depuis les sources};
  \node[callout, minimum width=15.6cm] at (7.9,-0.6)
    {Même outil (Ansible), deux façons d'isoler les services : des conteneurs d'application ou des conteneurs « machine ».};
\end{tikzpicture}}
\caption{Kolla-Ansible (à gauche) et OpenStack-Ansible (à droite). Le premier exécute chaque service dans un conteneur
d'application ; le second installe les services dans des environnements virtuels Python, au sein de conteneurs LXC.}
\label{fig:dep-osa}
\end{figure}

Le déploiement distingue l'hôte de déploiement, qui exécute Ansible, et les hôtes cibles. On y décrit le cloud dans
\texttt{/etc/openstack\_deploy/} : \texttt{openstack\_user\_config.yml} affecte les machines aux rôles (au moins trois hôtes
d'infrastructure dans \texttt{shared\_infra\_hosts}), \texttt{user\_variables.yml} règle les options globales et
\texttt{user\_secrets.yml} contient les mots de passe, que \texttt{pw-token-gen.py} génère \parencite{osa-config}. Le choix
\texttt{install\_method} (sources ou paquets) se fait au premier déploiement et ne peut plus changer ensuite \parencite{osa-config}.
Trois \glspl{playbook} font le travail, dans cet ordre \parencite{osa-run} : préparer les hôtes et construire les conteneurs, installer l'infrastructure (memcached, dépôt des sources, Galera, RabbitMQ), puis les services d'OpenStack.

\par\noindent\begin{minipage}{\linewidth}
\begin{lstlisting}[style=terminal]
# openstack-ansible openstack.osa.setup_hosts
# openstack-ansible openstack.osa.setup_infrastructure
# openstack-ansible openstack.osa.setup_openstack
\end{lstlisting}
\end{minipage}\par

Un déploiement « \gls{aio} » (\emph{all-in-one}, AIO), qui réunit tout sur une machine, sert à l'essai ou au laboratoire ; il demande
au minimum 8~vCPU, 8~Go de mémoire et 50~Go de disque (le projet en recommande 16~Go et 80~Go) \parencite{osa-aio}. OSA automatise
aussi les mises à niveau entre versions majeures. Le projet déconseille son choix si l'on veut des conteneurs d'application à 100~\%, si l'on
préfère Puppet, ou si l'on veut installer tous les services depuis les paquets de la distribution, dont la couverture reste partielle
\parencite{osa-about}.

% ---------------------------------------------------------------------
\subsection{Installation manuelle : ce que les outils automatisent}\label{sec:dep-manuel}
% ---------------------------------------------------------------------

Le guide d'installation officiel décrit, commande par commande, un cloud minimal installé avec les paquets de la distribution
\parencite{install-guide,install-overview}. Aucun outil ne s'interpose : on tape les commandes sur chaque nœud. On commence par
l'environnement commun (horloge synchronisée, dépôt des paquets OpenStack, MariaDB, RabbitMQ, memcached)
\parencite{install-env,install-packages-ubuntu,install-sql-ubuntu,install-messaging-ubuntu}, puis on installe les services dans
l'ordre que leurs dépendances imposent : Keystone, Glance, Placement, Nova, Neutron, puis Horizon et, en option, Cinder
\parencite{install-overview}. Chaque service suit le même scénario en six étapes (figure~\ref{fig:dep-manuel}).

\begin{figure}[!ht]
\centering
\adjustbox{max width=\linewidth}{%
\begin{tikzpicture}[>=Stealth, font=\footnotesize,
    cmd/.style={font=\scriptsize\ttfamily, text width=2.55cm, align=center}]
  \node[pgris, minimum width=15.6cm] (env) at (8.0,3.9)
    {Environnement commun, une fois sur le contrôleur : \ NTP \ $\cdot$ \ paquets OpenStack \ $\cdot$ \ MariaDB \ $\cdot$ \ RabbitMQ \ $\cdot$ \ memcached};
  \node[anchor=east, font=\footnotesize\bfseries, text=insarouge] at (15.9,2.85) {Pour chaque service (Glance en exemple) :};
  \foreach \x/\t/\c [count=\i] in {1.35/{Base de\\ données}/{CREATE DATABASE glance;},
        4.05/{Enregistrement\\ dans Keystone}/{openstack endpoint create},
        6.75/{Paquets}/{apt install glance},
        9.45/{Configuration}/{glance-api.conf},
        12.15/{Schéma de\\ la base}/{glance-manage db\_sync},
        14.85/{Redémarrage}/{systemctl restart glance-api}}
    { \node[pfoc, minimum width=2.45cm, minimum height=1.0cm] (s\i) at (\x,1.7) {\t};
      \node[num] at (\x-1.0,2.3) {\i};
      \node[cmd] at (\x,0.55) {\c}; }
  \foreach \a/\b in {1/2,2/3,3/4,4/5,5/6} \draw[fluxr] (s\a.east) -- (s\b.west);
  \draw[fluxd] (s1.north |- env.south) -- (s1.north);
  \node[callout, minimum width=15.6cm] at (8.0,-0.55)
    {Six étapes à répéter pour chaque service, sur chaque nœud : c'est précisément ce que les outils de déploiement automatisent.};
\end{tikzpicture}}
\caption{Le scénario d'installation manuelle d'un service, illustré par Glance \parencite{glance-install}. L'étape~2 inscrit le service
dans le catalogue de Keystone : un utilisateur de service, le rôle \texttt{admin}, une entrée de service et trois points d'accès.}
\label{fig:dep-manuel}
\end{figure}

\par\noindent\begin{minipage}{\linewidth}
\begin{lstlisting}[style=terminal]
$ sudo rabbitmqctl add_user openstack RABBIT_PASS
$ sudo rabbitmqctl set_permissions openstack ".*" ".*" ".*"
$ openstack user create --domain default --password-prompt glance
$ openstack role add --project service --user glance admin
$ openstack service create --name glance --description "OpenStack Image" image
$ openstack endpoint create --region RegionOne \
    image public http://controller:9292
\end{lstlisting}
\end{minipage}\par

Les deux premières commandes préparent RabbitMQ une fois pour toutes \parencite{install-messaging-ubuntu} ; les quatre suivantes
enregistrent Glance dans le catalogue (chapitre~\ref{chap:keystone}), puis on répète la dernière pour les points d'accès interne et
d'administration \parencite{glance-install}. Un nœud de calcul reçoit ensuite \texttt{nova-compute} et l'agent de Neutron, et doit être
associé à sa cellule, depuis le contrôleur, par \texttt{nova-manage cell\_v2 discover\_hosts} \parencite{nova-install-compute}.

Cette méthode a une vraie valeur pédagogique : on voit où chaque mot de passe, chaque base et chaque point d'accès apparaît. Elle ne
tient pas dans la durée : des dizaines de fichiers à garder cohérents, aucune idempotence (relancer une commande qui a déjà réussi
échoue ou duplique), aucune description unique du cloud, et des mises à niveau à la main. Au-delà d'un laboratoire, on délègue ce
travail à un outil.

\begin{remarque}[Deux sens de « bare metal »]
Le support parle d'installation « bare metal ». Le terme recouvre deux idées. Installer les services \emph{directement sur le
système des machines}, sans conteneur, c'est l'installation manuelle (ou OpenStack-Ansible sans LXC). Fournir \emph{des machines
physiques aux utilisateurs} est le rôle d'Ironic (chapitre~\ref{chap:ironic}) ; Kayobe s'en sert en outre pour installer le cloud lui-même
sur des serveurs nus (section~\ref{sec:dep-kolla}).
\end{remarque}

% ---------------------------------------------------------------------
\subsection{OpenStack sur Kubernetes}\label{sec:dep-autres}
% ---------------------------------------------------------------------

Deux approches installent OpenStack \emph{sur} Kubernetes, pour des publics différents. À ne pas confondre avec Magnum
(chapitre~\ref{chap:magnum}), qui fait l'inverse : il fournit des clusters Kubernetes \emph{sur} OpenStack.

\textbf{\gls{openstackhelm}} est une collection de charts Helm qui déploient OpenStack et les services associés sur Kubernetes
\parencite{helm-index,helm-readme}. Sa documentation vise des versions récentes de Kubernetes (1.33 à 1.35) et des hôtes Ubuntu
\parencite{helm-readme}. Elle s'adresse à une équipe qui exploite déjà Kubernetes et sait le faire vivre.

\textbf{\gls{sunbeam}} (Canonical OpenStack) s'installe avec un \emph{snap} et s'appuie sur \gls{juju} et Kubernetes. Le tutoriel officiel
déploie un cloud à un nœud (4~cœurs, 16~Gio de mémoire, 100~Gio de disque, Ubuntu 24.04) en cinq commandes \parencite{sunbeam-tuto} :

\par\noindent\begin{minipage}{\linewidth}
\begin{lstlisting}[style=terminal]
$ sudo snap install openstack
$ sunbeam prepare-node-script --bootstrap | bash -x && newgrp snap_daemon
$ sunbeam cluster bootstrap --accept-defaults --role control,compute,storage
$ sunbeam configure --accept-defaults --openrc demo-openrc
$ sunbeam launch ubuntu --name test
\end{lstlisting}
\end{minipage}\par

Pour agrandir le cloud, \texttt{sunbeam cluster add <FQDN>} produit un jeton que la nouvelle machine donne à
\texttt{sunbeam cluster join --role <rôles> -}. Les rôles sont \texttt{control}, \texttt{compute}, \texttt{network} et
\texttt{storage}, et \texttt{sunbeam cluster resize} répartit les fonctions de contrôle après l'ajout de nœuds de contrôle ;
les nœuds ne se désignent que par leur nom complet (FQDN) \parencite{sunbeam-multinode}. Red Hat suit une logique voisine avec
RHOSO, dont les opérateurs tournent sur OpenShift \parencite{rhoso-deploy}.

% ---------------------------------------------------------------------
\subsection{Choisir une méthode}\label{sec:dep-choix}
% ---------------------------------------------------------------------

Tous ces outils déploient le même OpenStack : les API, les services et les concepts des chapitres précédents ne changent pas. Seules
varient l'installation et l'exploitation, et c'est là que se joue le choix (figure~\ref{fig:dep-choix}).

\begin{figure}[!ht]
\centering
\adjustbox{max width=\linewidth}{%
\begin{tikzpicture}[>=Stealth, font=\footnotesize]
  \node[pcli, text width=6.0cm] (q1) at (3.4,5.3) {Apprendre ou développer, sur une machine jetable ?};
  \node[pcli, text width=6.0cm] (q2) at (3.4,3.1) {La plateforme existante repose-t-elle déjà sur Kubernetes ?};
  \node[pcli, text width=6.0cm] (q3) at (3.4,0.8) {Préférez-vous des conteneurs d'application, un par service ?};
  \draw[fluxd] (q1.south) -- (q2.north) node[etiqf, midway, right=2pt] {non};
  \draw[fluxd] (q2.south) -- (q3.north) node[etiqf, midway, right=2pt] {non};
  \node[pdat, text width=6.2cm] (r1a) at (12.2,5.9) {\textbf{DevStack} : développer, tester};
  \node[pdat, text width=6.2cm] (r1b) at (12.2,4.7) {\textbf{Installation manuelle}, \textbf{Sunbeam} à un nœud : apprendre};
  \node[psvc, text width=6.2cm] (r2)  at (12.2,3.1) {\textbf{OpenStack-Helm}, \textbf{Sunbeam}, \textbf{RHOSO}};
  \node[psvc, text width=6.2cm] (r3a) at (12.2,1.5) {\textbf{Kolla-Ansible} (Kayobe si le matériel est à provisionner)};
  \node[psvc, text width=6.2cm] (r3b) at (12.2,0.1) {\textbf{OpenStack-Ansible} (sources, conteneurs LXC)};
  \draw[fluxr] (q1.east) -- ++(0.75,0) |- (r1a.west);
  \draw[fluxr] (q1.east) -- ++(0.75,0) |- (r1b.west);
  \node[etiqf] at (7.0,5.55) {oui};
  \draw[fluxr] (q2.east) -- (r2.west) node[etiqf, midway, above] {oui};
  \draw[fluxr] (q3.east) -- ++(0.75,0) |- (r3a.west) node[etiqf, pos=0.8, above] {oui};
  \draw[fluxr] (q3.east) -- ++(0.75,0) |- (r3b.west) node[etiqf, pos=0.8, above] {non};
  \node[pdat, minimum width=0.5cm, inner xsep=4pt] at (2.6,-1.1) {};
  \node[anchor=west] at (2.95,-1.1) {essai, apprentissage};
  \node[psvc, minimum width=0.5cm, inner xsep=4pt] at (8.0,-1.1) {};
  \node[anchor=west] at (8.35,-1.1) {production};
  \node[callout, minimum width=15.6cm] at (7.9,-2.2)
    {Le bon choix dépend des compétences de l'équipe (Ansible, Kubernetes) et du support attendu, pas d'un classement général.};
\end{tikzpicture}}
\caption{Arbre de décision pour choisir une méthode de déploiement. Ce n'est pas la seule lecture possible : un fournisseur
(Canonical, Red Hat) peut imposer son offre, un client déjà équipé d'Ansible ira vers Kolla-Ansible ou OpenStack-Ansible.}
\label{fig:dep-choix}
\end{figure}

Quelques critères aident à départager Kolla-Ansible et OpenStack-Ansible, les deux outils de production les plus répandus :

\begin{itemize}
  \item \textbf{Isolation} : Kolla-Ansible isole chaque service dans une image figée ; OpenStack-Ansible installe depuis les sources
        dans des environnements virtuels, au sein de conteneurs LXC \parencite{osa-about}.
  \item \textbf{Configuration} : Kolla-Ansible est directif (quelques variables, puis des fichiers fusionnés) ; OpenStack-Ansible
        expose toute la configurabilité de ses rôles \parencite{kolla-ansible-philosophy,osa-about}.
  \item \textbf{Images} : avec Kolla-Ansible, il faut construire et héberger ses propres images pour la production ; avec
        OpenStack-Ansible, on gère un dépôt de sources et des environnements virtuels.
  \item \textbf{Mises à niveau} : toutes deux les automatisent (\texttt{kolla-ansible upgrade}, playbooks d'OpenStack-Ansible) ;
        aucune ne dispense de lire les notes de version.
\end{itemize}

% ---------------------------------------------------------------------
\subsection{Exploiter un cloud déployé}\label{sec:dep-exploitation}
% ---------------------------------------------------------------------

Déployer n'est que le début. Ce qui distingue un cloud fiable, c'est ce qu'on met en place autour : vérification, haute
disponibilité, sécurité, sauvegardes, supervision, mises à niveau.

\textbf{Vérifier.} Avant d'ouvrir le cloud aux utilisateurs, on contrôle que chaque service est enregistré et vivant. Avec
Kolla-Ansible, \texttt{post-deploy} écrit les accès administrateur dans \texttt{/etc/kolla/clouds.yaml}, et le script \texttt{init-runonce}
crée des réseaux, une image et d'autres ressources d'exemple (démonstration seulement) \parencite{kolla-ansible-quickstart} ; avec OpenStack-Ansible,
on utilise le conteneur utilitaire \parencite{osa-verify}.

\par\noindent\begin{minipage}{\linewidth}
\begin{lstlisting}[style=terminal]
$ openstack endpoint list                # chaque service a ses points d'accès
$ openstack compute service list         # conductor, scheduler, compute : up
$ openstack network agent list           # agents Neutron : Alive
$ openstack volume service list          # cinder-scheduler, cinder-volume : up
$ nova-status upgrade check              # prérequis d'une mise à niveau de Nova
\end{lstlisting}
\end{minipage}\par

\textbf{Haute disponibilité.} Elle repose sur deux principes \parencite{ha-concepts}. Les services sans état (API, ordonnanceurs)
tournent en actif/actif derrière HAProxy et une adresse virtuelle que \gls{keepalived} déplace en cas de panne. Les services à état
(MariaDB avec \gls{galera}, \gls{rabbitmq}) se répliquent en cluster d'un nombre impair de nœuds, dont le quorum est la majorité
(\(\lfloor n/2\rfloor+1\)) ; le \gls{plancontrole} compte donc au moins trois contrôleurs. La haute disponibilité ne protège pas
les machines des utilisateurs : si un hyperviseur tombe, Nova ne les relance pas seul (on évacue les instances avec
\texttt{openstack server evacuate}).

\begin{table}[!ht]
\centering
\renewcommand{\arraystretch}{1.25}
\small
\begin{tabular}{@{}>{\raggedright\arraybackslash}p{2.3cm}>{\raggedright\arraybackslash}p{13.1cm}@{}}
\toprule
\textbf{Sujet} & \textbf{Mise en œuvre avec Kolla-Ansible} \\
\midrule
Répartiteur & HAProxy et keepalived (groupe \texttt{loadbalancer}, qui reprend par défaut les nœuds \texttt{network}) ; variables \texttt{enable\_haproxy} et \texttt{enable\_keepalived} \parencite{kolla-ansible-haproxy}. \\
TLS & \texttt{kolla\_enable\_tls\_external} et \texttt{kolla\_enable\_tls\_internal} ; certificats dans \texttt{/etc/kolla/certificates} (\texttt{haproxy.pem}, \texttt{haproxy-internal.pem}) ; \texttt{kolla-ansible certificates} produit une autorité de test ; \texttt{enable\_letsencrypt} automatise le renouvellement \parencite{kolla-ansible-tls}. \\
Secrets & Mots de passe de \texttt{passwords.yml} (à sauvegarder avant tout changement). Rotation : vider les valeurs, \texttt{kolla-genpwd -p}, puis \texttt{kolla-ansible reconfigure} ; pour RabbitMQ, \texttt{stop} puis \texttt{deploy} \parencite{kolla-ansible-passwords}. \\
Sauvegarde & \texttt{kolla-ansible mariadb-backup} (complète ou incrémentale) \parencite{kolla-ansible-mariadb-backup} ; copie de \texttt{/etc/kolla} et de l'inventaire ; volumes et images selon les chapitres~\ref{chap:cinder} et~\ref{chap:glance}. \\
Journaux & Fichiers dans \texttt{/var/log/kolla/<service>} et \texttt{docker logs} ; \texttt{enable\_central\_logging} ajoute Fluentd et OpenSearch \parencite{kolla-ansible-troubleshooting,kolla-ansible-logging}. \\
Supervision & \texttt{enable\_prometheus} déploie Prometheus, que l'on peut relier à Grafana \parencite{kolla-ansible-prometheus}. \\
\bottomrule
\end{tabular}
\caption{L'exploitation d'un cloud Kolla-Ansible : ce que l'outil apporte.}
\label{tab:dep-exploit}
\end{table}

\textbf{Sécurité et sauvegardes.} Trois règles reviennent partout : garder le réseau de gestion interne (il porte les mots de passe
et les messages), chiffrer les API exposées (TLS) et protéger les secrets (tableau~\ref{tab:dep-exploit}). Une sauvegarde ne vaut que
si on sait restaurer : la base de données de chaque service, la configuration de l'outil et les mots de passe sont à sauvegarder
régulièrement et à restaurer lors d'un exercice.

\textbf{Mises à niveau.} Une version sort tous les six mois, mais on n'est pas obligé de les suivre toutes. Une version sur deux,
la première de l'année (2025.1, 2026.1\dots), est une version \gls{slurp} : on peut passer directement de l'une à la suivante, alors
qu'une version intermédiaire ne se met à niveau que vers la version qui la suit \parencite{openstack-slurp} (figure~\ref{fig:dep-slurp}).

\begin{figure}[!ht]
\centering
\adjustbox{max width=\linewidth}{%
\begin{tikzpicture}[>=Stealth, font=\footnotesize, adj/.style={-{Stealth[length=2mm]}, thick, insagris}]
  \foreach \x/\v/\n/\k [count=\i] in {1.5/2025.1/{}/psvc, 4.5/2025.2/{}/pgris, 7.5/2026.1/{Gazpacho}/psvc, 10.5/2026.2/{Hibiscus}/pgris, 13.5/2027.1/{Indri}/psvc}
    { \node[\k, minimum width=2.5cm, minimum height=1.1cm] (v\i) at (\x,1.0) {\textbf{\v}\\ {\scriptsize\mdseries \n}}; }
  \node[font=\scriptsize] at (1.5,0.2) {SLURP};
  \node[font=\scriptsize] at (7.5,0.2) {SLURP};
  \node[font=\scriptsize] at (13.5,0.2) {SLURP, en développement};
  \foreach \a/\b in {1/2,2/3,3/4,4/5} \draw[adj] (v\a.east) -- (v\b.west);
  \draw[fluxo] (v1.north) to[out=65,in=115] node[etiqf, midway, above] {saut SLURP $\to$ SLURP} (v3.north);
  \draw[fluxo] (v3.north) to[out=65,in=115] node[etiqf, midway, above] {saut SLURP $\to$ SLURP} (v5.north);
  \node[callout, minimum width=15.6cm] at (7.5,-0.8)
    {Depuis 2025.2, il faut passer par 2026.1 : une version non SLURP ne se met à niveau que vers la suivante.};
\end{tikzpicture}}
\caption{Chemins de mise à niveau (flèches grises : version suivante, toujours possible). Les noms des versions
2026.1, 2026.2 et 2027.1 sont ceux de la documentation \parencite{openstack-releases}.}
\label{fig:dep-slurp}
\end{figure}

On met toujours le plan de contrôle à niveau avant les nœuds de calcul, que l'on traite par lots : services anciens et récents
doivent coexister un temps. Nova impose un ordre (Placement d'abord, puis \texttt{nova-conductor} en premier et \texttt{nova-api} en
dernier), puis des migrations de données « en ligne » (\texttt{nova-manage db online\_data\_migrations}) \parencite{nova-upgrades}.
Kolla-Ansible enchaîne tout cela dans une seule commande \parencite{kolla-ansible-operating} :

\par\noindent\begin{minipage}{\linewidth}
\begin{lstlisting}[style=terminal]
$ cp -r /etc/kolla /etc/kolla.sauvegarde        # configuration d'abord
$ pip install --upgrade \
    git+https://opendev.org/openstack/kolla-ansible@stable/2026.2
$ kolla-ansible install-deps
$ kolla-mergepwd --old passwords.yml.old --new passwords.yml.new \
    --final /etc/kolla/passwords.yml
$ kolla-ansible pull && kolla-ansible prechecks
$ kolla-ansible upgrade
\end{lstlisting}
\end{minipage}\par

Avant cela, on lit les notes de version de la série visée, on met à jour l'inventaire et \texttt{globals.yml} d'après les exemples
fournis, et l'on fusionne les mots de passe avec \texttt{kolla-mergepwd}. Pour un saut SLURP, RabbitMQ ne supporte pas un saut de deux
versions : il faut d'abord le mettre à niveau vers une version intermédiaire. La documentation ne décrit pas de retour arrière :
on se protège par une sauvegarde de la base et de la configuration, et par un essai préalable sur un environnement de test
\parencite{kolla-ansible-operating}.

% ---------------------------------------------------------------------
\subsection{Exercices corrigés}\label{sec:dep-exercices}
% ---------------------------------------------------------------------

\begin{exercice}{Choisir un outil et une topologie}\label{ex:dep-choix}
Une PME de vingt utilisateurs veut une plateforme interne de développement sur quatre serveurs physiques. Une personne de l'équipe
maîtrise Ansible, personne ne connaît Kubernetes. Proposer un outil, une répartition des rôles et les points à prévoir avant
l'ouverture.
\end{exercice}
\begin{correction}
Pas de réponse unique ; l'argumentation compte. Production interne, compétence Ansible, aucun Kubernetes : on écarte DevStack,
Packstack, l'installation manuelle et les approches fondées sur Kubernetes, et l'on retient \textbf{Kolla-Ansible} ou
\textbf{OpenStack-Ansible}. Trois serveurs forment les contrôleurs (quorum de MariaDB et de RabbitMQ) et portent aussi le réseau et
le calcul ; le quatrième est un nœud de calcul. Les machines des utilisateurs partagent alors les serveurs avec le plan de
contrôle : acceptable à cette échelle, à surveiller. À prévoir : un hôte de déploiement distinct, un registre d'images local, une
adresse virtuelle inutilisée sur le réseau de gestion, TLS, la sauvegarde de MariaDB et de la configuration, une mise à niveau
annuelle sur les versions SLURP.
\end{correction}

\begin{exercice}{Relire une configuration}\label{ex:dep-relire}
Cet extrait d'un déploiement Kolla-Ansible contient trois erreurs de conception ; lesquelles ? On sait que \texttt{node1} a l'adresse
\texttt{10.1.0.11} sur \texttt{eth0}.

\begin{lstlisting}
[control]
node1
node2
[network]
node1
# globals.yml
kolla_internal_vip_address: "10.1.0.11"
network_interface: "eth0"
neutron_external_interface: "eth0"
\end{lstlisting}
\end{exercice}
\begin{correction}
(1) \textbf{Deux contrôleurs} n'apportent aucune haute disponibilité : la panne de l'un fait perdre le quorum ; il en faut trois.
(2) L'\textbf{adresse virtuelle est celle de \texttt{node1}} : elle doit être inutilisée sur le réseau de gestion, car keepalived la
porte et la déplace. (3) L'\textbf{interface externe est celle de gestion} : \texttt{neutron\_external\_interface} doit être une
interface dédiée, active et sans adresse IP, qui porte \texttt{br-ex}. Au passage, le groupe \texttt{network} se réduit à
\texttt{node1} : HAProxy et keepalived, qui suivent ce groupe, n'ont alors aucun nœud de secours.
\end{correction}

\begin{exercice}{Ajouter un service à un cloud en production}\label{ex:dep-service}
Le cloud, déployé avec Kolla-Ansible, doit fournir des répartiteurs de charge (chapitre~\ref{chap:octavia}). Quelles modifications et
quelles commandes faut-il ? Que reste-t-il à faire hors de l'outil ?
\end{exercice}
\begin{correction}
\begin{lstlisting}[style=terminal]
$ vi /etc/kolla/globals.yml                # enable_octavia: true
$ kolla-ansible octavia-certificates
$ kolla-ansible prechecks
$ kolla-ansible deploy
\end{lstlisting}
\texttt{deploy} est rejouable : il applique l'état décrit par les variables, sans repartir de zéro. Les certificats propres à
Octavia sont indispensables. Hors de l'outil, il reste à fournir ce que le chapitre~\ref{chap:octavia} décrit : l'image
d'amphore, le réseau et le groupe de sécurité de gestion, la \emph{flavor} des amphores.
\end{correction}

\begin{exercice}{Planifier une mise à niveau}\label{ex:dep-niveau}
Un cloud tourne en 2025.2 et l'équipe vise 2026.2. Peut-elle y passer directement ? Décrire la première mise à niveau avec
Kolla-Ansible.
\end{exercice}
\begin{correction}
Non : 2025.2 n'est pas une version SLURP, et ne se met à niveau que vers la suivante, 2026.1 ; de là, 2026.2 est la version
suivante. Il faut donc \textbf{deux mises à niveau successives}, chacune suivie d'une validation. Pour la première : lire les notes
de version de 2026.1 ; sauvegarder \texttt{/etc/kolla} et la base ; installer la branche \texttt{stable/2026.1} de
\texttt{kolla-ansible} ; \texttt{kolla-ansible install-deps} ; mettre à jour l'inventaire et \texttt{globals.yml} d'après les
exemples et fusionner les mots de passe avec \texttt{kolla-mergepwd} ; \texttt{kolla-ansible pull} puis \texttt{prechecks} ; enfin
\texttt{kolla-ansible upgrade}, qui recrée les conteneurs et migre les schémas. Aucun retour arrière n'est documenté : la
sauvegarde et l'essai sur un environnement de test sont la seule protection.
\end{correction}

% ---------------------------------------------------------------------
\subsection{Récapitulatif du chapitre}\label{sec:dep-recap}
% ---------------------------------------------------------------------

\begin{aretenir}
\begin{itemize}
  \item Un déploiement assemble des services, une infrastructure partagée (base, file de messages, cache, répartiteur) et des
        machines en rôles ; l'outil de déploiement automatise toute la pile.
  \item Un nombre impair de contrôleurs (trois au moins) pour le quorum ; séparer au moins les réseaux de gestion, de tunnel et
        externe, et garder le réseau de gestion interne.
  \item DevStack (sources, machine jetable) : développement et apprentissage. Kolla-Ansible (images de conteneurs, à construire
        soi-même en production) et OpenStack-Ansible (sources, conteneurs LXC) : production.
  \item TripleO est retiré ; Kolla-Ansible ne déploie plus Swift et ne déploie pas Ceph.
  \item L'exploitation compte autant que l'installation : haute disponibilité, TLS, secrets, sauvegardes, supervision, mises à
        niveau (SLURP : une version sur deux, avec saut direct).
\end{itemize}
\end{aretenir}

\begin{center}
\renewcommand{\arraystretch}{1.3}
\begin{tabular}{@{}>{\raggedright\arraybackslash}p{4.0cm}>{\raggedright\arraybackslash}p{9.8cm}r@{}}
\toprule
\textbf{Tâche} & \textbf{Commande} & \textbf{§} \\
\midrule
Préparer hôtes et secrets & \texttt{kolla-ansible bootstrap-servers}, \texttt{kolla-genpwd} & \ref{sec:dep-kolla} \\
Vérifier puis déployer & \texttt{kolla-ansible prechecks}, \texttt{pull}, \texttt{deploy} & \ref{sec:dep-kolla} \\
Appliquer une modification & \texttt{kolla-ansible reconfigure [-t SERVICE]} & \ref{sec:dep-kolla} \\
Sauvegarder MariaDB & \texttt{kolla-ansible mariadb-backup [--incremental]} & \ref{sec:dep-exploitation} \\
Mettre à niveau & \texttt{kolla-ansible upgrade} & \ref{sec:dep-exploitation} \\
Contrôler les services & \texttt{openstack compute service list}, \texttt{network agent list} & \ref{sec:dep-exploitation} \\
\bottomrule
\end{tabular}
\end{center}

\begin{savoirfaire}
\begin{itemize}
  \item Choisir un outil et une topologie à partir d'un contexte, en justifiant (exercice~\ref{ex:dep-choix}).
  \item Relire un inventaire et un \texttt{globals.yml} pour y repérer les erreurs de quorum, d'adresse virtuelle et d'interfaces
        (exercice~\ref{ex:dep-relire}).
  \item Ajouter un service à un cloud Kolla-Ansible et nommer ce qui reste à faire hors de l'outil (exercice~\ref{ex:dep-service}).
  \item Planifier un chemin de mise à niveau avec les versions SLURP (exercice~\ref{ex:dep-niveau}).
\end{itemize}
\end{savoirfaire}
.
%  Il ne contient ni préambule ni \begin{document} : les environnements
%  (definition, resultat, remarque, pointcle, aretenir, savoirfaire, exercice,
%  correction), les styles TikZ (pcli, pfoc, psvc, pdat, pgris, fluxo, fluxr,
%  fluxd, etiqf, callout, num, acteur, vie, msg, rep) et les couleurs (insarouge,
%  insagris, insaorange, insapeche) sont définis dans main.tex. Un chapitre = une \section.
%  Sources : support « Méthodes de déploiement » (chapitre 4) et documentation
%  officielle de DevStack, Kolla, Kolla-Ansible, OpenStack-Ansible, OpenStack-Helm,
%  du guide d'installation, des guides de haute disponibilité et de mise à niveau
%  (voir sources.bib). Versions décrites : celles de la documentation en vigueur
%  début octobre 2026 (version 2026.2 « Hibiscus » sortie le 30 septembre 2026).
%  Écarts avec le support : TripleO est retiré du projet (dépôt archivé en 2024) ;
%  Kolla-Ansible ne déploie plus Swift et ne déploie pas Ceph ; les images publiées par
%  Kolla ne sont pas destinées à la production ; la documentation de Kolla-Ansible ne
%  décrit pas de retour arrière après une mise à niveau.
%  Chapitre détaillé : 15 pages au plus.
% =====================================================================

\part{Mettre en œuvre}\label{part:miseenoeuvre}
\partdesc{Passer de la théorie à un cloud en fonctionnement : choix de l'outil de déploiement, topologie, installation, exploitation et mises à niveau.}

% =====================================================================
\section{Méthodes de déploiement d'OpenStack}\label{chap:deploiement}
% =====================================================================

Les chapitres précédents ont présenté chaque service du point de vue de l'utilisateur, puis de son architecture. Celui-ci se place
du point de vue de l'administrateur, qui doit installer la plateforme et la faire vivre. Il explique pourquoi ce travail est délicat,
décrit ce que tout déploiement assemble, présente les outils les plus répandus (DevStack, Kolla-Ansible, OpenStack-Ansible,
installation manuelle, approches fondées sur Kubernetes), aide à choisir, puis termine par l'exploitation d'un cloud installé :
haute disponibilité, sécurité, mises à niveau.

% ---------------------------------------------------------------------
\subsection{Pourquoi déployer OpenStack est délicat}\label{sec:dep-pourquoi}
% ---------------------------------------------------------------------

OpenStack n'est pas un logiciel unique mais une constellation de services (chapitres~\ref{chap:keystone} à~\ref{chap:trove}) :
chacun a son code, ses processus, son fichier de configuration et, le plus souvent, sa propre base de données. Tous reposent sur la
même infrastructure : une base SQL, une file de messages et un cache (\gls{memcached}) \parencite{install-env}. Un cloud réel
s'étend de plus sur plusieurs machines aux rôles différents, dont les versions doivent rester cohérentes entre elles. Tout installer,
configurer et mettre à jour à la main devient vite ingérable au-delà de quelques nœuds : c'est ce que les outils de déploiement
automatisent (figure~\ref{fig:dep-pile}).

\begin{definition}{Outil de déploiement}
Logiciel qui installe, configure et relie les services d'un cloud sur un ensemble de machines à partir d'une description
(inventaire, fichier de variables, modèles), de façon reproductible ; la plupart savent aussi le mettre à jour.
\end{definition}

\begin{figure}[!ht]
\centering
\adjustbox{max width=\linewidth}{%
\begin{tikzpicture}[>=Stealth, font=\footnotesize]
  \node[pcli, minimum width=9.6cm] (u) at (5.0,7.0) {Utilisateurs \ $\cdot$ \ ligne de commande \ $\cdot$ \ Horizon \ $\cdot$ \ API REST};
  \node[pfoc, minimum width=9.6cm, minimum height=1.0cm] (s) at (5.0,5.4)
    {Services OpenStack\\ {\scriptsize\mdseries Keystone, Nova, Neutron, Glance, Cinder, Octavia, Heat\dots}};
  \node[psvc, minimum width=9.6cm, minimum height=1.0cm] (i) at (5.0,3.8)
    {Services d'infrastructure partagés\\ {\scriptsize\mdseries MariaDB \ $\cdot$ \ RabbitMQ \ $\cdot$ \ memcached \ $\cdot$ \ HAProxy}};
  \node[pgris, minimum width=9.6cm, minimum height=1.0cm] (e) at (5.0,2.2)
    {Mode d'exécution\\ {\scriptsize\mdseries paquets \ $\cdot$ \ environnements Python \ $\cdot$ \ conteneurs \ $\cdot$ \ Kubernetes}};
  \node[pgris, minimum width=9.6cm, minimum height=1.0cm] (m) at (5.0,0.6)
    {Machines et réseaux\\ {\scriptsize\mdseries contrôleurs, calcul, stockage ; réseaux de gestion, de tunnel et externe}};
  \node[draw=insarouge, thick, rounded corners=4pt, fill=insarouge!8, minimum width=1.7cm, minimum height=6.0cm] (o) at (11.9,3.0) {};
  \node[rotate=90, font=\footnotesize\bfseries, text=insarouge] at (11.9,3.0) {Outil de déploiement};
  \foreach \b in {s,i,e,m} \draw[fluxr] (o.west |- \b.east) -- (\b.east);
  \node[callout, minimum width=13.2cm] at (6.3,-1.0)
    {L'outil de déploiement automatise toute la pile, pas seulement l'installation des services.};
\end{tikzpicture}}
\caption{Ce qu'un déploiement doit assembler. Les services (chapitres~\ref{chap:keystone} à~\ref{chap:trove}) reposent sur une
infrastructure partagée, installée sur des machines selon un mode d'exécution au choix ; l'outil de déploiement prend en charge
les quatre étages.}
\label{fig:dep-pile}
\end{figure}

Avant d'installer quoi que ce soit, il faut trancher les questions du tableau~\ref{tab:dep-decisions}. La version compte
autant que l'outil : OpenStack publie une version tous les six mois, la plus récente (2026.2 \emph{Hibiscus}) est sortie le
30~septembre 2026 et 2027.1 \emph{Indri} est en développement \parencite{openstack-releases}. Les outils sont eux-mêmes
versionnés par série, et l'on installe la branche de l'outil qui correspond à la version visée \parencite{kolla-ansible-operating}.

\begin{table}[!ht]
\centering
\renewcommand{\arraystretch}{1.3}
\begin{tabular}{@{}>{\raggedright\arraybackslash}p{3.0cm}>{\raggedright\arraybackslash}p{6.3cm}>{\raggedright\arraybackslash}p{6.1cm}@{}}
\toprule
\textbf{Décision} & \textbf{Question à trancher} & \textbf{Exemples de réponses} \\
\midrule
Topologie & Quels rôles, sur combien de nœuds ? & Un nœud pour l'essai ; trois contrôleurs pour la production. \\
Réseaux & Quelles interfaces pour la gestion, les tunnels, l'accès externe ? & Une interface par réseau, ou des VLAN sur un lien. \\
Stockage & Où vivent les volumes, les images, les objets ? & Disques locaux (LVM), cluster Ceph, NFS. \\
Mode d'installation & Paquets, sources, conteneurs ou Kubernetes ? & Voir le tableau~\ref{tab:dep-outils}. \\
Version & Quelle version, à quel rythme la mettre à niveau ? & Toutes les versions (six mois) ou une sur deux (\gls{slurp}). \\
Exploitation & Haute disponibilité, TLS, secrets, sauvegardes, supervision ? & Section~\ref{sec:dep-exploitation}. \\
\bottomrule
\end{tabular}
\caption{Les décisions à prendre avant de déployer.}
\label{tab:dep-decisions}
\end{table}

% ---------------------------------------------------------------------
\subsection{Anatomie d'un déploiement}\label{sec:dep-anatomie}
% ---------------------------------------------------------------------

Quelle que soit la méthode, on retrouve les mêmes rôles de nœuds, les mêmes services d'infrastructure et les mêmes réseaux
(figure~\ref{fig:dep-topo}).

\begin{itemize}
  \item \textbf{Hôte de déploiement} : la machine d'où l'on lance l'outil (\gls{ansible}, le plus souvent) ; en production, une machine
        distincte du cloud \parencite{osa-deploy-host,kolla-ansible-production}.
  \item \textbf{Contrôleurs} : API des services, ordonnanceurs, base de données et file de messages ; un nombre impair, pour
        établir un quorum \parencite{kolla-ansible-production}.
  \item \textbf{Réseau} : agents Neutron, routeurs ; avec Kolla-Ansible, HAProxy et keepalived s'y exécutent aussi
        \parencite{kolla-ansible-production}.
  \item \textbf{Calcul} : hyperviseur et \texttt{nova-compute}, là où vivent les machines virtuelles.
  \item \textbf{Stockage} : \texttt{cinder-volume} et ses disques ; les volumes peuvent aussi résider sur un cluster \gls{ceph} externe.
\end{itemize}

Le guide d'installation décrit une architecture minimale : un contrôleur, qui exécute l'identité, les images, Placement, les parties
de gestion de Nova et de Neutron, des agents réseau et le tableau de bord, et un nœud de calcul, avec en option des nœuds de stockage
bloc et objet. Il précise qu'elle n'est pas destinée à la production : les agents réseau y résident sur le contrôleur et le trafic
des tunnels emprunte le réseau de gestion \parencite{install-overview}.

\begin{figure}[!ht]
\centering
\adjustbox{max width=\linewidth}{%
\begin{tikzpicture}[>=Stealth, font=\footnotesize,
    nd/.style={draw=insaorange, fill=insapeche, thick, rounded corners=3pt, text width=3.3cm, align=center,
               inner sep=4pt, minimum height=2.3cm}]
  \node[pcli, minimum width=5.2cm] (dh) at (8.4,3.5) {Hôte de déploiement\\ {\scriptsize\mdseries Ansible, inventaire, variables}};
  \node[nd] (c) at (1.8,0)  {\textbf{Contrôleurs} {\scriptsize$\times 3$}\\[2pt] {\scriptsize API des services, ordonnanceurs, MariaDB, RabbitMQ, memcached}};
  \node[nd] (r) at (6.0,0)  {\textbf{Réseau}\\[2pt] {\scriptsize agents Neutron, routeurs, HAProxy et keepalived}};
  \node[nd] (k) at (10.2,0) {\textbf{Calcul} {\scriptsize$\times N$}\\[2pt] {\scriptsize hyperviseur, \texttt{nova-compute}, agent du réseau virtuel}};
  \node[nd] (s) at (14.4,0) {\textbf{Stockage}\\[2pt] {\scriptsize \texttt{cinder-volume}, disques locaux ou accès à Ceph}};
  \foreach \n in {c,r,k,s} \draw[fluxd] (dh.south) -- ++(0,-0.4) -| (\n.north);
  \node[etiqf] at (4.6,2.775) {Ansible, par SSH};
  % réseaux
  \foreach \y/\col/\tc/\nom in {-2.5/insarouge/insarouge/Gestion, -3.3/insaorange/insaorange/Tunnel,
      -4.1/insarouge!45/insarouge/Stockage, -4.9/insagris/insagris/Externe}
    { \draw[line width=2pt, \col] (0.2,\y) -- (16.1,\y);
      \node[anchor=east, font=\scriptsize\bfseries, text=\tc] at (-0.1,\y) {\nom}; }
  % branchements : décalage selon le réseau
  \foreach \cx/\y/\off/\col in {
      1.8/-2.5/-1.2/insarouge, 1.8/-4.9/1.2/insagris,
      6.0/-2.5/-1.2/insarouge, 6.0/-3.3/-0.4/insaorange, 6.0/-4.9/1.2/insagris,
      10.2/-2.5/-1.2/insarouge, 10.2/-3.3/-0.4/insaorange, 10.2/-4.1/0.4/insarouge!45,
      14.4/-2.5/-1.2/insarouge, 14.4/-4.1/0.4/insarouge!45}
    { \draw[thick, \col] ({\cx+\off},-1.15) -- ({\cx+\off},\y);
      \fill[\col] ({\cx+\off},\y) circle (2.6pt); }
  \node[callout, minimum width=15.2cm] at (8.2,-5.9)
    {Un même nœud peut tenir plusieurs rôles : l'inventaire de l'outil les décrit, le logiciel ne les fige pas.};
\end{tikzpicture}}
\caption{Topologie d'un déploiement multi-nœuds. L'hôte de déploiement configure les nœuds par SSH (pointillés) ; chaque rôle se
branche (points) sur les réseaux dont il a besoin.}
\label{fig:dep-topo}
\end{figure}

Les services d'infrastructure demandent une attention particulière. La base de données est le plus souvent un cluster MariaDB avec
\gls{galera} (multimaître), la file de messages un cluster \gls{rabbitmq}, et \gls{haproxy} répartit les requêtes des API vers les contrôleurs
\parencite{osa-service-arch}. Ces composants à état imposent un \gls{quorum} : deux instances n'apportent aucune haute
disponibilité (la panne de l'une paralyse l'autre), trois en tolèrent une en panne, cinq en tolèrent deux, et au-delà le gain
disparaît \parencite{kolla-ansible-production}. L'horloge de tous les nœuds doit en outre être synchronisée (\gls{ntp})
\parencite{install-env}.

Les réseaux se séparent par leur rôle (tableau~\ref{tab:dep-reseaux}) ; le trafic de gestion, qui porte les mots de passe et les
messages entre services, doit rester interne \parencite{kolla-ansible-production}.

\begin{table}[!ht]
\centering
\renewcommand{\arraystretch}{1.3}
\begin{tabular}{@{}>{\raggedright\arraybackslash}p{2.4cm}>{\raggedright\arraybackslash}p{7.6cm}>{\raggedright\arraybackslash}p{5.5cm}@{}}
\toprule
\textbf{Réseau} & \textbf{Rôle} & \textbf{Variable Kolla-Ansible} \\
\midrule
Gestion & Dialogue des services entre eux, avec la base et la file ; à garder interne. & \texttt{api\_\allowbreak interface} (défaut : \texttt{network\_\allowbreak interface}) \\
Tunnel & Trafic des réseaux de projets encapsulé (VXLAN, Geneve). & \texttt{tunnel\_\allowbreak interface} \\
Externe & Réseaux externes et adresses flottantes ; l'interface porte \texttt{br-ex} et n'a pas d'adresse IP. & \texttt{neutron\_\allowbreak external\_\allowbreak interface} \\
API publique & Points d'accès exposés aux utilisateurs derrière HAProxy. & \texttt{kolla\_\allowbreak external\_\allowbreak vip\_\allowbreak interface} \\
\bottomrule
\end{tabular}
\caption{Les réseaux d'un déploiement et leurs interfaces dans Kolla-Ansible \parencite{kolla-ansible-production}.}
\label{tab:dep-reseaux}
\end{table}

% ---------------------------------------------------------------------
\subsection{Panorama des méthodes}\label{sec:dep-panorama}
% ---------------------------------------------------------------------

Les outils se distinguent d'abord par la manière dont ils installent les services (figure~\ref{fig:dep-carte}), puis par leur public
(tableau~\ref{tab:dep-outils}). Aucun n'est « le meilleur » : chacun répond à un contexte.

\begin{figure}[!ht]
\centering
\adjustbox{max width=\linewidth}{%
\begin{tikzpicture}[>=Stealth, font=\footnotesize]
  \foreach \x/\t in {1.9/{Paquets de la\\ distribution}, 6.3/{Sources Git}, 10.7/{Images de\\ conteneurs}, 15.1/{Kubernetes}}
    \node[pgris, minimum width=3.9cm, minimum height=1.0cm] at (\x,4.6) {\t};
  \node[pdat, minimum width=3.7cm]  at (1.9,3.1)  {Installation manuelle};
  \node[pdat, minimum width=3.7cm]  at (1.9,2.2)  {Packstack (RDO)};
  \node[pdat, minimum width=3.7cm]  at (6.3,3.1)  {DevStack};
  \node[psvc, minimum width=3.7cm]  at (6.3,2.2)  {OpenStack-Ansible};
  \node[psvc, minimum width=3.7cm]  at (10.7,3.1) {Kolla-Ansible};
  \node[psvc, minimum width=3.7cm]  at (10.7,2.2) {Kayobe};
  \node[psvc, minimum width=3.7cm]  at (15.1,3.1) {OpenStack-Helm};
  \node[psvc, minimum width=3.7cm]  at (15.1,2.2) {Sunbeam (Canonical)};
  \node[psvc, minimum width=3.7cm]  at (15.1,1.3) {RHOSO (Red Hat)};
  \node[pdat, minimum width=0.5cm, inner xsep=4pt] at (3.2,0.0) {};
  \node[anchor=west] at (3.55,0.0) {essai, apprentissage, développement};
  \node[psvc, minimum width=0.5cm, inner xsep=4pt] at (9.4,0.0) {};
  \node[anchor=west] at (9.75,0.0) {production};
  \node[callout, minimum width=16.4cm] at (8.5,-1.2)
    {Ce qui change d'un outil à l'autre n'est pas la liste des services, mais la façon dont ils sont installés et mis à jour.};
\end{tikzpicture}}
\caption{Les outils de déploiement selon le mode d'installation des services. OpenStack-Ansible installe depuis les sources par
défaut (paquets possibles pour certains services) et peut isoler les services dans des conteneurs LXC \parencite{osa-config}.}
\label{fig:dep-carte}
\end{figure}

\begin{table}[!ht]
\centering
\renewcommand{\arraystretch}{1.25}
\small
\begin{tabular}{@{}>{\raggedright\arraybackslash}p{2.5cm}>{\raggedright\arraybackslash}p{5.5cm}>{\raggedright\arraybackslash}p{3.4cm}>{\raggedright\arraybackslash}p{3.8cm}@{}}
\toprule
\textbf{Outil} & \textbf{Principe} & \textbf{Usage visé} & \textbf{Ressources minimales documentées} \\
\midrule
Installation manuelle & Paquets de la distribution, guide d'installation officiel, service par service & Apprentissage, environnement isolé & Contrôleur : 4~Go ; calcul : 2~Go (preuve de concept) \parencite{install-env} \\
\gls{packstack} & Installeur du projet RDO (CentOS Stream) ; haute disponibilité et mises à niveau hors périmètre & Preuve de concept & 16~Go de mémoire \parencite{packstack} \\
DevStack & Scripts shell, installation depuis les sources Git & Développement, tests, démonstration & Machine dédiée (non précisé) \parencite{devstack-index} \\
OpenStack-Ansible & Ansible ; sources dans des environnements virtuels Python, conteneurs LXC & Production & Test (AIO) : 8~vCPU, 8~Go, 50~Go \parencite{osa-aio} \\
Kolla-Ansible & Ansible ; images de conteneurs Docker ou Podman & Production, évaluation & 2~interfaces, 8~Go, 40~Go \parencite{kolla-ansible-quickstart} \\
Kayobe & Kolla-Ansible, plus configuration des hôtes et provisionnement du matériel (Bifrost) & Production, depuis le serveur nu & --- \parencite{kayobe-index} \\
OpenStack-Helm & Charts Helm sur Kubernetes & Cloud sur Kubernetes & Kubernetes 1.33 à 1.35 \parencite{helm-readme} \\
Sunbeam & Snap, Juju, Kubernetes (Canonical OpenStack) \parencite{sunbeam-multinode} & Petits clouds, de un à plusieurs nœuds & Un nœud : 4~cœurs, 16~Gio, 100~Gio \parencite{sunbeam-tuto} \\
RHOSO & Opérateurs sur OpenShift (Red Hat) & Clients de Red Hat & --- \parencite{rhoso-deploy} \\
TripleO & Retiré du projet (dépôt archivé en 2024) & --- & --- \parencite{tripleo-docs} \\
\bottomrule
\end{tabular}
\caption{Les principaux outils de déploiement d'OpenStack.}
\label{tab:dep-outils}
\end{table}

TripleO (« OpenStack sur OpenStack » : un cloud d'administration, l'\emph{undercloud}, déployait le cloud de production,
l'\emph{overcloud}) figure encore dans de nombreux supports ; il n'est plus maintenu. Red Hat, qui en tirait son outil
\emph{director}, déploie aujourd'hui son offre par des opérateurs sur OpenShift (RHOSO) \parencite{tripleo-docs,rhoso-deploy}. La suite
détaille DevStack (section~\ref{sec:dep-devstack}), Kolla-Ansible (section~\ref{sec:dep-kolla}) et OpenStack-Ansible
(section~\ref{sec:dep-osa}), revient sur l'installation manuelle (section~\ref{sec:dep-manuel}) pour comprendre ce que les outils
automatisent, situe les approches fondées sur Kubernetes (section~\ref{sec:dep-autres}), puis aide à choisir
(section~\ref{sec:dep-choix}) et à exploiter (section~\ref{sec:dep-exploitation}).

% ---------------------------------------------------------------------
\subsection{DevStack : l'environnement tout-en-un}\label{sec:dep-devstack}
% ---------------------------------------------------------------------

\begin{definition}{DevStack}
Série de scripts qui déploient rapidement un cloud OpenStack complet sur une seule machine, à partir des dépôts Git de chaque projet
\parencite{devstack-index}.
\end{definition}

\gls{devstack} sert à construire vite des environnements de développement, à décrire des configurations qui fonctionnent, à faciliter les
contributions et à fournir l'environnement des tests d'intégration de chaque modification proposée \parencite{devstack-source}. Il
installe depuis les sources, sur la branche \texttt{master} par défaut ou sur une branche \texttt{stable/<série>} pour une version plus
ancienne : c'est ce qui en fait l'outil idéal pour essayer une modification de code (dans Nova, par exemple) avant de la proposer. Il
demande une machine dédiée et jetable. Ubuntu 24.04 est la plus testée ; le projet suit aussi les deux dernières versions LTS
d'Ubuntu, Rocky Linux 9 et openEuler. L'installation dure de 15 à 30 minutes \parencite{devstack-index}.

\begin{figure}[!ht]
\centering
\adjustbox{max width=\linewidth}{%
\begin{tikzpicture}[>=Stealth, font=\footnotesize]
  \node[pcli, minimum width=3.6cm] (lc) at (1.9,2.4) {\texttt{local.conf}\\ {\scriptsize\mdseries mots de passe, services}};
  \node[pgris, minimum width=3.6cm] (git) at (1.9,0.3) {Dépôts Git des projets\\ {\scriptsize\mdseries opendev.org}};
  \node[pfoc, minimum width=2.6cm, minimum height=1.3cm] (st) at (7.4,1.35) {\texttt{stack.sh}};
  \node[pgris, minimum width=5.6cm] (pk) at (13.0,2.8) {Paquets de la distribution\\ {\scriptsize\mdseries MySQL, RabbitMQ, Apache}};
  \node[psvc, minimum width=5.6cm] (sv) at (13.0,1.35) {Services installés dans \texttt{/opt/stack}\\ {\scriptsize\mdseries fichiers de configuration, bases}};
  \node[psvc, minimum width=5.6cm] (sd) at (13.0,-0.1) {Unités systemd \texttt{devstack@service}};
  \node[pdat, minimum width=9.4cm] (cl) at (9.4,-1.9) {Cloud à un nœud : API, Horizon, fichier \texttt{openrc}};
  \draw[fluxr] (lc.east) -| (st.north) node[num, pos=0.25] {1};
  \draw[fluxr] (git.east) -| (st.south) node[num, pos=0.25] {2};
  \foreach \t in {pk,sv,sd} \draw[fluxo] (st.east) -- (\t.west);
  \node[num] at (9.7,1.35) {3};
  \draw[fluxo] (sd.south) -- (sd.south |- cl.north);
  \node[callout, minimum width=15.6cm] at (8.5,-3.3)
    {DevStack installe depuis les sources, sur une machine dédiée : pour apprendre et développer, pas pour héberger un service réel.};
\end{tikzpicture}}
\caption{Ce que fait \texttt{stack.sh}. Il lit \texttt{local.conf} (1) et les dépôts des projets (2), puis installe les paquets de la
distribution, installe et configure les services dans \texttt{/opt/stack} et les démarre comme unités systemd (3).}
\label{fig:dep-devstack}
\end{figure}

Le script \texttt{stack.sh} ne doit pas être lancé en tant que \texttt{root} : il s'exécute sous un compte dédié, \texttt{stack}, qui
dispose de \texttt{sudo} sans mot de passe \parencite{devstack-index}. Ses services de base (Keystone, Nova, Placement, Glance,
Cinder, Neutron avec OVN et Horizon) sont actifs par défaut ; les autres s'ajoutent avec \texttt{enable\_service} (c'est le cas de Swift)
ou, pour un projet qui n'est pas dans le dépôt de DevStack, avec \texttt{enable\_plugin} \parencite{devstack-source,devstack-config,devstack-plugins}.
Chaque service tourne en unité systemd, dont on lit le journal avec \texttt{journalctl} \parencite{devstack-systemd}.

\par\noindent\begin{minipage}{\linewidth}
\begin{lstlisting}[style=terminal]
$ sudo useradd -s /bin/bash -d /opt/stack -m stack
$ sudo chmod +x /opt/stack
$ echo "stack ALL=(ALL) NOPASSWD: ALL" | sudo tee /etc/sudoers.d/stack
$ sudo -u stack -i
$ git clone https://opendev.org/openstack/devstack
$ cd devstack && nano local.conf
$ ./stack.sh
$ sudo journalctl -f --unit devstack@n-cpu.service
\end{lstlisting}
\end{minipage}\par

\par\noindent\begin{minipage}{\linewidth}
\begin{lstlisting}
[[local|localrc]]
ADMIN_PASSWORD=secret
DATABASE_PASSWORD=$ADMIN_PASSWORD
RABBIT_PASSWORD=$ADMIN_PASSWORD
SERVICE_PASSWORD=$ADMIN_PASSWORD
HOST_IP=192.168.1.10
# ajouter un projet absent des services par défaut (ici Heat)
enable_plugin heat https://opendev.org/openstack/heat
# retirer un service dont on n'a pas besoin
disable_service horizon
\end{lstlisting}
\end{minipage}\par

Le fichier \texttt{local.conf} est le seul à personnaliser : il suit un format de type INI dont la section \texttt{[[local|localrc]]}
fixe les mots de passe (des caractères alphanumériques seulement), l'adresse de l'hôte, les services actifs, et même les dépôts et
les branches à cloner (\texttt{NOVA\_REPO}, \texttt{NOVA\_BRANCH}) \parencite{devstack-config}. Une fois le script terminé, on charge les
identifiants avec \texttt{. openrc} puis on utilise la ligne de commande \parencite{devstack-source}.

Pour plusieurs machines, le contrôleur reçoit un \texttt{local.conf} complet et chaque nœud de calcul un fichier réduit, qui
désigne le contrôleur (\texttt{SERVICE\_HOST}, \texttt{MYSQL\_HOST}, \texttt{RABBIT\_HOST}) et ne garde que des services
comme \texttt{n-cpu}. Il faut ensuite associer chaque nouvel hôte à la cellule, depuis le contrôleur, avec
\texttt{tools/discover\_hosts.sh} (section~\ref{sec:nova-archi}) \parencite{devstack-multinode}.

DevStack n'est pas un installateur généraliste : la plupart des combinaisons d'options sont rarement testées, et il modifie en
profondeur le système hôte \parencite{devstack-overview,devstack-source}. Conçu pour un nœud, il n'offre pas de haute disponibilité.
\texttt{./unstack.sh} arrête les services (en laissant MySQL et RabbitMQ), \texttt{./clean.sh} efface les données.

% ---------------------------------------------------------------------
\subsection{Kolla-Ansible : OpenStack en conteneurs}\label{sec:dep-kolla}
% ---------------------------------------------------------------------

Le projet \gls{kolla} a pour mission de fournir des outils qui créent des conteneurs prêts pour la production et des outils qui déploient
des clouds OpenStack \parencite{kolla-source}. Il faut distinguer ses livrables :

\begin{itemize}
  \item \textbf{Kolla} construit les images de conteneurs de chaque service (\texttt{kolla-build}) ;
  \item \textbf{Kolla-Ansible} fournit les playbooks Ansible qui déploient ces images sur les nœuds ;
  \item \textbf{\gls{kayobe}} ajoute la configuration des hôtes et le provisionnement du matériel (avec Bifrost, qui s'appuie sur
        Ironic) à Kolla-Ansible \parencite{kayobe-index}.
\end{itemize}

\gls{kollaansible} fonctionne avec Docker ou Podman. Il prend en charge comme hôtes Debian 13, Rocky Linux 10 et Ubuntu 24.04, et les
images existent dans les mêmes trois distributions (\texttt{kolla\_base\_distro}) \parencite{kolla-ansible-support}. Par défaut il
tire les images publiées sur Quay.io (espace \texttt{openstack.kolla}). Ses auteurs avertissent que ces images, destinées aux
essais et aux démonstrations, ne subissent aucun examen de sécurité : en production, on construit et on maintient son propre jeu
d'images avec \texttt{kolla-build}, que l'on place dans un \gls{registre} local (recommandé dès que le déploiement compte plusieurs nœuds)
\parencite{kolla-source,kolla-image-building,kolla-ansible-multinode}.

\begin{figure}[!ht]
\centering
\adjustbox{max width=\linewidth}{%
\begin{tikzpicture}[>=Stealth, font=\footnotesize,
    ct/.style={draw=insaorange, fill=insapeche, rounded corners=2pt, font=\scriptsize\ttfamily, minimum height=0.55cm, inner xsep=3pt}]
  \node[pcli, minimum width=7.2cm] (dh) at (3.9,7.1)
    {Hôte de déploiement\\ {\scriptsize\mdseries \texttt{kolla-ansible}, inventaire, \texttt{globals.yml}, \texttt{passwords.yml}}};
  \node[pgris, minimum width=5.6cm] (reg) at (12.2,7.1) {Registre d'images\\ {\scriptsize\mdseries Quay.io ou registre local}};
  \foreach \y/\nom in {5.0/{Contrôleurs $\times 3$}, 3.6/{Réseau}, 2.2/{Calcul}, 0.8/{Stockage}}
    { \draw[draw=insagris, dashed, rounded corners=3pt, fill=insagris!5] (0,{\y-0.58}) rectangle (15.7,{\y+0.58});
      \node[anchor=west, font=\footnotesize\bfseries] at (0.15,\y) {\nom}; }
  \foreach \x/\t in {3.9/keystone, 5.9/nova\_api, 8.5/neutron\_server, 10.9/horizon, 12.7/mariadb, 14.5/rabbitmq}
    \node[ct] at (\x,5.0) {\t};
  \foreach \x/\t in {3.9/haproxy, 6.0/keepalived, 8.9/neutron\_l3\_agent, 12.1/neutron\_dhcp\_agent}
    \node[ct] at (\x,3.6) {\t};
  \foreach \x/\t in {4.2/nova\_compute, 6.8/nova\_libvirt, 10.0/openvswitch\_vswitchd}
    \node[ct] at (\x,2.2) {\t};
  \node[ct] at (4.2,0.8) {cinder\_volume};
  \node[pcli, minimum width=1.0cm, inner xsep=4pt] (vip) at (14.8,3.6) {VIP};
  \node[pgris, dashed, minimum width=3.8cm] (ceph) at (11.8,0.8) {Ceph (externe)};
  \draw[fluxd] (6.0,0.8) -- (ceph.west);
  \draw[fluxd] (dh.south) -- (3.9,5.65) node[etiqf, midway, right=3pt] {Ansible, par SSH};
  \draw[fluxo] (reg.south) -- (12.2,5.65) node[etiqf, midway, right=3pt] {tirer les images};
  \node[callout, minimum width=15.4cm] at (7.85,-0.75)
    {Chaque service tourne dans son conteneur, avec ses dépendances figées dans l'image : pas de conflit de versions entre services.};
\end{tikzpicture}}
\caption{Kolla-Ansible sur plusieurs nœuds. L'hôte de déploiement pilote les nœuds par Ansible (pointillés) ; ceux-ci tirent leurs
images d'un registre. Les groupes de l'inventaire fixent les conteneurs de chaque nœud. HAProxy et keepalived, qui portent
l'adresse virtuelle (VIP), s'exécutent par défaut sur les nœuds réseau. Kolla-Ansible ne déploie pas Ceph : il s'y connecte.}
\label{fig:dep-kolla}
\end{figure}

La configuration tient en trois éléments, placés sur l'hôte de déploiement, qu'installe \texttt{pip} dans un environnement virtuel
\parencite{kolla-ansible-quickstart} :
\begin{itemize}
  \item l'\textbf{\gls{inventaireansible}} Ansible, qui affecte les machines aux groupes \texttt{control}, \texttt{network}, \texttt{compute},
        \texttt{monitoring} et \texttt{storage} (le groupe \texttt{control} compte un nombre impair de nœuds) ; chaque service a lui-même
        son groupe, mais en modifier la répartition peut casser le déploiement \parencite{kolla-ansible-multinode,kolla-ansible-production} ;
  \item \texttt{/etc/kolla/globals.yml}, le fichier principal ;
  \item \texttt{/etc/kolla/passwords.yml}, dont les mots de passe vides se remplissent avec \texttt{kolla-genpwd}.
\end{itemize}

\par\noindent\begin{minipage}{\linewidth}
\begin{lstlisting}[style=terminal]
$ python3 -m venv /opt/kolla-venv && source /opt/kolla-venv/bin/activate
$ pip install git+https://opendev.org/openstack/kolla-ansible@stable/2026.1
$ sudo mkdir -p /etc/kolla && sudo chown $USER:$USER /etc/kolla
$ cp -r /opt/kolla-venv/share/kolla-ansible/etc_examples/kolla/* /etc/kolla
$ cp /opt/kolla-venv/share/kolla-ansible/ansible/inventory/multinode .
$ kolla-ansible install-deps
$ kolla-genpwd
\end{lstlisting}
\end{minipage}\par

\par\noindent\begin{minipage}{\linewidth}
\begin{lstlisting}
# /etc/kolla/globals.yml
kolla_base_distro: "ubuntu"
network_interface: "eth0"
neutron_external_interface: "eth1"
kolla_internal_vip_address: "10.1.0.250"
enable_cinder: true
enable_octavia: true
\end{lstlisting}
\end{minipage}\par

La valeur de \texttt{kolla\_internal\_vip\_address} doit être une adresse inutilisée du réseau de gestion : keepalived la porte et la
déplace en cas de panne ; l'interface externe doit être active, mais sans adresse IP \parencite{kolla-ansible-quickstart}. Kolla-Ansible
est volontairement très directif : quelques variables suffisent à un déploiement fonctionnel, et il ne multiplie pas les options.
Pour tout le reste, les fichiers déposés dans \texttt{/etc/kolla/config/} sont fusionnés avec la configuration par défaut du service
(par exemple, \texttt{nova.conf} avec \texttt{virt\_type = qemu} pour faire tourner le cloud dans des machines virtuelles sans
virtualisation imbriquée) \parencite{kolla-ansible-philosophy}. Les services de base (Keystone, Glance, Nova avec Placement, Neutron,
Heat, Horizon) sont actifs par défaut ; on ajoute les autres avec les variables \texttt{enable\_*} (tableau~\ref{tab:dep-kolla-services})
\parencite{kolla-ansible-source}.

\begin{table}[!ht]
\centering
\renewcommand{\arraystretch}{1.3}
\begin{tabular}{@{}>{\raggedright\arraybackslash}p{2.6cm}>{\raggedright\arraybackslash}p{3.4cm}>{\raggedright\arraybackslash}p{9.5cm}@{}}
\toprule
\textbf{Service} & \textbf{Variable} & \textbf{À savoir} \\
\midrule
Cinder (chap.~\ref{chap:cinder}) & \texttt{enable\_cinder} & Inactif par défaut ; disques LVM (\texttt{enable\_cinder\_\allowbreak backend\_\allowbreak lvm}) ou cluster Ceph externe. \\
Octavia (chap.~\ref{chap:octavia}) & \texttt{enable\_octavia} & Demande des certificats propres (\texttt{kolla-ansible octavia-certificates}) \parencite{kolla-ansible-source}. \\
Ironic (chap.~\ref{chap:ironic}) & \texttt{enable\_ironic} & S'appuie sur la même base de déploiement que le reste. \\
Magnum (chap.~\ref{chap:magnum}) & \texttt{enable\_magnum} & Le pilote de cluster reste à choisir (section~\ref{sec:ma-pilotes}). \\
Trove (chap.~\ref{chap:trove}) & \texttt{enable\_trove} & Les images invitées restent à construire (section~\ref{sec:tr-datastores}). \\
Swift (chap.~\ref{chap:swift}) & --- & Plus déployé par Kolla-Ansible (retiré faute de mainteneurs) ; Glance peut utiliser un Swift externe \parencite{kolla-ansible-source}. \\
\bottomrule
\end{tabular}
\caption{Activer les autres services avec Kolla-Ansible.}
\label{tab:dep-kolla-services}
\end{table}

Le déploiement suit un enchaînement fixe (figure~\ref{fig:dep-kolla-flux} et tableau~\ref{tab:dep-kolla-cmd}). La commande
\texttt{deploy} peut être relancée autant de fois que nécessaire, mais si une tâche d'amorçage échoue, un nouvel appel ne la répare pas ;
le plus rapide est alors de détruire ce déploiement raté (\texttt{kolla-ansible destroy}) \parencite{kolla-ansible-troubleshooting}.

\begin{figure}[!ht]
\centering
\adjustbox{max width=\linewidth}{%
\begin{tikzpicture}[>=Stealth, font=\footnotesize]
  \foreach \x/\c/\d [count=\i] in {1.6/bootstrap-servers/{prépare les hôtes}, 4.8/prechecks/{vérifie les prérequis}, 8.0/pull/{tire les images},
                                   11.2/deploy/{démarre les conteneurs}, 14.4/post-deploy/{écrit \texttt{clouds.yaml}}}
    { \node[pfoc, minimum width=2.9cm, font=\scriptsize\bfseries\ttfamily] (s\i) at (\x,2.0) {\c};
      \node[num] at (\x-1.25,2.55) {\i};
      \node[etiqf] at (\x,1.3) {\d}; }
  \foreach \a/\b in {1/2,2/3,3/4,4/5} \draw[fluxr] (s\a.east) -- (s\b.west);
  \node[anchor=west, font=\footnotesize\bfseries, text=insagris] at (0,-0.1) {Ensuite :};
  \foreach \x/\c in {3.6/reconfigure, 7.0/upgrade, 10.4/mariadb-backup, 13.8/destroy}
    \node[psvc, minimum width=2.9cm, font=\scriptsize\bfseries\ttfamily] at (\x,-0.1) {\c};
  \node[callout, minimum width=15.6cm] at (8.0,-1.4)
    {Les mêmes commandes servent à installer, à reconfigurer et à mettre à niveau : l'inventaire et les variables restent la seule description du cloud.};
\end{tikzpicture}}
\caption{Les cinq étapes d'un déploiement avec Kolla-Ansible, puis les commandes d'exploitation.}
\label{fig:dep-kolla-flux}
\end{figure}

\begin{table}[!ht]
\centering
\renewcommand{\arraystretch}{1.25}
\begin{tabular}{@{}>{\raggedright\arraybackslash\ttfamily\small}p{3.3cm}>{\raggedright\arraybackslash}p{12.2cm}@{}}
\toprule
\textrm{\textbf{Commande}} & \textbf{Effet} \\
\midrule
bootstrap-servers & Prépare chaque hôte : paquets système, moteur de conteneurs, utilisateurs \parencite{kolla-ansible-quickstart}. \\
prechecks & Vérifie que les hôtes sont dans un état où l'on peut déployer \parencite{kolla-ansible-troubleshooting}. \\
pull & Tire les images sur les hôtes, sans toucher aux conteneurs. \\
deploy & Génère les configurations, amorce les services et démarre tous les conteneurs. \\
post-deploy & Écrit les accès administrateur dans \texttt{/etc/kolla/clouds.yaml} \parencite{kolla-ansible-quickstart}. \\
reconfigure & Applique une modification de \texttt{globals.yml} ou de \texttt{config/}, éventuellement pour un seul service (\texttt{-t mariadb}). \\
validate-config & Valide les fichiers de configuration générés (après un premier déploiement). \\
upgrade & Met à niveau d'une série à la suivante : nouvelles images, migrations de schéma \parencite{kolla-ansible-operating}. \\
mariadb-backup & Sauvegarde complète ou incrémentale (\texttt{--incremental}) de MariaDB \parencite{kolla-ansible-mariadb-backup}. \\
destroy & Détruit les conteneurs, les volumes et la configuration des hôtes. \\
\bottomrule
\end{tabular}
\caption{Les commandes de \texttt{kolla-ansible} \parencite{kolla-ansible-operating,kolla-ansible-source}.}
\label{tab:dep-kolla-cmd}
\end{table}

% ---------------------------------------------------------------------
\subsection{OpenStack-Ansible : sources et conteneurs LXC}\label{sec:dep-osa}
% ---------------------------------------------------------------------

\gls{openstackansible} (OSA) déploie OpenStack avec Ansible, lui aussi pour la production. Par défaut, tous les services sont installés
depuis les sources dans des \glspl{venv} Python, puis isolés dans des conteneurs \gls{lxc}. Ce sont des conteneurs « machine »,
proches de petites machines virtuelles, et non des conteneurs d'application comme ceux de Docker ; ils sont facultatifs, et un
déploiement directement sur les hôtes est possible \parencite{osa-about}. Les hôtes pris en charge sont Ubuntu 24.04 et 26.04, Debian 13,
CentOS Stream 10 et Rocky Linux 10 \parencite{osa-about}.

\begin{figure}[!ht]
\centering
\adjustbox{max width=\linewidth}{%
\begin{tikzpicture}[>=Stealth, font=\footnotesize,
    ct/.style={draw=insaorange, fill=insapeche, rounded corners=2pt, font=\scriptsize\ttfamily, minimum height=0.6cm, minimum width=2.1cm},
    lx/.style={draw=insarouge!80!black, fill=insarouge!12, rounded corners=3pt, text width=1.8cm, align=center, font=\scriptsize,
               minimum height=1.35cm, inner sep=2pt}]
  \node[pcli, minimum width=3.4cm] at (3.7,5.4) {Kolla-Ansible};
  \node[pcli, minimum width=3.4cm] at (12.0,5.4) {OpenStack-Ansible};
  \draw[draw=insagris, dashed, rounded corners=3pt, fill=insagris!5] (0,0.6) rectangle (7.4,4.7);
  \node[anchor=north west, font=\footnotesize\bfseries] at (0.15,4.6) {Hôte : Docker ou Podman};
  \node[ct] at (1.4,2.9) {keystone};
  \node[ct] at (3.7,2.9) {nova\_api};
  \node[ct] at (6.0,2.9) {mariadb};
  \node[font=\scriptsize, text width=6.6cm, align=center] at (3.7,1.4)
    {un conteneur d'application par service, créé à partir d'une image (Kolla)};
  \draw[draw=insagris, dashed, rounded corners=3pt, fill=insagris!5] (8.4,0.6) rectangle (15.8,4.7);
  \node[anchor=north west, font=\footnotesize\bfseries] at (8.55,4.6) {Hôte : LXC (ou directement sur l'hôte)};
  \node[lx] at (9.9,2.8) {\textbf{galera}\\ MariaDB};
  \node[lx] at (12.1,2.8) {\textbf{keystone}\\ environnement Python};
  \node[lx] at (14.3,2.8) {\textbf{nova\_api}\\ environnement Python};
  \node[font=\scriptsize, text width=6.8cm, align=center] at (12.1,1.2)
    {conteneurs « machine » LXC ; services installés depuis les sources};
  \node[callout, minimum width=15.6cm] at (7.9,-0.6)
    {Même outil (Ansible), deux façons d'isoler les services : des conteneurs d'application ou des conteneurs « machine ».};
\end{tikzpicture}}
\caption{Kolla-Ansible (à gauche) et OpenStack-Ansible (à droite). Le premier exécute chaque service dans un conteneur
d'application ; le second installe les services dans des environnements virtuels Python, au sein de conteneurs LXC.}
\label{fig:dep-osa}
\end{figure}

Le déploiement distingue l'hôte de déploiement, qui exécute Ansible, et les hôtes cibles. On y décrit le cloud dans
\texttt{/etc/openstack\_deploy/} : \texttt{openstack\_user\_config.yml} affecte les machines aux rôles (au moins trois hôtes
d'infrastructure dans \texttt{shared\_infra\_hosts}), \texttt{user\_variables.yml} règle les options globales et
\texttt{user\_secrets.yml} contient les mots de passe, que \texttt{pw-token-gen.py} génère \parencite{osa-config}. Le choix
\texttt{install\_method} (sources ou paquets) se fait au premier déploiement et ne peut plus changer ensuite \parencite{osa-config}.
Trois \glspl{playbook} font le travail, dans cet ordre \parencite{osa-run} : préparer les hôtes et construire les conteneurs, installer l'infrastructure (memcached, dépôt des sources, Galera, RabbitMQ), puis les services d'OpenStack.

\par\noindent\begin{minipage}{\linewidth}
\begin{lstlisting}[style=terminal]
# openstack-ansible openstack.osa.setup_hosts
# openstack-ansible openstack.osa.setup_infrastructure
# openstack-ansible openstack.osa.setup_openstack
\end{lstlisting}
\end{minipage}\par

Un déploiement « \gls{aio} » (\emph{all-in-one}, AIO), qui réunit tout sur une machine, sert à l'essai ou au laboratoire ; il demande
au minimum 8~vCPU, 8~Go de mémoire et 50~Go de disque (le projet en recommande 16~Go et 80~Go) \parencite{osa-aio}. OSA automatise
aussi les mises à niveau entre versions majeures. Le projet déconseille son choix si l'on veut des conteneurs d'application à 100~\%, si l'on
préfère Puppet, ou si l'on veut installer tous les services depuis les paquets de la distribution, dont la couverture reste partielle
\parencite{osa-about}.

% ---------------------------------------------------------------------
\subsection{Installation manuelle : ce que les outils automatisent}\label{sec:dep-manuel}
% ---------------------------------------------------------------------

Le guide d'installation officiel décrit, commande par commande, un cloud minimal installé avec les paquets de la distribution
\parencite{install-guide,install-overview}. Aucun outil ne s'interpose : on tape les commandes sur chaque nœud. On commence par
l'environnement commun (horloge synchronisée, dépôt des paquets OpenStack, MariaDB, RabbitMQ, memcached)
\parencite{install-env,install-packages-ubuntu,install-sql-ubuntu,install-messaging-ubuntu}, puis on installe les services dans
l'ordre que leurs dépendances imposent : Keystone, Glance, Placement, Nova, Neutron, puis Horizon et, en option, Cinder
\parencite{install-overview}. Chaque service suit le même scénario en six étapes (figure~\ref{fig:dep-manuel}).

\begin{figure}[!ht]
\centering
\adjustbox{max width=\linewidth}{%
\begin{tikzpicture}[>=Stealth, font=\footnotesize,
    cmd/.style={font=\scriptsize\ttfamily, text width=2.55cm, align=center}]
  \node[pgris, minimum width=15.6cm] (env) at (8.0,3.9)
    {Environnement commun, une fois sur le contrôleur : \ NTP \ $\cdot$ \ paquets OpenStack \ $\cdot$ \ MariaDB \ $\cdot$ \ RabbitMQ \ $\cdot$ \ memcached};
  \node[anchor=east, font=\footnotesize\bfseries, text=insarouge] at (15.9,2.85) {Pour chaque service (Glance en exemple) :};
  \foreach \x/\t/\c [count=\i] in {1.35/{Base de\\ données}/{CREATE DATABASE glance;},
        4.05/{Enregistrement\\ dans Keystone}/{openstack endpoint create},
        6.75/{Paquets}/{apt install glance},
        9.45/{Configuration}/{glance-api.conf},
        12.15/{Schéma de\\ la base}/{glance-manage db\_sync},
        14.85/{Redémarrage}/{systemctl restart glance-api}}
    { \node[pfoc, minimum width=2.45cm, minimum height=1.0cm] (s\i) at (\x,1.7) {\t};
      \node[num] at (\x-1.0,2.3) {\i};
      \node[cmd] at (\x,0.55) {\c}; }
  \foreach \a/\b in {1/2,2/3,3/4,4/5,5/6} \draw[fluxr] (s\a.east) -- (s\b.west);
  \draw[fluxd] (s1.north |- env.south) -- (s1.north);
  \node[callout, minimum width=15.6cm] at (8.0,-0.55)
    {Six étapes à répéter pour chaque service, sur chaque nœud : c'est précisément ce que les outils de déploiement automatisent.};
\end{tikzpicture}}
\caption{Le scénario d'installation manuelle d'un service, illustré par Glance \parencite{glance-install}. L'étape~2 inscrit le service
dans le catalogue de Keystone : un utilisateur de service, le rôle \texttt{admin}, une entrée de service et trois points d'accès.}
\label{fig:dep-manuel}
\end{figure}

\par\noindent\begin{minipage}{\linewidth}
\begin{lstlisting}[style=terminal]
$ sudo rabbitmqctl add_user openstack RABBIT_PASS
$ sudo rabbitmqctl set_permissions openstack ".*" ".*" ".*"
$ openstack user create --domain default --password-prompt glance
$ openstack role add --project service --user glance admin
$ openstack service create --name glance --description "OpenStack Image" image
$ openstack endpoint create --region RegionOne \
    image public http://controller:9292
\end{lstlisting}
\end{minipage}\par

Les deux premières commandes préparent RabbitMQ une fois pour toutes \parencite{install-messaging-ubuntu} ; les quatre suivantes
enregistrent Glance dans le catalogue (chapitre~\ref{chap:keystone}), puis on répète la dernière pour les points d'accès interne et
d'administration \parencite{glance-install}. Un nœud de calcul reçoit ensuite \texttt{nova-compute} et l'agent de Neutron, et doit être
associé à sa cellule, depuis le contrôleur, par \texttt{nova-manage cell\_v2 discover\_hosts} \parencite{nova-install-compute}.

Cette méthode a une vraie valeur pédagogique : on voit où chaque mot de passe, chaque base et chaque point d'accès apparaît. Elle ne
tient pas dans la durée : des dizaines de fichiers à garder cohérents, aucune idempotence (relancer une commande qui a déjà réussi
échoue ou duplique), aucune description unique du cloud, et des mises à niveau à la main. Au-delà d'un laboratoire, on délègue ce
travail à un outil.

\begin{remarque}[Deux sens de « bare metal »]
Le support parle d'installation « bare metal ». Le terme recouvre deux idées. Installer les services \emph{directement sur le
système des machines}, sans conteneur, c'est l'installation manuelle (ou OpenStack-Ansible sans LXC). Fournir \emph{des machines
physiques aux utilisateurs} est le rôle d'Ironic (chapitre~\ref{chap:ironic}) ; Kayobe s'en sert en outre pour installer le cloud lui-même
sur des serveurs nus (section~\ref{sec:dep-kolla}).
\end{remarque}

% ---------------------------------------------------------------------
\subsection{OpenStack sur Kubernetes}\label{sec:dep-autres}
% ---------------------------------------------------------------------

Deux approches installent OpenStack \emph{sur} Kubernetes, pour des publics différents. À ne pas confondre avec Magnum
(chapitre~\ref{chap:magnum}), qui fait l'inverse : il fournit des clusters Kubernetes \emph{sur} OpenStack.

\textbf{\gls{openstackhelm}} est une collection de charts Helm qui déploient OpenStack et les services associés sur Kubernetes
\parencite{helm-index,helm-readme}. Sa documentation vise des versions récentes de Kubernetes (1.33 à 1.35) et des hôtes Ubuntu
\parencite{helm-readme}. Elle s'adresse à une équipe qui exploite déjà Kubernetes et sait le faire vivre.

\textbf{\gls{sunbeam}} (Canonical OpenStack) s'installe avec un \emph{snap} et s'appuie sur \gls{juju} et Kubernetes. Le tutoriel officiel
déploie un cloud à un nœud (4~cœurs, 16~Gio de mémoire, 100~Gio de disque, Ubuntu 24.04) en cinq commandes \parencite{sunbeam-tuto} :

\par\noindent\begin{minipage}{\linewidth}
\begin{lstlisting}[style=terminal]
$ sudo snap install openstack
$ sunbeam prepare-node-script --bootstrap | bash -x && newgrp snap_daemon
$ sunbeam cluster bootstrap --accept-defaults --role control,compute,storage
$ sunbeam configure --accept-defaults --openrc demo-openrc
$ sunbeam launch ubuntu --name test
\end{lstlisting}
\end{minipage}\par

Pour agrandir le cloud, \texttt{sunbeam cluster add <FQDN>} produit un jeton que la nouvelle machine donne à
\texttt{sunbeam cluster join --role <rôles> -}. Les rôles sont \texttt{control}, \texttt{compute}, \texttt{network} et
\texttt{storage}, et \texttt{sunbeam cluster resize} répartit les fonctions de contrôle après l'ajout de nœuds de contrôle ;
les nœuds ne se désignent que par leur nom complet (FQDN) \parencite{sunbeam-multinode}. Red Hat suit une logique voisine avec
RHOSO, dont les opérateurs tournent sur OpenShift \parencite{rhoso-deploy}.

% ---------------------------------------------------------------------
\subsection{Choisir une méthode}\label{sec:dep-choix}
% ---------------------------------------------------------------------

Tous ces outils déploient le même OpenStack : les API, les services et les concepts des chapitres précédents ne changent pas. Seules
varient l'installation et l'exploitation, et c'est là que se joue le choix (figure~\ref{fig:dep-choix}).

\begin{figure}[!ht]
\centering
\adjustbox{max width=\linewidth}{%
\begin{tikzpicture}[>=Stealth, font=\footnotesize]
  \node[pcli, text width=6.0cm] (q1) at (3.4,5.3) {Apprendre ou développer, sur une machine jetable ?};
  \node[pcli, text width=6.0cm] (q2) at (3.4,3.1) {La plateforme existante repose-t-elle déjà sur Kubernetes ?};
  \node[pcli, text width=6.0cm] (q3) at (3.4,0.8) {Préférez-vous des conteneurs d'application, un par service ?};
  \draw[fluxd] (q1.south) -- (q2.north) node[etiqf, midway, right=2pt] {non};
  \draw[fluxd] (q2.south) -- (q3.north) node[etiqf, midway, right=2pt] {non};
  \node[pdat, text width=6.2cm] (r1a) at (12.2,5.9) {\textbf{DevStack} : développer, tester};
  \node[pdat, text width=6.2cm] (r1b) at (12.2,4.7) {\textbf{Installation manuelle}, \textbf{Sunbeam} à un nœud : apprendre};
  \node[psvc, text width=6.2cm] (r2)  at (12.2,3.1) {\textbf{OpenStack-Helm}, \textbf{Sunbeam}, \textbf{RHOSO}};
  \node[psvc, text width=6.2cm] (r3a) at (12.2,1.5) {\textbf{Kolla-Ansible} (Kayobe si le matériel est à provisionner)};
  \node[psvc, text width=6.2cm] (r3b) at (12.2,0.1) {\textbf{OpenStack-Ansible} (sources, conteneurs LXC)};
  \draw[fluxr] (q1.east) -- ++(0.75,0) |- (r1a.west);
  \draw[fluxr] (q1.east) -- ++(0.75,0) |- (r1b.west);
  \node[etiqf] at (7.0,5.55) {oui};
  \draw[fluxr] (q2.east) -- (r2.west) node[etiqf, midway, above] {oui};
  \draw[fluxr] (q3.east) -- ++(0.75,0) |- (r3a.west) node[etiqf, pos=0.8, above] {oui};
  \draw[fluxr] (q3.east) -- ++(0.75,0) |- (r3b.west) node[etiqf, pos=0.8, above] {non};
  \node[pdat, minimum width=0.5cm, inner xsep=4pt] at (2.6,-1.1) {};
  \node[anchor=west] at (2.95,-1.1) {essai, apprentissage};
  \node[psvc, minimum width=0.5cm, inner xsep=4pt] at (8.0,-1.1) {};
  \node[anchor=west] at (8.35,-1.1) {production};
  \node[callout, minimum width=15.6cm] at (7.9,-2.2)
    {Le bon choix dépend des compétences de l'équipe (Ansible, Kubernetes) et du support attendu, pas d'un classement général.};
\end{tikzpicture}}
\caption{Arbre de décision pour choisir une méthode de déploiement. Ce n'est pas la seule lecture possible : un fournisseur
(Canonical, Red Hat) peut imposer son offre, un client déjà équipé d'Ansible ira vers Kolla-Ansible ou OpenStack-Ansible.}
\label{fig:dep-choix}
\end{figure}

Quelques critères aident à départager Kolla-Ansible et OpenStack-Ansible, les deux outils de production les plus répandus :

\begin{itemize}
  \item \textbf{Isolation} : Kolla-Ansible isole chaque service dans une image figée ; OpenStack-Ansible installe depuis les sources
        dans des environnements virtuels, au sein de conteneurs LXC \parencite{osa-about}.
  \item \textbf{Configuration} : Kolla-Ansible est directif (quelques variables, puis des fichiers fusionnés) ; OpenStack-Ansible
        expose toute la configurabilité de ses rôles \parencite{kolla-ansible-philosophy,osa-about}.
  \item \textbf{Images} : avec Kolla-Ansible, il faut construire et héberger ses propres images pour la production ; avec
        OpenStack-Ansible, on gère un dépôt de sources et des environnements virtuels.
  \item \textbf{Mises à niveau} : toutes deux les automatisent (\texttt{kolla-ansible upgrade}, playbooks d'OpenStack-Ansible) ;
        aucune ne dispense de lire les notes de version.
\end{itemize}

% ---------------------------------------------------------------------
\subsection{Exploiter un cloud déployé}\label{sec:dep-exploitation}
% ---------------------------------------------------------------------

Déployer n'est que le début. Ce qui distingue un cloud fiable, c'est ce qu'on met en place autour : vérification, haute
disponibilité, sécurité, sauvegardes, supervision, mises à niveau.

\textbf{Vérifier.} Avant d'ouvrir le cloud aux utilisateurs, on contrôle que chaque service est enregistré et vivant. Avec
Kolla-Ansible, \texttt{post-deploy} écrit les accès administrateur dans \texttt{/etc/kolla/clouds.yaml}, et le script \texttt{init-runonce}
crée des réseaux, une image et d'autres ressources d'exemple (démonstration seulement) \parencite{kolla-ansible-quickstart} ; avec OpenStack-Ansible,
on utilise le conteneur utilitaire \parencite{osa-verify}.

\par\noindent\begin{minipage}{\linewidth}
\begin{lstlisting}[style=terminal]
$ openstack endpoint list                # chaque service a ses points d'accès
$ openstack compute service list         # conductor, scheduler, compute : up
$ openstack network agent list           # agents Neutron : Alive
$ openstack volume service list          # cinder-scheduler, cinder-volume : up
$ nova-status upgrade check              # prérequis d'une mise à niveau de Nova
\end{lstlisting}
\end{minipage}\par

\textbf{Haute disponibilité.} Elle repose sur deux principes \parencite{ha-concepts}. Les services sans état (API, ordonnanceurs)
tournent en actif/actif derrière HAProxy et une adresse virtuelle que \gls{keepalived} déplace en cas de panne. Les services à état
(MariaDB avec \gls{galera}, \gls{rabbitmq}) se répliquent en cluster d'un nombre impair de nœuds, dont le quorum est la majorité
(\(\lfloor n/2\rfloor+1\)) ; le \gls{plancontrole} compte donc au moins trois contrôleurs. La haute disponibilité ne protège pas
les machines des utilisateurs : si un hyperviseur tombe, Nova ne les relance pas seul (on évacue les instances avec
\texttt{openstack server evacuate}).

\begin{table}[!ht]
\centering
\renewcommand{\arraystretch}{1.25}
\small
\begin{tabular}{@{}>{\raggedright\arraybackslash}p{2.3cm}>{\raggedright\arraybackslash}p{13.1cm}@{}}
\toprule
\textbf{Sujet} & \textbf{Mise en œuvre avec Kolla-Ansible} \\
\midrule
Répartiteur & HAProxy et keepalived (groupe \texttt{loadbalancer}, qui reprend par défaut les nœuds \texttt{network}) ; variables \texttt{enable\_haproxy} et \texttt{enable\_keepalived} \parencite{kolla-ansible-haproxy}. \\
TLS & \texttt{kolla\_enable\_tls\_external} et \texttt{kolla\_enable\_tls\_internal} ; certificats dans \texttt{/etc/kolla/certificates} (\texttt{haproxy.pem}, \texttt{haproxy-internal.pem}) ; \texttt{kolla-ansible certificates} produit une autorité de test ; \texttt{enable\_letsencrypt} automatise le renouvellement \parencite{kolla-ansible-tls}. \\
Secrets & Mots de passe de \texttt{passwords.yml} (à sauvegarder avant tout changement). Rotation : vider les valeurs, \texttt{kolla-genpwd -p}, puis \texttt{kolla-ansible reconfigure} ; pour RabbitMQ, \texttt{stop} puis \texttt{deploy} \parencite{kolla-ansible-passwords}. \\
Sauvegarde & \texttt{kolla-ansible mariadb-backup} (complète ou incrémentale) \parencite{kolla-ansible-mariadb-backup} ; copie de \texttt{/etc/kolla} et de l'inventaire ; volumes et images selon les chapitres~\ref{chap:cinder} et~\ref{chap:glance}. \\
Journaux & Fichiers dans \texttt{/var/log/kolla/<service>} et \texttt{docker logs} ; \texttt{enable\_central\_logging} ajoute Fluentd et OpenSearch \parencite{kolla-ansible-troubleshooting,kolla-ansible-logging}. \\
Supervision & \texttt{enable\_prometheus} déploie Prometheus, que l'on peut relier à Grafana \parencite{kolla-ansible-prometheus}. \\
\bottomrule
\end{tabular}
\caption{L'exploitation d'un cloud Kolla-Ansible : ce que l'outil apporte.}
\label{tab:dep-exploit}
\end{table}

\textbf{Sécurité et sauvegardes.} Trois règles reviennent partout : garder le réseau de gestion interne (il porte les mots de passe
et les messages), chiffrer les API exposées (TLS) et protéger les secrets (tableau~\ref{tab:dep-exploit}). Une sauvegarde ne vaut que
si on sait restaurer : la base de données de chaque service, la configuration de l'outil et les mots de passe sont à sauvegarder
régulièrement et à restaurer lors d'un exercice.

\textbf{Mises à niveau.} Une version sort tous les six mois, mais on n'est pas obligé de les suivre toutes. Une version sur deux,
la première de l'année (2025.1, 2026.1\dots), est une version \gls{slurp} : on peut passer directement de l'une à la suivante, alors
qu'une version intermédiaire ne se met à niveau que vers la version qui la suit \parencite{openstack-slurp} (figure~\ref{fig:dep-slurp}).

\begin{figure}[!ht]
\centering
\adjustbox{max width=\linewidth}{%
\begin{tikzpicture}[>=Stealth, font=\footnotesize, adj/.style={-{Stealth[length=2mm]}, thick, insagris}]
  \foreach \x/\v/\n/\k [count=\i] in {1.5/2025.1/{}/psvc, 4.5/2025.2/{}/pgris, 7.5/2026.1/{Gazpacho}/psvc, 10.5/2026.2/{Hibiscus}/pgris, 13.5/2027.1/{Indri}/psvc}
    { \node[\k, minimum width=2.5cm, minimum height=1.1cm] (v\i) at (\x,1.0) {\textbf{\v}\\ {\scriptsize\mdseries \n}}; }
  \node[font=\scriptsize] at (1.5,0.2) {SLURP};
  \node[font=\scriptsize] at (7.5,0.2) {SLURP};
  \node[font=\scriptsize] at (13.5,0.2) {SLURP, en développement};
  \foreach \a/\b in {1/2,2/3,3/4,4/5} \draw[adj] (v\a.east) -- (v\b.west);
  \draw[fluxo] (v1.north) to[out=65,in=115] node[etiqf, midway, above] {saut SLURP $\to$ SLURP} (v3.north);
  \draw[fluxo] (v3.north) to[out=65,in=115] node[etiqf, midway, above] {saut SLURP $\to$ SLURP} (v5.north);
  \node[callout, minimum width=15.6cm] at (7.5,-0.8)
    {Depuis 2025.2, il faut passer par 2026.1 : une version non SLURP ne se met à niveau que vers la suivante.};
\end{tikzpicture}}
\caption{Chemins de mise à niveau (flèches grises : version suivante, toujours possible). Les noms des versions
2026.1, 2026.2 et 2027.1 sont ceux de la documentation \parencite{openstack-releases}.}
\label{fig:dep-slurp}
\end{figure}

On met toujours le plan de contrôle à niveau avant les nœuds de calcul, que l'on traite par lots : services anciens et récents
doivent coexister un temps. Nova impose un ordre (Placement d'abord, puis \texttt{nova-conductor} en premier et \texttt{nova-api} en
dernier), puis des migrations de données « en ligne » (\texttt{nova-manage db online\_data\_migrations}) \parencite{nova-upgrades}.
Kolla-Ansible enchaîne tout cela dans une seule commande \parencite{kolla-ansible-operating} :

\par\noindent\begin{minipage}{\linewidth}
\begin{lstlisting}[style=terminal]
$ cp -r /etc/kolla /etc/kolla.sauvegarde        # configuration d'abord
$ pip install --upgrade \
    git+https://opendev.org/openstack/kolla-ansible@stable/2026.2
$ kolla-ansible install-deps
$ kolla-mergepwd --old passwords.yml.old --new passwords.yml.new \
    --final /etc/kolla/passwords.yml
$ kolla-ansible pull && kolla-ansible prechecks
$ kolla-ansible upgrade
\end{lstlisting}
\end{minipage}\par

Avant cela, on lit les notes de version de la série visée, on met à jour l'inventaire et \texttt{globals.yml} d'après les exemples
fournis, et l'on fusionne les mots de passe avec \texttt{kolla-mergepwd}. Pour un saut SLURP, RabbitMQ ne supporte pas un saut de deux
versions : il faut d'abord le mettre à niveau vers une version intermédiaire. La documentation ne décrit pas de retour arrière :
on se protège par une sauvegarde de la base et de la configuration, et par un essai préalable sur un environnement de test
\parencite{kolla-ansible-operating}.

% ---------------------------------------------------------------------
\subsection{Exercices corrigés}\label{sec:dep-exercices}
% ---------------------------------------------------------------------

\begin{exercice}{Choisir un outil et une topologie}\label{ex:dep-choix}
Une PME de vingt utilisateurs veut une plateforme interne de développement sur quatre serveurs physiques. Une personne de l'équipe
maîtrise Ansible, personne ne connaît Kubernetes. Proposer un outil, une répartition des rôles et les points à prévoir avant
l'ouverture.
\end{exercice}
\begin{correction}
Pas de réponse unique ; l'argumentation compte. Production interne, compétence Ansible, aucun Kubernetes : on écarte DevStack,
Packstack, l'installation manuelle et les approches fondées sur Kubernetes, et l'on retient \textbf{Kolla-Ansible} ou
\textbf{OpenStack-Ansible}. Trois serveurs forment les contrôleurs (quorum de MariaDB et de RabbitMQ) et portent aussi le réseau et
le calcul ; le quatrième est un nœud de calcul. Les machines des utilisateurs partagent alors les serveurs avec le plan de
contrôle : acceptable à cette échelle, à surveiller. À prévoir : un hôte de déploiement distinct, un registre d'images local, une
adresse virtuelle inutilisée sur le réseau de gestion, TLS, la sauvegarde de MariaDB et de la configuration, une mise à niveau
annuelle sur les versions SLURP.
\end{correction}

\begin{exercice}{Relire une configuration}\label{ex:dep-relire}
Cet extrait d'un déploiement Kolla-Ansible contient trois erreurs de conception ; lesquelles ? On sait que \texttt{node1} a l'adresse
\texttt{10.1.0.11} sur \texttt{eth0}.

\begin{lstlisting}
[control]
node1
node2
[network]
node1
# globals.yml
kolla_internal_vip_address: "10.1.0.11"
network_interface: "eth0"
neutron_external_interface: "eth0"
\end{lstlisting}
\end{exercice}
\begin{correction}
(1) \textbf{Deux contrôleurs} n'apportent aucune haute disponibilité : la panne de l'un fait perdre le quorum ; il en faut trois.
(2) L'\textbf{adresse virtuelle est celle de \texttt{node1}} : elle doit être inutilisée sur le réseau de gestion, car keepalived la
porte et la déplace. (3) L'\textbf{interface externe est celle de gestion} : \texttt{neutron\_external\_interface} doit être une
interface dédiée, active et sans adresse IP, qui porte \texttt{br-ex}. Au passage, le groupe \texttt{network} se réduit à
\texttt{node1} : HAProxy et keepalived, qui suivent ce groupe, n'ont alors aucun nœud de secours.
\end{correction}

\begin{exercice}{Ajouter un service à un cloud en production}\label{ex:dep-service}
Le cloud, déployé avec Kolla-Ansible, doit fournir des répartiteurs de charge (chapitre~\ref{chap:octavia}). Quelles modifications et
quelles commandes faut-il ? Que reste-t-il à faire hors de l'outil ?
\end{exercice}
\begin{correction}
\begin{lstlisting}[style=terminal]
$ vi /etc/kolla/globals.yml                # enable_octavia: true
$ kolla-ansible octavia-certificates
$ kolla-ansible prechecks
$ kolla-ansible deploy
\end{lstlisting}
\texttt{deploy} est rejouable : il applique l'état décrit par les variables, sans repartir de zéro. Les certificats propres à
Octavia sont indispensables. Hors de l'outil, il reste à fournir ce que le chapitre~\ref{chap:octavia} décrit : l'image
d'amphore, le réseau et le groupe de sécurité de gestion, la \emph{flavor} des amphores.
\end{correction}

\begin{exercice}{Planifier une mise à niveau}\label{ex:dep-niveau}
Un cloud tourne en 2025.2 et l'équipe vise 2026.2. Peut-elle y passer directement ? Décrire la première mise à niveau avec
Kolla-Ansible.
\end{exercice}
\begin{correction}
Non : 2025.2 n'est pas une version SLURP, et ne se met à niveau que vers la suivante, 2026.1 ; de là, 2026.2 est la version
suivante. Il faut donc \textbf{deux mises à niveau successives}, chacune suivie d'une validation. Pour la première : lire les notes
de version de 2026.1 ; sauvegarder \texttt{/etc/kolla} et la base ; installer la branche \texttt{stable/2026.1} de
\texttt{kolla-ansible} ; \texttt{kolla-ansible install-deps} ; mettre à jour l'inventaire et \texttt{globals.yml} d'après les
exemples et fusionner les mots de passe avec \texttt{kolla-mergepwd} ; \texttt{kolla-ansible pull} puis \texttt{prechecks} ; enfin
\texttt{kolla-ansible upgrade}, qui recrée les conteneurs et migre les schémas. Aucun retour arrière n'est documenté : la
sauvegarde et l'essai sur un environnement de test sont la seule protection.
\end{correction}

% ---------------------------------------------------------------------
\subsection{Récapitulatif du chapitre}\label{sec:dep-recap}
% ---------------------------------------------------------------------

\begin{aretenir}
\begin{itemize}
  \item Un déploiement assemble des services, une infrastructure partagée (base, file de messages, cache, répartiteur) et des
        machines en rôles ; l'outil de déploiement automatise toute la pile.
  \item Un nombre impair de contrôleurs (trois au moins) pour le quorum ; séparer au moins les réseaux de gestion, de tunnel et
        externe, et garder le réseau de gestion interne.
  \item DevStack (sources, machine jetable) : développement et apprentissage. Kolla-Ansible (images de conteneurs, à construire
        soi-même en production) et OpenStack-Ansible (sources, conteneurs LXC) : production.
  \item TripleO est retiré ; Kolla-Ansible ne déploie plus Swift et ne déploie pas Ceph.
  \item L'exploitation compte autant que l'installation : haute disponibilité, TLS, secrets, sauvegardes, supervision, mises à
        niveau (SLURP : une version sur deux, avec saut direct).
\end{itemize}
\end{aretenir}

\begin{center}
\renewcommand{\arraystretch}{1.3}
\begin{tabular}{@{}>{\raggedright\arraybackslash}p{4.0cm}>{\raggedright\arraybackslash}p{9.8cm}r@{}}
\toprule
\textbf{Tâche} & \textbf{Commande} & \textbf{§} \\
\midrule
Préparer hôtes et secrets & \texttt{kolla-ansible bootstrap-servers}, \texttt{kolla-genpwd} & \ref{sec:dep-kolla} \\
Vérifier puis déployer & \texttt{kolla-ansible prechecks}, \texttt{pull}, \texttt{deploy} & \ref{sec:dep-kolla} \\
Appliquer une modification & \texttt{kolla-ansible reconfigure [-t SERVICE]} & \ref{sec:dep-kolla} \\
Sauvegarder MariaDB & \texttt{kolla-ansible mariadb-backup [--incremental]} & \ref{sec:dep-exploitation} \\
Mettre à niveau & \texttt{kolla-ansible upgrade} & \ref{sec:dep-exploitation} \\
Contrôler les services & \texttt{openstack compute service list}, \texttt{network agent list} & \ref{sec:dep-exploitation} \\
\bottomrule
\end{tabular}
\end{center}

\begin{savoirfaire}
\begin{itemize}
  \item Choisir un outil et une topologie à partir d'un contexte, en justifiant (exercice~\ref{ex:dep-choix}).
  \item Relire un inventaire et un \texttt{globals.yml} pour y repérer les erreurs de quorum, d'adresse virtuelle et d'interfaces
        (exercice~\ref{ex:dep-relire}).
  \item Ajouter un service à un cloud Kolla-Ansible et nommer ce qui reste à faire hors de l'outil (exercice~\ref{ex:dep-service}).
  \item Planifier un chemin de mise à niveau avec les versions SLURP (exercice~\ref{ex:dep-niveau}).
\end{itemize}
\end{savoirfaire}
.
%  Il ne contient ni préambule ni \begin{document} : les environnements
%  (definition, resultat, remarque, pointcle, aretenir, savoirfaire, exercice,
%  correction), les styles TikZ (pcli, pfoc, psvc, pdat, pgris, fluxo, fluxr,
%  fluxd, etiqf, callout, num, acteur, vie, msg, rep) et les couleurs (insarouge,
%  insagris, insaorange, insapeche) sont définis dans main.tex. Un chapitre = une \section.
%  Sources : support « Méthodes de déploiement » (chapitre 4) et documentation
%  officielle de DevStack, Kolla, Kolla-Ansible, OpenStack-Ansible, OpenStack-Helm,
%  du guide d'installation, des guides de haute disponibilité et de mise à niveau
%  (voir sources.bib). Versions décrites : celles de la documentation en vigueur
%  début octobre 2026 (version 2026.2 « Hibiscus » sortie le 30 septembre 2026).
%  Écarts avec le support : TripleO est retiré du projet (dépôt archivé en 2024) ;
%  Kolla-Ansible ne déploie plus Swift et ne déploie pas Ceph ; les images publiées par
%  Kolla ne sont pas destinées à la production ; la documentation de Kolla-Ansible ne
%  décrit pas de retour arrière après une mise à niveau.
%  Chapitre détaillé : 15 pages au plus.
% =====================================================================

\part{Mettre en œuvre}\label{part:miseenoeuvre}
\partdesc{Passer de la théorie à un cloud en fonctionnement : choix de l'outil de déploiement, topologie, installation, exploitation et mises à niveau.}

% =====================================================================
\section{Méthodes de déploiement d'OpenStack}\label{chap:deploiement}
% =====================================================================

Les chapitres précédents ont présenté chaque service du point de vue de l'utilisateur, puis de son architecture. Celui-ci se place
du point de vue de l'administrateur, qui doit installer la plateforme et la faire vivre. Il explique pourquoi ce travail est délicat,
décrit ce que tout déploiement assemble, présente les outils les plus répandus (DevStack, Kolla-Ansible, OpenStack-Ansible,
installation manuelle, approches fondées sur Kubernetes), aide à choisir, puis termine par l'exploitation d'un cloud installé :
haute disponibilité, sécurité, mises à niveau.

% ---------------------------------------------------------------------
\subsection{Pourquoi déployer OpenStack est délicat}\label{sec:dep-pourquoi}
% ---------------------------------------------------------------------

OpenStack n'est pas un logiciel unique mais une constellation de services (chapitres~\ref{chap:keystone} à~\ref{chap:trove}) :
chacun a son code, ses processus, son fichier de configuration et, le plus souvent, sa propre base de données. Tous reposent sur la
même infrastructure : une base SQL, une file de messages et un cache (\gls{memcached}) \parencite{install-env}. Un cloud réel
s'étend de plus sur plusieurs machines aux rôles différents, dont les versions doivent rester cohérentes entre elles. Tout installer,
configurer et mettre à jour à la main devient vite ingérable au-delà de quelques nœuds : c'est ce que les outils de déploiement
automatisent (figure~\ref{fig:dep-pile}).

\begin{definition}{Outil de déploiement}
Logiciel qui installe, configure et relie les services d'un cloud sur un ensemble de machines à partir d'une description
(inventaire, fichier de variables, modèles), de façon reproductible ; la plupart savent aussi le mettre à jour.
\end{definition}

\begin{figure}[!ht]
\centering
\adjustbox{max width=\linewidth}{%
\begin{tikzpicture}[>=Stealth, font=\footnotesize]
  \node[pcli, minimum width=9.6cm] (u) at (5.0,7.0) {Utilisateurs \ $\cdot$ \ ligne de commande \ $\cdot$ \ Horizon \ $\cdot$ \ API REST};
  \node[pfoc, minimum width=9.6cm, minimum height=1.0cm] (s) at (5.0,5.4)
    {Services OpenStack\\ {\scriptsize\mdseries Keystone, Nova, Neutron, Glance, Cinder, Octavia, Heat\dots}};
  \node[psvc, minimum width=9.6cm, minimum height=1.0cm] (i) at (5.0,3.8)
    {Services d'infrastructure partagés\\ {\scriptsize\mdseries MariaDB \ $\cdot$ \ RabbitMQ \ $\cdot$ \ memcached \ $\cdot$ \ HAProxy}};
  \node[pgris, minimum width=9.6cm, minimum height=1.0cm] (e) at (5.0,2.2)
    {Mode d'exécution\\ {\scriptsize\mdseries paquets \ $\cdot$ \ environnements Python \ $\cdot$ \ conteneurs \ $\cdot$ \ Kubernetes}};
  \node[pgris, minimum width=9.6cm, minimum height=1.0cm] (m) at (5.0,0.6)
    {Machines et réseaux\\ {\scriptsize\mdseries contrôleurs, calcul, stockage ; réseaux de gestion, de tunnel et externe}};
  \node[draw=insarouge, thick, rounded corners=4pt, fill=insarouge!8, minimum width=1.7cm, minimum height=6.0cm] (o) at (11.9,3.0) {};
  \node[rotate=90, font=\footnotesize\bfseries, text=insarouge] at (11.9,3.0) {Outil de déploiement};
  \foreach \b in {s,i,e,m} \draw[fluxr] (o.west |- \b.east) -- (\b.east);
  \node[callout, minimum width=13.2cm] at (6.3,-1.0)
    {L'outil de déploiement automatise toute la pile, pas seulement l'installation des services.};
\end{tikzpicture}}
\caption{Ce qu'un déploiement doit assembler. Les services (chapitres~\ref{chap:keystone} à~\ref{chap:trove}) reposent sur une
infrastructure partagée, installée sur des machines selon un mode d'exécution au choix ; l'outil de déploiement prend en charge
les quatre étages.}
\label{fig:dep-pile}
\end{figure}

Avant d'installer quoi que ce soit, il faut trancher les questions du tableau~\ref{tab:dep-decisions}. La version compte
autant que l'outil : OpenStack publie une version tous les six mois, la plus récente (2026.2 \emph{Hibiscus}) est sortie le
30~septembre 2026 et 2027.1 \emph{Indri} est en développement \parencite{openstack-releases}. Les outils sont eux-mêmes
versionnés par série, et l'on installe la branche de l'outil qui correspond à la version visée \parencite{kolla-ansible-operating}.

\begin{table}[!ht]
\centering
\renewcommand{\arraystretch}{1.3}
\begin{tabular}{@{}>{\raggedright\arraybackslash}p{3.0cm}>{\raggedright\arraybackslash}p{6.3cm}>{\raggedright\arraybackslash}p{6.1cm}@{}}
\toprule
\textbf{Décision} & \textbf{Question à trancher} & \textbf{Exemples de réponses} \\
\midrule
Topologie & Quels rôles, sur combien de nœuds ? & Un nœud pour l'essai ; trois contrôleurs pour la production. \\
Réseaux & Quelles interfaces pour la gestion, les tunnels, l'accès externe ? & Une interface par réseau, ou des VLAN sur un lien. \\
Stockage & Où vivent les volumes, les images, les objets ? & Disques locaux (LVM), cluster Ceph, NFS. \\
Mode d'installation & Paquets, sources, conteneurs ou Kubernetes ? & Voir le tableau~\ref{tab:dep-outils}. \\
Version & Quelle version, à quel rythme la mettre à niveau ? & Toutes les versions (six mois) ou une sur deux (\gls{slurp}). \\
Exploitation & Haute disponibilité, TLS, secrets, sauvegardes, supervision ? & Section~\ref{sec:dep-exploitation}. \\
\bottomrule
\end{tabular}
\caption{Les décisions à prendre avant de déployer.}
\label{tab:dep-decisions}
\end{table}

% ---------------------------------------------------------------------
\subsection{Anatomie d'un déploiement}\label{sec:dep-anatomie}
% ---------------------------------------------------------------------

Quelle que soit la méthode, on retrouve les mêmes rôles de nœuds, les mêmes services d'infrastructure et les mêmes réseaux
(figure~\ref{fig:dep-topo}).

\begin{itemize}
  \item \textbf{Hôte de déploiement} : la machine d'où l'on lance l'outil (\gls{ansible}, le plus souvent) ; en production, une machine
        distincte du cloud \parencite{osa-deploy-host,kolla-ansible-production}.
  \item \textbf{Contrôleurs} : API des services, ordonnanceurs, base de données et file de messages ; un nombre impair, pour
        établir un quorum \parencite{kolla-ansible-production}.
  \item \textbf{Réseau} : agents Neutron, routeurs ; avec Kolla-Ansible, HAProxy et keepalived s'y exécutent aussi
        \parencite{kolla-ansible-production}.
  \item \textbf{Calcul} : hyperviseur et \texttt{nova-compute}, là où vivent les machines virtuelles.
  \item \textbf{Stockage} : \texttt{cinder-volume} et ses disques ; les volumes peuvent aussi résider sur un cluster \gls{ceph} externe.
\end{itemize}

Le guide d'installation décrit une architecture minimale : un contrôleur, qui exécute l'identité, les images, Placement, les parties
de gestion de Nova et de Neutron, des agents réseau et le tableau de bord, et un nœud de calcul, avec en option des nœuds de stockage
bloc et objet. Il précise qu'elle n'est pas destinée à la production : les agents réseau y résident sur le contrôleur et le trafic
des tunnels emprunte le réseau de gestion \parencite{install-overview}.

\begin{figure}[!ht]
\centering
\adjustbox{max width=\linewidth}{%
\begin{tikzpicture}[>=Stealth, font=\footnotesize,
    nd/.style={draw=insaorange, fill=insapeche, thick, rounded corners=3pt, text width=3.3cm, align=center,
               inner sep=4pt, minimum height=2.3cm}]
  \node[pcli, minimum width=5.2cm] (dh) at (8.4,3.5) {Hôte de déploiement\\ {\scriptsize\mdseries Ansible, inventaire, variables}};
  \node[nd] (c) at (1.8,0)  {\textbf{Contrôleurs} {\scriptsize$\times 3$}\\[2pt] {\scriptsize API des services, ordonnanceurs, MariaDB, RabbitMQ, memcached}};
  \node[nd] (r) at (6.0,0)  {\textbf{Réseau}\\[2pt] {\scriptsize agents Neutron, routeurs, HAProxy et keepalived}};
  \node[nd] (k) at (10.2,0) {\textbf{Calcul} {\scriptsize$\times N$}\\[2pt] {\scriptsize hyperviseur, \texttt{nova-compute}, agent du réseau virtuel}};
  \node[nd] (s) at (14.4,0) {\textbf{Stockage}\\[2pt] {\scriptsize \texttt{cinder-volume}, disques locaux ou accès à Ceph}};
  \foreach \n in {c,r,k,s} \draw[fluxd] (dh.south) -- ++(0,-0.4) -| (\n.north);
  \node[etiqf] at (4.6,2.775) {Ansible, par SSH};
  % réseaux
  \foreach \y/\col/\tc/\nom in {-2.5/insarouge/insarouge/Gestion, -3.3/insaorange/insaorange/Tunnel,
      -4.1/insarouge!45/insarouge/Stockage, -4.9/insagris/insagris/Externe}
    { \draw[line width=2pt, \col] (0.2,\y) -- (16.1,\y);
      \node[anchor=east, font=\scriptsize\bfseries, text=\tc] at (-0.1,\y) {\nom}; }
  % branchements : décalage selon le réseau
  \foreach \cx/\y/\off/\col in {
      1.8/-2.5/-1.2/insarouge, 1.8/-4.9/1.2/insagris,
      6.0/-2.5/-1.2/insarouge, 6.0/-3.3/-0.4/insaorange, 6.0/-4.9/1.2/insagris,
      10.2/-2.5/-1.2/insarouge, 10.2/-3.3/-0.4/insaorange, 10.2/-4.1/0.4/insarouge!45,
      14.4/-2.5/-1.2/insarouge, 14.4/-4.1/0.4/insarouge!45}
    { \draw[thick, \col] ({\cx+\off},-1.15) -- ({\cx+\off},\y);
      \fill[\col] ({\cx+\off},\y) circle (2.6pt); }
  \node[callout, minimum width=15.2cm] at (8.2,-5.9)
    {Un même nœud peut tenir plusieurs rôles : l'inventaire de l'outil les décrit, le logiciel ne les fige pas.};
\end{tikzpicture}}
\caption{Topologie d'un déploiement multi-nœuds. L'hôte de déploiement configure les nœuds par SSH (pointillés) ; chaque rôle se
branche (points) sur les réseaux dont il a besoin.}
\label{fig:dep-topo}
\end{figure}

Les services d'infrastructure demandent une attention particulière. La base de données est le plus souvent un cluster MariaDB avec
\gls{galera} (multimaître), la file de messages un cluster \gls{rabbitmq}, et \gls{haproxy} répartit les requêtes des API vers les contrôleurs
\parencite{osa-service-arch}. Ces composants à état imposent un \gls{quorum} : deux instances n'apportent aucune haute
disponibilité (la panne de l'une paralyse l'autre), trois en tolèrent une en panne, cinq en tolèrent deux, et au-delà le gain
disparaît \parencite{kolla-ansible-production}. L'horloge de tous les nœuds doit en outre être synchronisée (\gls{ntp})
\parencite{install-env}.

Les réseaux se séparent par leur rôle (tableau~\ref{tab:dep-reseaux}) ; le trafic de gestion, qui porte les mots de passe et les
messages entre services, doit rester interne \parencite{kolla-ansible-production}.

\begin{table}[!ht]
\centering
\renewcommand{\arraystretch}{1.3}
\begin{tabular}{@{}>{\raggedright\arraybackslash}p{2.4cm}>{\raggedright\arraybackslash}p{7.6cm}>{\raggedright\arraybackslash}p{5.5cm}@{}}
\toprule
\textbf{Réseau} & \textbf{Rôle} & \textbf{Variable Kolla-Ansible} \\
\midrule
Gestion & Dialogue des services entre eux, avec la base et la file ; à garder interne. & \texttt{api\_\allowbreak interface} (défaut : \texttt{network\_\allowbreak interface}) \\
Tunnel & Trafic des réseaux de projets encapsulé (VXLAN, Geneve). & \texttt{tunnel\_\allowbreak interface} \\
Externe & Réseaux externes et adresses flottantes ; l'interface porte \texttt{br-ex} et n'a pas d'adresse IP. & \texttt{neutron\_\allowbreak external\_\allowbreak interface} \\
API publique & Points d'accès exposés aux utilisateurs derrière HAProxy. & \texttt{kolla\_\allowbreak external\_\allowbreak vip\_\allowbreak interface} \\
\bottomrule
\end{tabular}
\caption{Les réseaux d'un déploiement et leurs interfaces dans Kolla-Ansible \parencite{kolla-ansible-production}.}
\label{tab:dep-reseaux}
\end{table}

% ---------------------------------------------------------------------
\subsection{Panorama des méthodes}\label{sec:dep-panorama}
% ---------------------------------------------------------------------

Les outils se distinguent d'abord par la manière dont ils installent les services (figure~\ref{fig:dep-carte}), puis par leur public
(tableau~\ref{tab:dep-outils}). Aucun n'est « le meilleur » : chacun répond à un contexte.

\begin{figure}[!ht]
\centering
\adjustbox{max width=\linewidth}{%
\begin{tikzpicture}[>=Stealth, font=\footnotesize]
  \foreach \x/\t in {1.9/{Paquets de la\\ distribution}, 6.3/{Sources Git}, 10.7/{Images de\\ conteneurs}, 15.1/{Kubernetes}}
    \node[pgris, minimum width=3.9cm, minimum height=1.0cm] at (\x,4.6) {\t};
  \node[pdat, minimum width=3.7cm]  at (1.9,3.1)  {Installation manuelle};
  \node[pdat, minimum width=3.7cm]  at (1.9,2.2)  {Packstack (RDO)};
  \node[pdat, minimum width=3.7cm]  at (6.3,3.1)  {DevStack};
  \node[psvc, minimum width=3.7cm]  at (6.3,2.2)  {OpenStack-Ansible};
  \node[psvc, minimum width=3.7cm]  at (10.7,3.1) {Kolla-Ansible};
  \node[psvc, minimum width=3.7cm]  at (10.7,2.2) {Kayobe};
  \node[psvc, minimum width=3.7cm]  at (15.1,3.1) {OpenStack-Helm};
  \node[psvc, minimum width=3.7cm]  at (15.1,2.2) {Sunbeam (Canonical)};
  \node[psvc, minimum width=3.7cm]  at (15.1,1.3) {RHOSO (Red Hat)};
  \node[pdat, minimum width=0.5cm, inner xsep=4pt] at (3.2,0.0) {};
  \node[anchor=west] at (3.55,0.0) {essai, apprentissage, développement};
  \node[psvc, minimum width=0.5cm, inner xsep=4pt] at (9.4,0.0) {};
  \node[anchor=west] at (9.75,0.0) {production};
  \node[callout, minimum width=16.4cm] at (8.5,-1.2)
    {Ce qui change d'un outil à l'autre n'est pas la liste des services, mais la façon dont ils sont installés et mis à jour.};
\end{tikzpicture}}
\caption{Les outils de déploiement selon le mode d'installation des services. OpenStack-Ansible installe depuis les sources par
défaut (paquets possibles pour certains services) et peut isoler les services dans des conteneurs LXC \parencite{osa-config}.}
\label{fig:dep-carte}
\end{figure}

\begin{table}[!ht]
\centering
\renewcommand{\arraystretch}{1.25}
\small
\begin{tabular}{@{}>{\raggedright\arraybackslash}p{2.5cm}>{\raggedright\arraybackslash}p{5.5cm}>{\raggedright\arraybackslash}p{3.4cm}>{\raggedright\arraybackslash}p{3.8cm}@{}}
\toprule
\textbf{Outil} & \textbf{Principe} & \textbf{Usage visé} & \textbf{Ressources minimales documentées} \\
\midrule
Installation manuelle & Paquets de la distribution, guide d'installation officiel, service par service & Apprentissage, environnement isolé & Contrôleur : 4~Go ; calcul : 2~Go (preuve de concept) \parencite{install-env} \\
\gls{packstack} & Installeur du projet RDO (CentOS Stream) ; haute disponibilité et mises à niveau hors périmètre & Preuve de concept & 16~Go de mémoire \parencite{packstack} \\
DevStack & Scripts shell, installation depuis les sources Git & Développement, tests, démonstration & Machine dédiée (non précisé) \parencite{devstack-index} \\
OpenStack-Ansible & Ansible ; sources dans des environnements virtuels Python, conteneurs LXC & Production & Test (AIO) : 8~vCPU, 8~Go, 50~Go \parencite{osa-aio} \\
Kolla-Ansible & Ansible ; images de conteneurs Docker ou Podman & Production, évaluation & 2~interfaces, 8~Go, 40~Go \parencite{kolla-ansible-quickstart} \\
Kayobe & Kolla-Ansible, plus configuration des hôtes et provisionnement du matériel (Bifrost) & Production, depuis le serveur nu & --- \parencite{kayobe-index} \\
OpenStack-Helm & Charts Helm sur Kubernetes & Cloud sur Kubernetes & Kubernetes 1.33 à 1.35 \parencite{helm-readme} \\
Sunbeam & Snap, Juju, Kubernetes (Canonical OpenStack) \parencite{sunbeam-multinode} & Petits clouds, de un à plusieurs nœuds & Un nœud : 4~cœurs, 16~Gio, 100~Gio \parencite{sunbeam-tuto} \\
RHOSO & Opérateurs sur OpenShift (Red Hat) & Clients de Red Hat & --- \parencite{rhoso-deploy} \\
TripleO & Retiré du projet (dépôt archivé en 2024) & --- & --- \parencite{tripleo-docs} \\
\bottomrule
\end{tabular}
\caption{Les principaux outils de déploiement d'OpenStack.}
\label{tab:dep-outils}
\end{table}

TripleO (« OpenStack sur OpenStack » : un cloud d'administration, l'\emph{undercloud}, déployait le cloud de production,
l'\emph{overcloud}) figure encore dans de nombreux supports ; il n'est plus maintenu. Red Hat, qui en tirait son outil
\emph{director}, déploie aujourd'hui son offre par des opérateurs sur OpenShift (RHOSO) \parencite{tripleo-docs,rhoso-deploy}. La suite
détaille DevStack (section~\ref{sec:dep-devstack}), Kolla-Ansible (section~\ref{sec:dep-kolla}) et OpenStack-Ansible
(section~\ref{sec:dep-osa}), revient sur l'installation manuelle (section~\ref{sec:dep-manuel}) pour comprendre ce que les outils
automatisent, situe les approches fondées sur Kubernetes (section~\ref{sec:dep-autres}), puis aide à choisir
(section~\ref{sec:dep-choix}) et à exploiter (section~\ref{sec:dep-exploitation}).

% ---------------------------------------------------------------------
\subsection{DevStack : l'environnement tout-en-un}\label{sec:dep-devstack}
% ---------------------------------------------------------------------

\begin{definition}{DevStack}
Série de scripts qui déploient rapidement un cloud OpenStack complet sur une seule machine, à partir des dépôts Git de chaque projet
\parencite{devstack-index}.
\end{definition}

\gls{devstack} sert à construire vite des environnements de développement, à décrire des configurations qui fonctionnent, à faciliter les
contributions et à fournir l'environnement des tests d'intégration de chaque modification proposée \parencite{devstack-source}. Il
installe depuis les sources, sur la branche \texttt{master} par défaut ou sur une branche \texttt{stable/<série>} pour une version plus
ancienne : c'est ce qui en fait l'outil idéal pour essayer une modification de code (dans Nova, par exemple) avant de la proposer. Il
demande une machine dédiée et jetable. Ubuntu 24.04 est la plus testée ; le projet suit aussi les deux dernières versions LTS
d'Ubuntu, Rocky Linux 9 et openEuler. L'installation dure de 15 à 30 minutes \parencite{devstack-index}.

\begin{figure}[!ht]
\centering
\adjustbox{max width=\linewidth}{%
\begin{tikzpicture}[>=Stealth, font=\footnotesize]
  \node[pcli, minimum width=3.6cm] (lc) at (1.9,2.4) {\texttt{local.conf}\\ {\scriptsize\mdseries mots de passe, services}};
  \node[pgris, minimum width=3.6cm] (git) at (1.9,0.3) {Dépôts Git des projets\\ {\scriptsize\mdseries opendev.org}};
  \node[pfoc, minimum width=2.6cm, minimum height=1.3cm] (st) at (7.4,1.35) {\texttt{stack.sh}};
  \node[pgris, minimum width=5.6cm] (pk) at (13.0,2.8) {Paquets de la distribution\\ {\scriptsize\mdseries MySQL, RabbitMQ, Apache}};
  \node[psvc, minimum width=5.6cm] (sv) at (13.0,1.35) {Services installés dans \texttt{/opt/stack}\\ {\scriptsize\mdseries fichiers de configuration, bases}};
  \node[psvc, minimum width=5.6cm] (sd) at (13.0,-0.1) {Unités systemd \texttt{devstack@service}};
  \node[pdat, minimum width=9.4cm] (cl) at (9.4,-1.9) {Cloud à un nœud : API, Horizon, fichier \texttt{openrc}};
  \draw[fluxr] (lc.east) -| (st.north) node[num, pos=0.25] {1};
  \draw[fluxr] (git.east) -| (st.south) node[num, pos=0.25] {2};
  \foreach \t in {pk,sv,sd} \draw[fluxo] (st.east) -- (\t.west);
  \node[num] at (9.7,1.35) {3};
  \draw[fluxo] (sd.south) -- (sd.south |- cl.north);
  \node[callout, minimum width=15.6cm] at (8.5,-3.3)
    {DevStack installe depuis les sources, sur une machine dédiée : pour apprendre et développer, pas pour héberger un service réel.};
\end{tikzpicture}}
\caption{Ce que fait \texttt{stack.sh}. Il lit \texttt{local.conf} (1) et les dépôts des projets (2), puis installe les paquets de la
distribution, installe et configure les services dans \texttt{/opt/stack} et les démarre comme unités systemd (3).}
\label{fig:dep-devstack}
\end{figure}

Le script \texttt{stack.sh} ne doit pas être lancé en tant que \texttt{root} : il s'exécute sous un compte dédié, \texttt{stack}, qui
dispose de \texttt{sudo} sans mot de passe \parencite{devstack-index}. Ses services de base (Keystone, Nova, Placement, Glance,
Cinder, Neutron avec OVN et Horizon) sont actifs par défaut ; les autres s'ajoutent avec \texttt{enable\_service} (c'est le cas de Swift)
ou, pour un projet qui n'est pas dans le dépôt de DevStack, avec \texttt{enable\_plugin} \parencite{devstack-source,devstack-config,devstack-plugins}.
Chaque service tourne en unité systemd, dont on lit le journal avec \texttt{journalctl} \parencite{devstack-systemd}.

\par\noindent\begin{minipage}{\linewidth}
\begin{lstlisting}[style=terminal]
$ sudo useradd -s /bin/bash -d /opt/stack -m stack
$ sudo chmod +x /opt/stack
$ echo "stack ALL=(ALL) NOPASSWD: ALL" | sudo tee /etc/sudoers.d/stack
$ sudo -u stack -i
$ git clone https://opendev.org/openstack/devstack
$ cd devstack && nano local.conf
$ ./stack.sh
$ sudo journalctl -f --unit devstack@n-cpu.service
\end{lstlisting}
\end{minipage}\par

\par\noindent\begin{minipage}{\linewidth}
\begin{lstlisting}
[[local|localrc]]
ADMIN_PASSWORD=secret
DATABASE_PASSWORD=$ADMIN_PASSWORD
RABBIT_PASSWORD=$ADMIN_PASSWORD
SERVICE_PASSWORD=$ADMIN_PASSWORD
HOST_IP=192.168.1.10
# ajouter un projet absent des services par défaut (ici Heat)
enable_plugin heat https://opendev.org/openstack/heat
# retirer un service dont on n'a pas besoin
disable_service horizon
\end{lstlisting}
\end{minipage}\par

Le fichier \texttt{local.conf} est le seul à personnaliser : il suit un format de type INI dont la section \texttt{[[local|localrc]]}
fixe les mots de passe (des caractères alphanumériques seulement), l'adresse de l'hôte, les services actifs, et même les dépôts et
les branches à cloner (\texttt{NOVA\_REPO}, \texttt{NOVA\_BRANCH}) \parencite{devstack-config}. Une fois le script terminé, on charge les
identifiants avec \texttt{. openrc} puis on utilise la ligne de commande \parencite{devstack-source}.

Pour plusieurs machines, le contrôleur reçoit un \texttt{local.conf} complet et chaque nœud de calcul un fichier réduit, qui
désigne le contrôleur (\texttt{SERVICE\_HOST}, \texttt{MYSQL\_HOST}, \texttt{RABBIT\_HOST}) et ne garde que des services
comme \texttt{n-cpu}. Il faut ensuite associer chaque nouvel hôte à la cellule, depuis le contrôleur, avec
\texttt{tools/discover\_hosts.sh} (section~\ref{sec:nova-archi}) \parencite{devstack-multinode}.

DevStack n'est pas un installateur généraliste : la plupart des combinaisons d'options sont rarement testées, et il modifie en
profondeur le système hôte \parencite{devstack-overview,devstack-source}. Conçu pour un nœud, il n'offre pas de haute disponibilité.
\texttt{./unstack.sh} arrête les services (en laissant MySQL et RabbitMQ), \texttt{./clean.sh} efface les données.

% ---------------------------------------------------------------------
\subsection{Kolla-Ansible : OpenStack en conteneurs}\label{sec:dep-kolla}
% ---------------------------------------------------------------------

Le projet \gls{kolla} a pour mission de fournir des outils qui créent des conteneurs prêts pour la production et des outils qui déploient
des clouds OpenStack \parencite{kolla-source}. Il faut distinguer ses livrables :

\begin{itemize}
  \item \textbf{Kolla} construit les images de conteneurs de chaque service (\texttt{kolla-build}) ;
  \item \textbf{Kolla-Ansible} fournit les playbooks Ansible qui déploient ces images sur les nœuds ;
  \item \textbf{\gls{kayobe}} ajoute la configuration des hôtes et le provisionnement du matériel (avec Bifrost, qui s'appuie sur
        Ironic) à Kolla-Ansible \parencite{kayobe-index}.
\end{itemize}

\gls{kollaansible} fonctionne avec Docker ou Podman. Il prend en charge comme hôtes Debian 13, Rocky Linux 10 et Ubuntu 24.04, et les
images existent dans les mêmes trois distributions (\texttt{kolla\_base\_distro}) \parencite{kolla-ansible-support}. Par défaut il
tire les images publiées sur Quay.io (espace \texttt{openstack.kolla}). Ses auteurs avertissent que ces images, destinées aux
essais et aux démonstrations, ne subissent aucun examen de sécurité : en production, on construit et on maintient son propre jeu
d'images avec \texttt{kolla-build}, que l'on place dans un \gls{registre} local (recommandé dès que le déploiement compte plusieurs nœuds)
\parencite{kolla-source,kolla-image-building,kolla-ansible-multinode}.

\begin{figure}[!ht]
\centering
\adjustbox{max width=\linewidth}{%
\begin{tikzpicture}[>=Stealth, font=\footnotesize,
    ct/.style={draw=insaorange, fill=insapeche, rounded corners=2pt, font=\scriptsize\ttfamily, minimum height=0.55cm, inner xsep=3pt}]
  \node[pcli, minimum width=7.2cm] (dh) at (3.9,7.1)
    {Hôte de déploiement\\ {\scriptsize\mdseries \texttt{kolla-ansible}, inventaire, \texttt{globals.yml}, \texttt{passwords.yml}}};
  \node[pgris, minimum width=5.6cm] (reg) at (12.2,7.1) {Registre d'images\\ {\scriptsize\mdseries Quay.io ou registre local}};
  \foreach \y/\nom in {5.0/{Contrôleurs $\times 3$}, 3.6/{Réseau}, 2.2/{Calcul}, 0.8/{Stockage}}
    { \draw[draw=insagris, dashed, rounded corners=3pt, fill=insagris!5] (0,{\y-0.58}) rectangle (15.7,{\y+0.58});
      \node[anchor=west, font=\footnotesize\bfseries] at (0.15,\y) {\nom}; }
  \foreach \x/\t in {3.9/keystone, 5.9/nova\_api, 8.5/neutron\_server, 10.9/horizon, 12.7/mariadb, 14.5/rabbitmq}
    \node[ct] at (\x,5.0) {\t};
  \foreach \x/\t in {3.9/haproxy, 6.0/keepalived, 8.9/neutron\_l3\_agent, 12.1/neutron\_dhcp\_agent}
    \node[ct] at (\x,3.6) {\t};
  \foreach \x/\t in {4.2/nova\_compute, 6.8/nova\_libvirt, 10.0/openvswitch\_vswitchd}
    \node[ct] at (\x,2.2) {\t};
  \node[ct] at (4.2,0.8) {cinder\_volume};
  \node[pcli, minimum width=1.0cm, inner xsep=4pt] (vip) at (14.8,3.6) {VIP};
  \node[pgris, dashed, minimum width=3.8cm] (ceph) at (11.8,0.8) {Ceph (externe)};
  \draw[fluxd] (6.0,0.8) -- (ceph.west);
  \draw[fluxd] (dh.south) -- (3.9,5.65) node[etiqf, midway, right=3pt] {Ansible, par SSH};
  \draw[fluxo] (reg.south) -- (12.2,5.65) node[etiqf, midway, right=3pt] {tirer les images};
  \node[callout, minimum width=15.4cm] at (7.85,-0.75)
    {Chaque service tourne dans son conteneur, avec ses dépendances figées dans l'image : pas de conflit de versions entre services.};
\end{tikzpicture}}
\caption{Kolla-Ansible sur plusieurs nœuds. L'hôte de déploiement pilote les nœuds par Ansible (pointillés) ; ceux-ci tirent leurs
images d'un registre. Les groupes de l'inventaire fixent les conteneurs de chaque nœud. HAProxy et keepalived, qui portent
l'adresse virtuelle (VIP), s'exécutent par défaut sur les nœuds réseau. Kolla-Ansible ne déploie pas Ceph : il s'y connecte.}
\label{fig:dep-kolla}
\end{figure}

La configuration tient en trois éléments, placés sur l'hôte de déploiement, qu'installe \texttt{pip} dans un environnement virtuel
\parencite{kolla-ansible-quickstart} :
\begin{itemize}
  \item l'\textbf{\gls{inventaireansible}} Ansible, qui affecte les machines aux groupes \texttt{control}, \texttt{network}, \texttt{compute},
        \texttt{monitoring} et \texttt{storage} (le groupe \texttt{control} compte un nombre impair de nœuds) ; chaque service a lui-même
        son groupe, mais en modifier la répartition peut casser le déploiement \parencite{kolla-ansible-multinode,kolla-ansible-production} ;
  \item \texttt{/etc/kolla/globals.yml}, le fichier principal ;
  \item \texttt{/etc/kolla/passwords.yml}, dont les mots de passe vides se remplissent avec \texttt{kolla-genpwd}.
\end{itemize}

\par\noindent\begin{minipage}{\linewidth}
\begin{lstlisting}[style=terminal]
$ python3 -m venv /opt/kolla-venv && source /opt/kolla-venv/bin/activate
$ pip install git+https://opendev.org/openstack/kolla-ansible@stable/2026.1
$ sudo mkdir -p /etc/kolla && sudo chown $USER:$USER /etc/kolla
$ cp -r /opt/kolla-venv/share/kolla-ansible/etc_examples/kolla/* /etc/kolla
$ cp /opt/kolla-venv/share/kolla-ansible/ansible/inventory/multinode .
$ kolla-ansible install-deps
$ kolla-genpwd
\end{lstlisting}
\end{minipage}\par

\par\noindent\begin{minipage}{\linewidth}
\begin{lstlisting}
# /etc/kolla/globals.yml
kolla_base_distro: "ubuntu"
network_interface: "eth0"
neutron_external_interface: "eth1"
kolla_internal_vip_address: "10.1.0.250"
enable_cinder: true
enable_octavia: true
\end{lstlisting}
\end{minipage}\par

La valeur de \texttt{kolla\_internal\_vip\_address} doit être une adresse inutilisée du réseau de gestion : keepalived la porte et la
déplace en cas de panne ; l'interface externe doit être active, mais sans adresse IP \parencite{kolla-ansible-quickstart}. Kolla-Ansible
est volontairement très directif : quelques variables suffisent à un déploiement fonctionnel, et il ne multiplie pas les options.
Pour tout le reste, les fichiers déposés dans \texttt{/etc/kolla/config/} sont fusionnés avec la configuration par défaut du service
(par exemple, \texttt{nova.conf} avec \texttt{virt\_type = qemu} pour faire tourner le cloud dans des machines virtuelles sans
virtualisation imbriquée) \parencite{kolla-ansible-philosophy}. Les services de base (Keystone, Glance, Nova avec Placement, Neutron,
Heat, Horizon) sont actifs par défaut ; on ajoute les autres avec les variables \texttt{enable\_*} (tableau~\ref{tab:dep-kolla-services})
\parencite{kolla-ansible-source}.

\begin{table}[!ht]
\centering
\renewcommand{\arraystretch}{1.3}
\begin{tabular}{@{}>{\raggedright\arraybackslash}p{2.6cm}>{\raggedright\arraybackslash}p{3.4cm}>{\raggedright\arraybackslash}p{9.5cm}@{}}
\toprule
\textbf{Service} & \textbf{Variable} & \textbf{À savoir} \\
\midrule
Cinder (chap.~\ref{chap:cinder}) & \texttt{enable\_cinder} & Inactif par défaut ; disques LVM (\texttt{enable\_cinder\_\allowbreak backend\_\allowbreak lvm}) ou cluster Ceph externe. \\
Octavia (chap.~\ref{chap:octavia}) & \texttt{enable\_octavia} & Demande des certificats propres (\texttt{kolla-ansible octavia-certificates}) \parencite{kolla-ansible-source}. \\
Ironic (chap.~\ref{chap:ironic}) & \texttt{enable\_ironic} & S'appuie sur la même base de déploiement que le reste. \\
Magnum (chap.~\ref{chap:magnum}) & \texttt{enable\_magnum} & Le pilote de cluster reste à choisir (section~\ref{sec:ma-pilotes}). \\
Trove (chap.~\ref{chap:trove}) & \texttt{enable\_trove} & Les images invitées restent à construire (section~\ref{sec:tr-datastores}). \\
Swift (chap.~\ref{chap:swift}) & --- & Plus déployé par Kolla-Ansible (retiré faute de mainteneurs) ; Glance peut utiliser un Swift externe \parencite{kolla-ansible-source}. \\
\bottomrule
\end{tabular}
\caption{Activer les autres services avec Kolla-Ansible.}
\label{tab:dep-kolla-services}
\end{table}

Le déploiement suit un enchaînement fixe (figure~\ref{fig:dep-kolla-flux} et tableau~\ref{tab:dep-kolla-cmd}). La commande
\texttt{deploy} peut être relancée autant de fois que nécessaire, mais si une tâche d'amorçage échoue, un nouvel appel ne la répare pas ;
le plus rapide est alors de détruire ce déploiement raté (\texttt{kolla-ansible destroy}) \parencite{kolla-ansible-troubleshooting}.

\begin{figure}[!ht]
\centering
\adjustbox{max width=\linewidth}{%
\begin{tikzpicture}[>=Stealth, font=\footnotesize]
  \foreach \x/\c/\d [count=\i] in {1.6/bootstrap-servers/{prépare les hôtes}, 4.8/prechecks/{vérifie les prérequis}, 8.0/pull/{tire les images},
                                   11.2/deploy/{démarre les conteneurs}, 14.4/post-deploy/{écrit \texttt{clouds.yaml}}}
    { \node[pfoc, minimum width=2.9cm, font=\scriptsize\bfseries\ttfamily] (s\i) at (\x,2.0) {\c};
      \node[num] at (\x-1.25,2.55) {\i};
      \node[etiqf] at (\x,1.3) {\d}; }
  \foreach \a/\b in {1/2,2/3,3/4,4/5} \draw[fluxr] (s\a.east) -- (s\b.west);
  \node[anchor=west, font=\footnotesize\bfseries, text=insagris] at (0,-0.1) {Ensuite :};
  \foreach \x/\c in {3.6/reconfigure, 7.0/upgrade, 10.4/mariadb-backup, 13.8/destroy}
    \node[psvc, minimum width=2.9cm, font=\scriptsize\bfseries\ttfamily] at (\x,-0.1) {\c};
  \node[callout, minimum width=15.6cm] at (8.0,-1.4)
    {Les mêmes commandes servent à installer, à reconfigurer et à mettre à niveau : l'inventaire et les variables restent la seule description du cloud.};
\end{tikzpicture}}
\caption{Les cinq étapes d'un déploiement avec Kolla-Ansible, puis les commandes d'exploitation.}
\label{fig:dep-kolla-flux}
\end{figure}

\begin{table}[!ht]
\centering
\renewcommand{\arraystretch}{1.25}
\begin{tabular}{@{}>{\raggedright\arraybackslash\ttfamily\small}p{3.3cm}>{\raggedright\arraybackslash}p{12.2cm}@{}}
\toprule
\textrm{\textbf{Commande}} & \textbf{Effet} \\
\midrule
bootstrap-servers & Prépare chaque hôte : paquets système, moteur de conteneurs, utilisateurs \parencite{kolla-ansible-quickstart}. \\
prechecks & Vérifie que les hôtes sont dans un état où l'on peut déployer \parencite{kolla-ansible-troubleshooting}. \\
pull & Tire les images sur les hôtes, sans toucher aux conteneurs. \\
deploy & Génère les configurations, amorce les services et démarre tous les conteneurs. \\
post-deploy & Écrit les accès administrateur dans \texttt{/etc/kolla/clouds.yaml} \parencite{kolla-ansible-quickstart}. \\
reconfigure & Applique une modification de \texttt{globals.yml} ou de \texttt{config/}, éventuellement pour un seul service (\texttt{-t mariadb}). \\
validate-config & Valide les fichiers de configuration générés (après un premier déploiement). \\
upgrade & Met à niveau d'une série à la suivante : nouvelles images, migrations de schéma \parencite{kolla-ansible-operating}. \\
mariadb-backup & Sauvegarde complète ou incrémentale (\texttt{--incremental}) de MariaDB \parencite{kolla-ansible-mariadb-backup}. \\
destroy & Détruit les conteneurs, les volumes et la configuration des hôtes. \\
\bottomrule
\end{tabular}
\caption{Les commandes de \texttt{kolla-ansible} \parencite{kolla-ansible-operating,kolla-ansible-source}.}
\label{tab:dep-kolla-cmd}
\end{table}

% ---------------------------------------------------------------------
\subsection{OpenStack-Ansible : sources et conteneurs LXC}\label{sec:dep-osa}
% ---------------------------------------------------------------------

\gls{openstackansible} (OSA) déploie OpenStack avec Ansible, lui aussi pour la production. Par défaut, tous les services sont installés
depuis les sources dans des \glspl{venv} Python, puis isolés dans des conteneurs \gls{lxc}. Ce sont des conteneurs « machine »,
proches de petites machines virtuelles, et non des conteneurs d'application comme ceux de Docker ; ils sont facultatifs, et un
déploiement directement sur les hôtes est possible \parencite{osa-about}. Les hôtes pris en charge sont Ubuntu 24.04 et 26.04, Debian 13,
CentOS Stream 10 et Rocky Linux 10 \parencite{osa-about}.

\begin{figure}[!ht]
\centering
\adjustbox{max width=\linewidth}{%
\begin{tikzpicture}[>=Stealth, font=\footnotesize,
    ct/.style={draw=insaorange, fill=insapeche, rounded corners=2pt, font=\scriptsize\ttfamily, minimum height=0.6cm, minimum width=2.1cm},
    lx/.style={draw=insarouge!80!black, fill=insarouge!12, rounded corners=3pt, text width=1.8cm, align=center, font=\scriptsize,
               minimum height=1.35cm, inner sep=2pt}]
  \node[pcli, minimum width=3.4cm] at (3.7,5.4) {Kolla-Ansible};
  \node[pcli, minimum width=3.4cm] at (12.0,5.4) {OpenStack-Ansible};
  \draw[draw=insagris, dashed, rounded corners=3pt, fill=insagris!5] (0,0.6) rectangle (7.4,4.7);
  \node[anchor=north west, font=\footnotesize\bfseries] at (0.15,4.6) {Hôte : Docker ou Podman};
  \node[ct] at (1.4,2.9) {keystone};
  \node[ct] at (3.7,2.9) {nova\_api};
  \node[ct] at (6.0,2.9) {mariadb};
  \node[font=\scriptsize, text width=6.6cm, align=center] at (3.7,1.4)
    {un conteneur d'application par service, créé à partir d'une image (Kolla)};
  \draw[draw=insagris, dashed, rounded corners=3pt, fill=insagris!5] (8.4,0.6) rectangle (15.8,4.7);
  \node[anchor=north west, font=\footnotesize\bfseries] at (8.55,4.6) {Hôte : LXC (ou directement sur l'hôte)};
  \node[lx] at (9.9,2.8) {\textbf{galera}\\ MariaDB};
  \node[lx] at (12.1,2.8) {\textbf{keystone}\\ environnement Python};
  \node[lx] at (14.3,2.8) {\textbf{nova\_api}\\ environnement Python};
  \node[font=\scriptsize, text width=6.8cm, align=center] at (12.1,1.2)
    {conteneurs « machine » LXC ; services installés depuis les sources};
  \node[callout, minimum width=15.6cm] at (7.9,-0.6)
    {Même outil (Ansible), deux façons d'isoler les services : des conteneurs d'application ou des conteneurs « machine ».};
\end{tikzpicture}}
\caption{Kolla-Ansible (à gauche) et OpenStack-Ansible (à droite). Le premier exécute chaque service dans un conteneur
d'application ; le second installe les services dans des environnements virtuels Python, au sein de conteneurs LXC.}
\label{fig:dep-osa}
\end{figure}

Le déploiement distingue l'hôte de déploiement, qui exécute Ansible, et les hôtes cibles. On y décrit le cloud dans
\texttt{/etc/openstack\_deploy/} : \texttt{openstack\_user\_config.yml} affecte les machines aux rôles (au moins trois hôtes
d'infrastructure dans \texttt{shared\_infra\_hosts}), \texttt{user\_variables.yml} règle les options globales et
\texttt{user\_secrets.yml} contient les mots de passe, que \texttt{pw-token-gen.py} génère \parencite{osa-config}. Le choix
\texttt{install\_method} (sources ou paquets) se fait au premier déploiement et ne peut plus changer ensuite \parencite{osa-config}.
Trois \glspl{playbook} font le travail, dans cet ordre \parencite{osa-run} : préparer les hôtes et construire les conteneurs, installer l'infrastructure (memcached, dépôt des sources, Galera, RabbitMQ), puis les services d'OpenStack.

\par\noindent\begin{minipage}{\linewidth}
\begin{lstlisting}[style=terminal]
# openstack-ansible openstack.osa.setup_hosts
# openstack-ansible openstack.osa.setup_infrastructure
# openstack-ansible openstack.osa.setup_openstack
\end{lstlisting}
\end{minipage}\par

Un déploiement « \gls{aio} » (\emph{all-in-one}, AIO), qui réunit tout sur une machine, sert à l'essai ou au laboratoire ; il demande
au minimum 8~vCPU, 8~Go de mémoire et 50~Go de disque (le projet en recommande 16~Go et 80~Go) \parencite{osa-aio}. OSA automatise
aussi les mises à niveau entre versions majeures. Le projet déconseille son choix si l'on veut des conteneurs d'application à 100~\%, si l'on
préfère Puppet, ou si l'on veut installer tous les services depuis les paquets de la distribution, dont la couverture reste partielle
\parencite{osa-about}.

% ---------------------------------------------------------------------
\subsection{Installation manuelle : ce que les outils automatisent}\label{sec:dep-manuel}
% ---------------------------------------------------------------------

Le guide d'installation officiel décrit, commande par commande, un cloud minimal installé avec les paquets de la distribution
\parencite{install-guide,install-overview}. Aucun outil ne s'interpose : on tape les commandes sur chaque nœud. On commence par
l'environnement commun (horloge synchronisée, dépôt des paquets OpenStack, MariaDB, RabbitMQ, memcached)
\parencite{install-env,install-packages-ubuntu,install-sql-ubuntu,install-messaging-ubuntu}, puis on installe les services dans
l'ordre que leurs dépendances imposent : Keystone, Glance, Placement, Nova, Neutron, puis Horizon et, en option, Cinder
\parencite{install-overview}. Chaque service suit le même scénario en six étapes (figure~\ref{fig:dep-manuel}).

\begin{figure}[!ht]
\centering
\adjustbox{max width=\linewidth}{%
\begin{tikzpicture}[>=Stealth, font=\footnotesize,
    cmd/.style={font=\scriptsize\ttfamily, text width=2.55cm, align=center}]
  \node[pgris, minimum width=15.6cm] (env) at (8.0,3.9)
    {Environnement commun, une fois sur le contrôleur : \ NTP \ $\cdot$ \ paquets OpenStack \ $\cdot$ \ MariaDB \ $\cdot$ \ RabbitMQ \ $\cdot$ \ memcached};
  \node[anchor=east, font=\footnotesize\bfseries, text=insarouge] at (15.9,2.85) {Pour chaque service (Glance en exemple) :};
  \foreach \x/\t/\c [count=\i] in {1.35/{Base de\\ données}/{CREATE DATABASE glance;},
        4.05/{Enregistrement\\ dans Keystone}/{openstack endpoint create},
        6.75/{Paquets}/{apt install glance},
        9.45/{Configuration}/{glance-api.conf},
        12.15/{Schéma de\\ la base}/{glance-manage db\_sync},
        14.85/{Redémarrage}/{systemctl restart glance-api}}
    { \node[pfoc, minimum width=2.45cm, minimum height=1.0cm] (s\i) at (\x,1.7) {\t};
      \node[num] at (\x-1.0,2.3) {\i};
      \node[cmd] at (\x,0.55) {\c}; }
  \foreach \a/\b in {1/2,2/3,3/4,4/5,5/6} \draw[fluxr] (s\a.east) -- (s\b.west);
  \draw[fluxd] (s1.north |- env.south) -- (s1.north);
  \node[callout, minimum width=15.6cm] at (8.0,-0.55)
    {Six étapes à répéter pour chaque service, sur chaque nœud : c'est précisément ce que les outils de déploiement automatisent.};
\end{tikzpicture}}
\caption{Le scénario d'installation manuelle d'un service, illustré par Glance \parencite{glance-install}. L'étape~2 inscrit le service
dans le catalogue de Keystone : un utilisateur de service, le rôle \texttt{admin}, une entrée de service et trois points d'accès.}
\label{fig:dep-manuel}
\end{figure}

\par\noindent\begin{minipage}{\linewidth}
\begin{lstlisting}[style=terminal]
$ sudo rabbitmqctl add_user openstack RABBIT_PASS
$ sudo rabbitmqctl set_permissions openstack ".*" ".*" ".*"
$ openstack user create --domain default --password-prompt glance
$ openstack role add --project service --user glance admin
$ openstack service create --name glance --description "OpenStack Image" image
$ openstack endpoint create --region RegionOne \
    image public http://controller:9292
\end{lstlisting}
\end{minipage}\par

Les deux premières commandes préparent RabbitMQ une fois pour toutes \parencite{install-messaging-ubuntu} ; les quatre suivantes
enregistrent Glance dans le catalogue (chapitre~\ref{chap:keystone}), puis on répète la dernière pour les points d'accès interne et
d'administration \parencite{glance-install}. Un nœud de calcul reçoit ensuite \texttt{nova-compute} et l'agent de Neutron, et doit être
associé à sa cellule, depuis le contrôleur, par \texttt{nova-manage cell\_v2 discover\_hosts} \parencite{nova-install-compute}.

Cette méthode a une vraie valeur pédagogique : on voit où chaque mot de passe, chaque base et chaque point d'accès apparaît. Elle ne
tient pas dans la durée : des dizaines de fichiers à garder cohérents, aucune idempotence (relancer une commande qui a déjà réussi
échoue ou duplique), aucune description unique du cloud, et des mises à niveau à la main. Au-delà d'un laboratoire, on délègue ce
travail à un outil.

\begin{remarque}[Deux sens de « bare metal »]
Le support parle d'installation « bare metal ». Le terme recouvre deux idées. Installer les services \emph{directement sur le
système des machines}, sans conteneur, c'est l'installation manuelle (ou OpenStack-Ansible sans LXC). Fournir \emph{des machines
physiques aux utilisateurs} est le rôle d'Ironic (chapitre~\ref{chap:ironic}) ; Kayobe s'en sert en outre pour installer le cloud lui-même
sur des serveurs nus (section~\ref{sec:dep-kolla}).
\end{remarque}

% ---------------------------------------------------------------------
\subsection{OpenStack sur Kubernetes}\label{sec:dep-autres}
% ---------------------------------------------------------------------

Deux approches installent OpenStack \emph{sur} Kubernetes, pour des publics différents. À ne pas confondre avec Magnum
(chapitre~\ref{chap:magnum}), qui fait l'inverse : il fournit des clusters Kubernetes \emph{sur} OpenStack.

\textbf{\gls{openstackhelm}} est une collection de charts Helm qui déploient OpenStack et les services associés sur Kubernetes
\parencite{helm-index,helm-readme}. Sa documentation vise des versions récentes de Kubernetes (1.33 à 1.35) et des hôtes Ubuntu
\parencite{helm-readme}. Elle s'adresse à une équipe qui exploite déjà Kubernetes et sait le faire vivre.

\textbf{\gls{sunbeam}} (Canonical OpenStack) s'installe avec un \emph{snap} et s'appuie sur \gls{juju} et Kubernetes. Le tutoriel officiel
déploie un cloud à un nœud (4~cœurs, 16~Gio de mémoire, 100~Gio de disque, Ubuntu 24.04) en cinq commandes \parencite{sunbeam-tuto} :

\par\noindent\begin{minipage}{\linewidth}
\begin{lstlisting}[style=terminal]
$ sudo snap install openstack
$ sunbeam prepare-node-script --bootstrap | bash -x && newgrp snap_daemon
$ sunbeam cluster bootstrap --accept-defaults --role control,compute,storage
$ sunbeam configure --accept-defaults --openrc demo-openrc
$ sunbeam launch ubuntu --name test
\end{lstlisting}
\end{minipage}\par

Pour agrandir le cloud, \texttt{sunbeam cluster add <FQDN>} produit un jeton que la nouvelle machine donne à
\texttt{sunbeam cluster join --role <rôles> -}. Les rôles sont \texttt{control}, \texttt{compute}, \texttt{network} et
\texttt{storage}, et \texttt{sunbeam cluster resize} répartit les fonctions de contrôle après l'ajout de nœuds de contrôle ;
les nœuds ne se désignent que par leur nom complet (FQDN) \parencite{sunbeam-multinode}. Red Hat suit une logique voisine avec
RHOSO, dont les opérateurs tournent sur OpenShift \parencite{rhoso-deploy}.

% ---------------------------------------------------------------------
\subsection{Choisir une méthode}\label{sec:dep-choix}
% ---------------------------------------------------------------------

Tous ces outils déploient le même OpenStack : les API, les services et les concepts des chapitres précédents ne changent pas. Seules
varient l'installation et l'exploitation, et c'est là que se joue le choix (figure~\ref{fig:dep-choix}).

\begin{figure}[!ht]
\centering
\adjustbox{max width=\linewidth}{%
\begin{tikzpicture}[>=Stealth, font=\footnotesize]
  \node[pcli, text width=6.0cm] (q1) at (3.4,5.3) {Apprendre ou développer, sur une machine jetable ?};
  \node[pcli, text width=6.0cm] (q2) at (3.4,3.1) {La plateforme existante repose-t-elle déjà sur Kubernetes ?};
  \node[pcli, text width=6.0cm] (q3) at (3.4,0.8) {Préférez-vous des conteneurs d'application, un par service ?};
  \draw[fluxd] (q1.south) -- (q2.north) node[etiqf, midway, right=2pt] {non};
  \draw[fluxd] (q2.south) -- (q3.north) node[etiqf, midway, right=2pt] {non};
  \node[pdat, text width=6.2cm] (r1a) at (12.2,5.9) {\textbf{DevStack} : développer, tester};
  \node[pdat, text width=6.2cm] (r1b) at (12.2,4.7) {\textbf{Installation manuelle}, \textbf{Sunbeam} à un nœud : apprendre};
  \node[psvc, text width=6.2cm] (r2)  at (12.2,3.1) {\textbf{OpenStack-Helm}, \textbf{Sunbeam}, \textbf{RHOSO}};
  \node[psvc, text width=6.2cm] (r3a) at (12.2,1.5) {\textbf{Kolla-Ansible} (Kayobe si le matériel est à provisionner)};
  \node[psvc, text width=6.2cm] (r3b) at (12.2,0.1) {\textbf{OpenStack-Ansible} (sources, conteneurs LXC)};
  \draw[fluxr] (q1.east) -- ++(0.75,0) |- (r1a.west);
  \draw[fluxr] (q1.east) -- ++(0.75,0) |- (r1b.west);
  \node[etiqf] at (7.0,5.55) {oui};
  \draw[fluxr] (q2.east) -- (r2.west) node[etiqf, midway, above] {oui};
  \draw[fluxr] (q3.east) -- ++(0.75,0) |- (r3a.west) node[etiqf, pos=0.8, above] {oui};
  \draw[fluxr] (q3.east) -- ++(0.75,0) |- (r3b.west) node[etiqf, pos=0.8, above] {non};
  \node[pdat, minimum width=0.5cm, inner xsep=4pt] at (2.6,-1.1) {};
  \node[anchor=west] at (2.95,-1.1) {essai, apprentissage};
  \node[psvc, minimum width=0.5cm, inner xsep=4pt] at (8.0,-1.1) {};
  \node[anchor=west] at (8.35,-1.1) {production};
  \node[callout, minimum width=15.6cm] at (7.9,-2.2)
    {Le bon choix dépend des compétences de l'équipe (Ansible, Kubernetes) et du support attendu, pas d'un classement général.};
\end{tikzpicture}}
\caption{Arbre de décision pour choisir une méthode de déploiement. Ce n'est pas la seule lecture possible : un fournisseur
(Canonical, Red Hat) peut imposer son offre, un client déjà équipé d'Ansible ira vers Kolla-Ansible ou OpenStack-Ansible.}
\label{fig:dep-choix}
\end{figure}

Quelques critères aident à départager Kolla-Ansible et OpenStack-Ansible, les deux outils de production les plus répandus :

\begin{itemize}
  \item \textbf{Isolation} : Kolla-Ansible isole chaque service dans une image figée ; OpenStack-Ansible installe depuis les sources
        dans des environnements virtuels, au sein de conteneurs LXC \parencite{osa-about}.
  \item \textbf{Configuration} : Kolla-Ansible est directif (quelques variables, puis des fichiers fusionnés) ; OpenStack-Ansible
        expose toute la configurabilité de ses rôles \parencite{kolla-ansible-philosophy,osa-about}.
  \item \textbf{Images} : avec Kolla-Ansible, il faut construire et héberger ses propres images pour la production ; avec
        OpenStack-Ansible, on gère un dépôt de sources et des environnements virtuels.
  \item \textbf{Mises à niveau} : toutes deux les automatisent (\texttt{kolla-ansible upgrade}, playbooks d'OpenStack-Ansible) ;
        aucune ne dispense de lire les notes de version.
\end{itemize}

% ---------------------------------------------------------------------
\subsection{Exploiter un cloud déployé}\label{sec:dep-exploitation}
% ---------------------------------------------------------------------

Déployer n'est que le début. Ce qui distingue un cloud fiable, c'est ce qu'on met en place autour : vérification, haute
disponibilité, sécurité, sauvegardes, supervision, mises à niveau.

\textbf{Vérifier.} Avant d'ouvrir le cloud aux utilisateurs, on contrôle que chaque service est enregistré et vivant. Avec
Kolla-Ansible, \texttt{post-deploy} écrit les accès administrateur dans \texttt{/etc/kolla/clouds.yaml}, et le script \texttt{init-runonce}
crée des réseaux, une image et d'autres ressources d'exemple (démonstration seulement) \parencite{kolla-ansible-quickstart} ; avec OpenStack-Ansible,
on utilise le conteneur utilitaire \parencite{osa-verify}.

\par\noindent\begin{minipage}{\linewidth}
\begin{lstlisting}[style=terminal]
$ openstack endpoint list                # chaque service a ses points d'accès
$ openstack compute service list         # conductor, scheduler, compute : up
$ openstack network agent list           # agents Neutron : Alive
$ openstack volume service list          # cinder-scheduler, cinder-volume : up
$ nova-status upgrade check              # prérequis d'une mise à niveau de Nova
\end{lstlisting}
\end{minipage}\par

\textbf{Haute disponibilité.} Elle repose sur deux principes \parencite{ha-concepts}. Les services sans état (API, ordonnanceurs)
tournent en actif/actif derrière HAProxy et une adresse virtuelle que \gls{keepalived} déplace en cas de panne. Les services à état
(MariaDB avec \gls{galera}, \gls{rabbitmq}) se répliquent en cluster d'un nombre impair de nœuds, dont le quorum est la majorité
(\(\lfloor n/2\rfloor+1\)) ; le \gls{plancontrole} compte donc au moins trois contrôleurs. La haute disponibilité ne protège pas
les machines des utilisateurs : si un hyperviseur tombe, Nova ne les relance pas seul (on évacue les instances avec
\texttt{openstack server evacuate}).

\begin{table}[!ht]
\centering
\renewcommand{\arraystretch}{1.25}
\small
\begin{tabular}{@{}>{\raggedright\arraybackslash}p{2.3cm}>{\raggedright\arraybackslash}p{13.1cm}@{}}
\toprule
\textbf{Sujet} & \textbf{Mise en œuvre avec Kolla-Ansible} \\
\midrule
Répartiteur & HAProxy et keepalived (groupe \texttt{loadbalancer}, qui reprend par défaut les nœuds \texttt{network}) ; variables \texttt{enable\_haproxy} et \texttt{enable\_keepalived} \parencite{kolla-ansible-haproxy}. \\
TLS & \texttt{kolla\_enable\_tls\_external} et \texttt{kolla\_enable\_tls\_internal} ; certificats dans \texttt{/etc/kolla/certificates} (\texttt{haproxy.pem}, \texttt{haproxy-internal.pem}) ; \texttt{kolla-ansible certificates} produit une autorité de test ; \texttt{enable\_letsencrypt} automatise le renouvellement \parencite{kolla-ansible-tls}. \\
Secrets & Mots de passe de \texttt{passwords.yml} (à sauvegarder avant tout changement). Rotation : vider les valeurs, \texttt{kolla-genpwd -p}, puis \texttt{kolla-ansible reconfigure} ; pour RabbitMQ, \texttt{stop} puis \texttt{deploy} \parencite{kolla-ansible-passwords}. \\
Sauvegarde & \texttt{kolla-ansible mariadb-backup} (complète ou incrémentale) \parencite{kolla-ansible-mariadb-backup} ; copie de \texttt{/etc/kolla} et de l'inventaire ; volumes et images selon les chapitres~\ref{chap:cinder} et~\ref{chap:glance}. \\
Journaux & Fichiers dans \texttt{/var/log/kolla/<service>} et \texttt{docker logs} ; \texttt{enable\_central\_logging} ajoute Fluentd et OpenSearch \parencite{kolla-ansible-troubleshooting,kolla-ansible-logging}. \\
Supervision & \texttt{enable\_prometheus} déploie Prometheus, que l'on peut relier à Grafana \parencite{kolla-ansible-prometheus}. \\
\bottomrule
\end{tabular}
\caption{L'exploitation d'un cloud Kolla-Ansible : ce que l'outil apporte.}
\label{tab:dep-exploit}
\end{table}

\textbf{Sécurité et sauvegardes.} Trois règles reviennent partout : garder le réseau de gestion interne (il porte les mots de passe
et les messages), chiffrer les API exposées (TLS) et protéger les secrets (tableau~\ref{tab:dep-exploit}). Une sauvegarde ne vaut que
si on sait restaurer : la base de données de chaque service, la configuration de l'outil et les mots de passe sont à sauvegarder
régulièrement et à restaurer lors d'un exercice.

\textbf{Mises à niveau.} Une version sort tous les six mois, mais on n'est pas obligé de les suivre toutes. Une version sur deux,
la première de l'année (2025.1, 2026.1\dots), est une version \gls{slurp} : on peut passer directement de l'une à la suivante, alors
qu'une version intermédiaire ne se met à niveau que vers la version qui la suit \parencite{openstack-slurp} (figure~\ref{fig:dep-slurp}).

\begin{figure}[!ht]
\centering
\adjustbox{max width=\linewidth}{%
\begin{tikzpicture}[>=Stealth, font=\footnotesize, adj/.style={-{Stealth[length=2mm]}, thick, insagris}]
  \foreach \x/\v/\n/\k [count=\i] in {1.5/2025.1/{}/psvc, 4.5/2025.2/{}/pgris, 7.5/2026.1/{Gazpacho}/psvc, 10.5/2026.2/{Hibiscus}/pgris, 13.5/2027.1/{Indri}/psvc}
    { \node[\k, minimum width=2.5cm, minimum height=1.1cm] (v\i) at (\x,1.0) {\textbf{\v}\\ {\scriptsize\mdseries \n}}; }
  \node[font=\scriptsize] at (1.5,0.2) {SLURP};
  \node[font=\scriptsize] at (7.5,0.2) {SLURP};
  \node[font=\scriptsize] at (13.5,0.2) {SLURP, en développement};
  \foreach \a/\b in {1/2,2/3,3/4,4/5} \draw[adj] (v\a.east) -- (v\b.west);
  \draw[fluxo] (v1.north) to[out=65,in=115] node[etiqf, midway, above] {saut SLURP $\to$ SLURP} (v3.north);
  \draw[fluxo] (v3.north) to[out=65,in=115] node[etiqf, midway, above] {saut SLURP $\to$ SLURP} (v5.north);
  \node[callout, minimum width=15.6cm] at (7.5,-0.8)
    {Depuis 2025.2, il faut passer par 2026.1 : une version non SLURP ne se met à niveau que vers la suivante.};
\end{tikzpicture}}
\caption{Chemins de mise à niveau (flèches grises : version suivante, toujours possible). Les noms des versions
2026.1, 2026.2 et 2027.1 sont ceux de la documentation \parencite{openstack-releases}.}
\label{fig:dep-slurp}
\end{figure}

On met toujours le plan de contrôle à niveau avant les nœuds de calcul, que l'on traite par lots : services anciens et récents
doivent coexister un temps. Nova impose un ordre (Placement d'abord, puis \texttt{nova-conductor} en premier et \texttt{nova-api} en
dernier), puis des migrations de données « en ligne » (\texttt{nova-manage db online\_data\_migrations}) \parencite{nova-upgrades}.
Kolla-Ansible enchaîne tout cela dans une seule commande \parencite{kolla-ansible-operating} :

\par\noindent\begin{minipage}{\linewidth}
\begin{lstlisting}[style=terminal]
$ cp -r /etc/kolla /etc/kolla.sauvegarde        # configuration d'abord
$ pip install --upgrade \
    git+https://opendev.org/openstack/kolla-ansible@stable/2026.2
$ kolla-ansible install-deps
$ kolla-mergepwd --old passwords.yml.old --new passwords.yml.new \
    --final /etc/kolla/passwords.yml
$ kolla-ansible pull && kolla-ansible prechecks
$ kolla-ansible upgrade
\end{lstlisting}
\end{minipage}\par

Avant cela, on lit les notes de version de la série visée, on met à jour l'inventaire et \texttt{globals.yml} d'après les exemples
fournis, et l'on fusionne les mots de passe avec \texttt{kolla-mergepwd}. Pour un saut SLURP, RabbitMQ ne supporte pas un saut de deux
versions : il faut d'abord le mettre à niveau vers une version intermédiaire. La documentation ne décrit pas de retour arrière :
on se protège par une sauvegarde de la base et de la configuration, et par un essai préalable sur un environnement de test
\parencite{kolla-ansible-operating}.

% ---------------------------------------------------------------------
\subsection{Exercices corrigés}\label{sec:dep-exercices}
% ---------------------------------------------------------------------

\begin{exercice}{Choisir un outil et une topologie}\label{ex:dep-choix}
Une PME de vingt utilisateurs veut une plateforme interne de développement sur quatre serveurs physiques. Une personne de l'équipe
maîtrise Ansible, personne ne connaît Kubernetes. Proposer un outil, une répartition des rôles et les points à prévoir avant
l'ouverture.
\end{exercice}
\begin{correction}
Pas de réponse unique ; l'argumentation compte. Production interne, compétence Ansible, aucun Kubernetes : on écarte DevStack,
Packstack, l'installation manuelle et les approches fondées sur Kubernetes, et l'on retient \textbf{Kolla-Ansible} ou
\textbf{OpenStack-Ansible}. Trois serveurs forment les contrôleurs (quorum de MariaDB et de RabbitMQ) et portent aussi le réseau et
le calcul ; le quatrième est un nœud de calcul. Les machines des utilisateurs partagent alors les serveurs avec le plan de
contrôle : acceptable à cette échelle, à surveiller. À prévoir : un hôte de déploiement distinct, un registre d'images local, une
adresse virtuelle inutilisée sur le réseau de gestion, TLS, la sauvegarde de MariaDB et de la configuration, une mise à niveau
annuelle sur les versions SLURP.
\end{correction}

\begin{exercice}{Relire une configuration}\label{ex:dep-relire}
Cet extrait d'un déploiement Kolla-Ansible contient trois erreurs de conception ; lesquelles ? On sait que \texttt{node1} a l'adresse
\texttt{10.1.0.11} sur \texttt{eth0}.

\begin{lstlisting}
[control]
node1
node2
[network]
node1
# globals.yml
kolla_internal_vip_address: "10.1.0.11"
network_interface: "eth0"
neutron_external_interface: "eth0"
\end{lstlisting}
\end{exercice}
\begin{correction}
(1) \textbf{Deux contrôleurs} n'apportent aucune haute disponibilité : la panne de l'un fait perdre le quorum ; il en faut trois.
(2) L'\textbf{adresse virtuelle est celle de \texttt{node1}} : elle doit être inutilisée sur le réseau de gestion, car keepalived la
porte et la déplace. (3) L'\textbf{interface externe est celle de gestion} : \texttt{neutron\_external\_interface} doit être une
interface dédiée, active et sans adresse IP, qui porte \texttt{br-ex}. Au passage, le groupe \texttt{network} se réduit à
\texttt{node1} : HAProxy et keepalived, qui suivent ce groupe, n'ont alors aucun nœud de secours.
\end{correction}

\begin{exercice}{Ajouter un service à un cloud en production}\label{ex:dep-service}
Le cloud, déployé avec Kolla-Ansible, doit fournir des répartiteurs de charge (chapitre~\ref{chap:octavia}). Quelles modifications et
quelles commandes faut-il ? Que reste-t-il à faire hors de l'outil ?
\end{exercice}
\begin{correction}
\begin{lstlisting}[style=terminal]
$ vi /etc/kolla/globals.yml                # enable_octavia: true
$ kolla-ansible octavia-certificates
$ kolla-ansible prechecks
$ kolla-ansible deploy
\end{lstlisting}
\texttt{deploy} est rejouable : il applique l'état décrit par les variables, sans repartir de zéro. Les certificats propres à
Octavia sont indispensables. Hors de l'outil, il reste à fournir ce que le chapitre~\ref{chap:octavia} décrit : l'image
d'amphore, le réseau et le groupe de sécurité de gestion, la \emph{flavor} des amphores.
\end{correction}

\begin{exercice}{Planifier une mise à niveau}\label{ex:dep-niveau}
Un cloud tourne en 2025.2 et l'équipe vise 2026.2. Peut-elle y passer directement ? Décrire la première mise à niveau avec
Kolla-Ansible.
\end{exercice}
\begin{correction}
Non : 2025.2 n'est pas une version SLURP, et ne se met à niveau que vers la suivante, 2026.1 ; de là, 2026.2 est la version
suivante. Il faut donc \textbf{deux mises à niveau successives}, chacune suivie d'une validation. Pour la première : lire les notes
de version de 2026.1 ; sauvegarder \texttt{/etc/kolla} et la base ; installer la branche \texttt{stable/2026.1} de
\texttt{kolla-ansible} ; \texttt{kolla-ansible install-deps} ; mettre à jour l'inventaire et \texttt{globals.yml} d'après les
exemples et fusionner les mots de passe avec \texttt{kolla-mergepwd} ; \texttt{kolla-ansible pull} puis \texttt{prechecks} ; enfin
\texttt{kolla-ansible upgrade}, qui recrée les conteneurs et migre les schémas. Aucun retour arrière n'est documenté : la
sauvegarde et l'essai sur un environnement de test sont la seule protection.
\end{correction}

% ---------------------------------------------------------------------
\subsection{Récapitulatif du chapitre}\label{sec:dep-recap}
% ---------------------------------------------------------------------

\begin{aretenir}
\begin{itemize}
  \item Un déploiement assemble des services, une infrastructure partagée (base, file de messages, cache, répartiteur) et des
        machines en rôles ; l'outil de déploiement automatise toute la pile.
  \item Un nombre impair de contrôleurs (trois au moins) pour le quorum ; séparer au moins les réseaux de gestion, de tunnel et
        externe, et garder le réseau de gestion interne.
  \item DevStack (sources, machine jetable) : développement et apprentissage. Kolla-Ansible (images de conteneurs, à construire
        soi-même en production) et OpenStack-Ansible (sources, conteneurs LXC) : production.
  \item TripleO est retiré ; Kolla-Ansible ne déploie plus Swift et ne déploie pas Ceph.
  \item L'exploitation compte autant que l'installation : haute disponibilité, TLS, secrets, sauvegardes, supervision, mises à
        niveau (SLURP : une version sur deux, avec saut direct).
\end{itemize}
\end{aretenir}

\begin{center}
\renewcommand{\arraystretch}{1.3}
\begin{tabular}{@{}>{\raggedright\arraybackslash}p{4.0cm}>{\raggedright\arraybackslash}p{9.8cm}r@{}}
\toprule
\textbf{Tâche} & \textbf{Commande} & \textbf{§} \\
\midrule
Préparer hôtes et secrets & \texttt{kolla-ansible bootstrap-servers}, \texttt{kolla-genpwd} & \ref{sec:dep-kolla} \\
Vérifier puis déployer & \texttt{kolla-ansible prechecks}, \texttt{pull}, \texttt{deploy} & \ref{sec:dep-kolla} \\
Appliquer une modification & \texttt{kolla-ansible reconfigure [-t SERVICE]} & \ref{sec:dep-kolla} \\
Sauvegarder MariaDB & \texttt{kolla-ansible mariadb-backup [--incremental]} & \ref{sec:dep-exploitation} \\
Mettre à niveau & \texttt{kolla-ansible upgrade} & \ref{sec:dep-exploitation} \\
Contrôler les services & \texttt{openstack compute service list}, \texttt{network agent list} & \ref{sec:dep-exploitation} \\
\bottomrule
\end{tabular}
\end{center}

\begin{savoirfaire}
\begin{itemize}
  \item Choisir un outil et une topologie à partir d'un contexte, en justifiant (exercice~\ref{ex:dep-choix}).
  \item Relire un inventaire et un \texttt{globals.yml} pour y repérer les erreurs de quorum, d'adresse virtuelle et d'interfaces
        (exercice~\ref{ex:dep-relire}).
  \item Ajouter un service à un cloud Kolla-Ansible et nommer ce qui reste à faire hors de l'outil (exercice~\ref{ex:dep-service}).
  \item Planifier un chemin de mise à niveau avec les versions SLURP (exercice~\ref{ex:dep-niveau}).
\end{itemize}
\end{savoirfaire}
